\documentclass[11pt]{article}
\usepackage[margin=1.08in]{geometry}
\usepackage{amsmath,amssymb,amsthm,mathtools,mathrsfs}
\usepackage{booktabs,array,longtable}
\usepackage{xcolor}
\usepackage[colorlinks=true,linkcolor=blue!60!black,citecolor=blue!60!black,urlcolor=blue!60!black]{hyperref}

\newtheorem{theorem}{Theorem}[section]
\newtheorem{proposition}[theorem]{Proposition}
\newtheorem{lemma}[theorem]{Lemma}
\newtheorem{corollary}[theorem]{Corollary}
\newtheorem{cproposition}[theorem]{Computational Proposition}
\theoremstyle{definition}
\newtheorem{definition}[theorem]{Definition}
\theoremstyle{remark}
\newtheorem{remark}[theorem]{Remark}

\newcommand{\C}{\mathbb C}
\newcommand{\R}{\mathbb R}
\newcommand{\g}{\mathfrak g}
\newcommand{\W}{\mathcal W}
\newcommand{\F}{\mathcal F}
\newcommand{\cH}{\mathcal H}
\newcommand{\cL}{\mathcal L}
\newcommand{\cE}{\mathcal E}
\newcommand{\id}{\mathrm{id}}
\newcommand{\im}{\operatorname{im}}
\newcommand{\rank}{\operatorname{rank}}
\newcommand{\Tr}{\operatorname{Tr}}
\newcommand{\vol}{\operatorname{vol}}
\newcommand{\CS}{\operatorname{CS}}

\title{Residual $SO(4,2)$ Cohomology of Free Weyl Gravity\\
on the Conformal Cylinder:\\
Cartan Homotopy, Weyl-Square Classes, and Their Residual Pairing}
\author{GPT-5.6.sol\thanks{OpenAI model and principal programme author.}
\and Asger Alstrup Palm\thanks{Programme orchestrator and corresponding human contact; Honorary Professor, DTU Compute, Technical University of Denmark. \texttt{asger@area9.dk}.}}
\date{Draft for external review, July 2026}

\begin{document}
\maketitle

\begin{abstract}
We compute the residual cohomology of free four-dimensional Weyl gravity on
the conformal cylinder $\R\times S^3$ in a selected closed-universe problem
that treats all fifteen conformal transformations---including compact time
translation $D$---as residual gauge transformations.  This is not the
fixed-cylinder quantization in which $D$ is the Hamiltonian, and it is not a
Hadamard or interacting quantum construction.  The coefficient space is the
parity-complete on-shell
Weyl-curvature Fock module, whose one-particle invariant form has signs
$+I_E\oplus(-I_A)\oplus(-I_L)$.  Cartan's identity
$d\iota_D+\iota_Dd=\mathcal L_D$ contracts every nonzero compact-degree
subcomplex.  In the centered four-ghost polarization, the vacuum and
one-particle sectors are acyclic and the complete two-particle calculation
gives
\[
H^4_{\rm res}=\operatorname{span}\{[W_+^2],[W_-^2]\},
\qquad \operatorname{sig}G_{\rm res}=(2,0),
\qquad G_{\rm res}=I_2.
\]
These are ghost-dressed deformation classes, not one-particle gravitons.
Their residual Hermitian form is positive-definite; the displayed $I_2$
uses the selected basis, orientation, and unit-overlap normalization.
Their parity combinations descend to $C^2$ and $C\widetilde C$; after
quotienting the Pontryagin boundary direction, the dynamical deformation
space is one-dimensional with Gram matrix $I_1$.

We connect this calculation to the covariant free BV complex.  BGG and
cylinder intertwiners identify the Bach quotient with the $E/A/L$ curvature
module, while support-local retracts preserve the BV current.  Null symbol
ranks $11$ and $9$ rule out a scalar-wave witness on the original field
bundle and isolate helicity $\pm2$.  A locally constraint-adjusted $26$-state
Weyl--Cotton system propagates all fourteen constraints.  Adjoint-tractor and
hybrid contractions give all-row causal homotopies whose causal difference is
a compact-to-spacelike-compact quasi-isomorphism.  It recovers the residual
endpoints, transports $SO(4,2)$ up to chain homotopy, and identifies Green and
current pairings.  Hence
$H^4_{\rm cov}=\operatorname{span}\{[W_+^2],[W_-^2]\}$ and
$\operatorname{sig}G_{\rm cov}=(2,0)$, with $G_{\rm cov}=I_2$ after the
same normalization choices.

We do not claim a positive graviton Hilbert space, an integrated $SO(4,2)$
representation, quantum anomaly cancellation, or interacting unitarity.
\end{abstract}

\section*{Authorship and generated inputs}
The named model author produced the mathematical programme, derivations,
verification code, referee responses, and manuscript.  Asger Alstrup Palm
commissioned and directed the work, initiated the questions, and served as
non-technical orchestrator and corresponding human contact; he claims no
technical contribution.  Repository scripts generate the identified tables,
certificates, and selected displayed summaries; those artifacts are
computational inputs rather than additional authors.

\tableofcontents

\section{Introduction}\label{sec:introduction}

Pure conformal gravity in four dimensions is defined, up to an overall
coupling and the Euler density, by the Weyl-square action
\begin{equation}\label{eq:weyl-action}
 S_W[g]=-\alpha_g\int_M d^4x\,\sqrt{-g}\,
 C_{\mu\nu\rho\sigma}C^{\mu\nu\rho\sigma}.
\end{equation}
Its fourth-order linearized equation has more solutions than the Einstein
equation, and a conventional local reduction produces an indefinite
oscillator form.  This familiar observation does not by itself determine the
state space after all gauge constraints have been imposed.  On a compact
spatial slice, background conformal symmetries are reducibilities of the
local Diff$\times$Weyl gauge system.  Treating those residual transformations
as gauge introduces a second, finite-dimensional reduction which is absent
from a naive oscillator count.  It is also a substantive physical choice.
In a Dirac or background-independence interpretation of a closed universe,
all residual conformal transformations can be quotiented.  In a fixed
cylinder conformal field theory, by contrast, $D$ is ordinarily the
Hamiltonian and the conformal group acts globally.  This paper computes the
first choice; it does not assert that strict pure-Weyl gravity uniquely
selects it.

The organizing question is therefore more basic than a sign assignment:
before asking whether a gravitational mode has positive or negative norm,
one must first determine whether it survives the full gauge reduction being
imposed.  The answer below concerns the residual reduction.  The
smooth metric-to-curvature bridge is proved globally, and a symbolic
all-energy cylinder construction identifies its algebraic positive-energy
part with the $E/A/L$ module used below.  The passage to the complete
selected algebraic local-plus-residual BV--BFV state complex is also proved
below.  Its natural energy-mode Krein--Fock completion is proved as well.
The reduced Lorentzian metric fields are then identified with that
completion by a field-induced branch Cauchy--Sobolev theorem.  The auxiliary
route now supplies the curved four-row operator identities, a
support-preserving all-row deformation retract, and the covariant current
comparison.  It also supplies a basis-independent scalar-wave no-go theorem:
at a null covector the action Hessian has rank $11$, whereas the gauge symbol
has rank $9$.  The uncancelled rank-two quotient is precisely helicity
$\pm2$, and the reduced Weyl symbol acts on it by $\frac14 I_2$.  Thus a
causal contraction must carry the physical sector through curvature or
tractor variables rather than hide it in a $24$-component scalar-wave block.
The
  Weyl--Cotton equations, their constraint-adjusted symmetric-hyperbolic
  closure, sourced constraint propagation, and the all-level $E/A/L$ audit
  are proved.  The support-local all-row prolongation retract and prolonged
  BV differential are coefficientwise complete as well.  The curved adjoint-tractor BGG
  retract transfers the parent causal homotopy to the trace-free endpoint;
  a local trace/Weyl shear and the hybrid algebraic projector extend it to
  all $386$ prolonged rows.  The induced causal map is now certified as a
  quasi-isomorphism, all fifteen residual endpoints and their duals are
  recovered without duplication, the $SO(4,2)$ action transfers up to an
  explicit compactly supported chain homotopy, and the direct cyclic
  Green/current comparison fixes the covariant Gram matrix.  Thus the
  covariant $H^4$ theorem is transport of the existing residual calculation,
  not a second cohomology computation.  Distributional/Hadamard completion
  and alternative boundary polarizations remain open.  A canonical endpoint
  Green inverse is not part of the transferred construction, while a direct
  invariant two-wave metric factorization is excluded in the certified
  family.

This paper starts from the classical linearized BV--BFV complex and then,
after the stated positive-frequency polarization, studies its completed free
BRST state representation on
\begin{equation}
 M=\R\times S^3,
 \qquad
 \bar g=-dt^2+d\Omega_3^2,
\end{equation}
with unit cylinder radius.  Three spaces must be distinguished:
\begin{equation}\label{eq:three-levels}
 \begin{aligned}
 \text{local oscillator cohomology}
 &\supseteq \text{classically integrable perturbations},\\
 \text{classically integrable perturbations}
 &\longrightarrow \text{residual-BRST classes}.
 \end{aligned}
\end{equation}
The first contains the familiar Einstein, vector-descendant, and logarithmic
or higher-derivative towers.  The last is obtained only after the residual
$SO(4,2)$ constraints and their ghosts are included.  A linearized
oscillator need not survive this final reduction.

Our result is most naturally stated in the language of residual vertex
cohomology.  The Hamada ghost polarization supplies a top-degree four-ghost
factor of compact degree $-4$.  It pairs with scalar matter primaries of
weight four.  The complete centered residual cohomology consists of two
chiral Weyl-square classes, and their residual CE Hermitian form is
positive-definite and normalized to the identity matrix $I_2$.  Their parity
combinations are the dynamical
Weyl-square density and the topological Pontryagin density.  Thus the
positive matrix is attached to classical vertex or theory-space directions.
It is not a positive norm on a propagating graviton Fock space; indeed, the
one-particle residual cohomology vanishes.

\paragraph{Relation to previous work.}
The conformal deformation detour complex on obstruction-flat backgrounds
was constructed by Branson and Gover \cite{BransonGover}; its formal
self-adjointness and related Yang--Mills detour constructions were developed
in \cite{GoverSombergSoucek}.  Gover and Peterson give the oriented
four-dimensional deformation sequence, while \v{C}ap proves that the flat
BGG sequence is a fine resolution and computes the same cohomology as the
flat twisted de Rham complex \cite{GoverPeterson,CapBGG}.  Cylinder
harmonics, oscillator charges, and
primary building blocks were developed by Hamada and Horata and by Hamada
\cite{HamadaHorata,HamadaBuildingBlocks}.  In the broader
conformal-mode/Riegert programme, Hamada identified the residual
diffeomorphisms with cylinder conformal transformations
\cite{HamadaCFT}, constructed the fifteen-generator residual BRST charge,
used the product of four proper-conformal ghosts as the centered ghost
vacuum, derived the scalar weight-four condition, and displayed a
Weyl-square BRST state \cite{HamadaBRST}.  None of those ingredients is
claimed as new here.  We use the harmonic and residual-algebra conventions,
but we do not import the Riegert field or the quantum claims of that model.
The Weyl-curvature module and its character are consistent with the general
conformal higher-spin description and partition-function literature
\cite{KuzenkoPonds,BeccariaBekaertTseytlin}.  Finally, the identification of
$C^2$ as the unique local four-derivative self-interaction of one linear
conformal graviton, under the assumptions stated there, is due to Boulanger
and Henneaux \cite{BoulangerHenneaux}.

The defensible new result is narrower and more concrete.  Starting with the
parity-complete pure-Weyl coefficient module, we compute the complete
minimal \emph{absolute} residual complex in the centered degree.  We prove
that the vacuum and one-particle sectors vanish, that precisely two
two-particle classes survive after the incoming-image test, and that their
normalized Hermitian residual CE Gram matrix is $I_2$.  The standard
Cartan homotopy makes the particle-number inventory exhaustive, and the
Euler--Lagrange map splits the two positive directions into one dynamical
and one topological class.

\paragraph{Claim boundary.}
Eight statements must remain separate.
\begin{enumerate}
\item The \emph{algebraic residual theorem} is exact for the minimal
absolute $SO(4,2)$ Chevalley--Eilenberg complex built on the on-shell Weyl
module, the stated centered ghost polarization, and full residual conformal
gauging.
\item The smooth complexified metric-to-geometric-curvature bridge is an
exact global theorem.  Its restriction to the $D$-finite cylinder category
is also exact: the explicit $E/A/L$ harmonics give same-energy metric
preimages, nonzero curvature intertwiners, and character exhaustion.
These statements concern the geometric and algebraic realizations; the
energy-mode completion is a separate theorem below.
\item In the finite polynomial conformal category, an exact metric-BV
calculation now includes the ghost and antifield detour rows, constructs a
homotopy-equivariant retract, transfers the strict residual action, and
reproduces the centered kernel and image.  Cross-energy recursion constructs
its invariant physical form in the raw basis; the raw cyclic retract and
dressed isometry are exact through the centered window.  The minimal rows
are derived from the field master action, and a specified gauge-fixed
nonminimal extension contracts back to them exactly.  For the selected
algebraic BV--BFV suspension and positive-frequency polarization, the endpoint
map has scalar $+1$, the polarized row ledger is exhaustive, and the
matter-times-ghost pairing transfers exactly.  This is not a statement about
a unique boundary polarization.
\item The normalized $E/A/L$ module has a canonical infinite-index
energy-mode Krein completion, its bosonic Fock completion is Krein, and the
residual BRST operator has a closed maximal block realization.  The bounded
Cartan contraction reduces its cohomology to the unchanged finite centered
block.  This mode-space theorem alone does not construct the reduced field
Cauchy realization below or integrate the unbounded Lie-algebra action to
$SO(4,2)$.
\item On the Lorentzian cylinder the tensor- and vector-curl identities give
local normally hyperbolic factors for the reduced TT and transverse-vector
operators.  The lower TT, vector, and upper TT branches are exactly the
$E/A/L$ towers.  Their field-action residues determine a branch
Cauchy--Sobolev space Krein-unitarily equivalent to the completed
one-particle module.  This reduced theorem is not yet a Green's witness for
the complete BV complex.
\item The action-derived curved auxiliary Hessian and four-row operator
identities, the local BV-canonical deformation retract, and the covariant
current comparison are exact.  They are recorded respectively by
\texttt{curved\_operator\_identity},
\texttt{curved\_deformation\_retract}, and
\texttt{curved\_current\_comparison}, all of which are true.  The same
calculation proves that the former scalar-wave witness ansatz is impossible;
it does not supply causal Green homotopies.
\item The null-symbol quotient is the two-dimensional helicity-$\pm2$
module and the reduced Weyl symbol is $\frac14 I_2$.  The $26$-state
Weyl--Cotton equations, constraint-adjusted symmetric hyperbolicity, sourced
constraint propagation and all-level $E/A/L$ spectrum are certified.  The
support-local all-row prolongation and BV operator identity are also certified.
The retained $30$-row metric endpoint and its canonical backward witness are
coefficient-complete.  A normalization audit gives scalar trace-free biwave
symbol and reduces the open primal channel to a rank-nine operator.  Its
complete invariant two-wave-factor family is ruled out by an
order-two polynomial certificate.  The canonical upper curvature lift is
coefficientwise complete, the canonical middle graph lift is obstructed, and the minimal cyclic
relative saddle has zero Schur correction.  A flat adjoint-tractor HPL
compression extends to a curved cyclic BGG retract: ten invariant
curvature channels cancel the former $48$ entries, and the compressed parent
middle operator equals the complete Bach block with normalization $-2$ at
every derivative order.  Parent causal transfer, an explicit trace/Weyl
shear, and the $356/30$ hybrid retract give a causal homotopy on all $386$
rows.  Thirteen curvature/current/causal/transport flags are true.  The two
remaining false flags concern only the superseded canonical endpoint Green
inverse and prolonged-witness implementation; the direct tractor route does
not require them.  The exact dependency DAG therefore promotes
\texttt{final\_covariant\_H4}.
\item Calling the resulting classes physical states additionally requires a
boundary and gauge principle which makes all fifteen conformal generators,
including $D$, gauge rather than global charges.  We call them
\emph{candidate physical classes under full residual conformal gauging}.
\end{enumerate}
No statement in this paper establishes quantum BRST nilpotency, anomaly
cancellation, interacting pseudo-unitarity, or a positive asymptotic
particle space.

\begin{figure}[htbp]
\centering
\small
\fbox{\begin{minipage}{0.94\textwidth}
\centering
\textbf{Selected boundary and gauging input}\par
$\Sigma=S^3$ closed; zero boundary/asymptotic charge; no external clock or
deparametrization; all fifteen residual reducibilities gauged; centered
positive-frequency polarization.

\medskip
\begin{minipage}[t]{0.47\linewidth}
\centering
\textbf{Paper A spine: residual calculation}\par
\medskip
$E/A/L$ Weyl module with its invariant Krein form

$\Downarrow$

polarized free symmetric algebra and absolute residual CE complex

$\Downarrow$

Cartan contraction and exhaustive centered $N=0,1,2$ calculation

$\Downarrow$

$H^4_{\rm res}=\langle[W_+^2],[W_-^2]\rangle$,\par
$\operatorname{sig}G_{\rm res}=(2,0)$,
$G_{\rm res}=I_2$ after normalization

$\Downarrow$

descent to $C^2,C\widetilde C$ and the one-dimensional dynamical quotient
\end{minipage}
\hfill
\begin{minipage}[t]{0.47\linewidth}
\centering
\textbf{Paper B spine: covariant transport}\par
\medskip
metric BV detour complex and global BGG bridge

$\Downarrow$

support-local auxiliary/curvature retracts and current comparison

$\Downarrow$

Weyl--Cotton evolution, sourced constraints, and all-row causal homotopy

$\Downarrow$

causal--Cauchy quasi-isomorphism, endpoint recovery, and
$SO(4,2)$-equivariant pairing transport

$\Downarrow$

$H^4_{\rm cov}\cong H^4_{\rm res}$ with the same signature and normalized
Gram form
\end{minipage}

\medskip
\hrule
\medskip
The right spine imports the already computed residual module and pairing
from the left spine; it does not repeat the Chevalley--Eilenberg calculation.
Both conclusions remain classical, free, and dependent on the selected
closed-universe boundary/gauging input above.
\end{minipage}}
\caption{Dependency diagram separating the residual computation from its
covariant causal transport.}
\label{fig:two-spine-dependency}
\end{figure}

The local-to-residual bridge belongs to the general subject of local BRST
cohomology with antifields, gauge fixing, and descent
\cite{BarnichBrandtHenneaux,BarnichHenneauxHurthSkenderis}.  Metsaev's
ordinary-derivative BRST--BV construction gives a concrete realization of
conformal fields and their conformal algebra on fields and antifields
\cite{MetsaevBRSTBV}.  Dneprov and Grigoriev construct a minimal nonlinear
Weyl-jet model and compatible presymplectic structure
\cite{DneprovGrigoriev}.  Their curvature-dependent ghost transformations
make that model a valuable comparison, but not an immediate proof that the
\emph{free} transferred residual differential is strict.  Neither work
supplies the global cylinder zero-mode normalization used here.

\section{The free local pure-Weyl complex}\label{sec:local-detour}

\subsection{Gauge, curvature, and equation maps}

Let $h_{\mu\nu}$ be a metric perturbation and let $(\xi^\mu,\sigma)$ be an
infinitesimal Diff$\times$Weyl parameter.  The full linear gauge map is
\begin{equation}\label{eq:full-K}
 K(\xi,\sigma)_{\mu\nu}
 =\cL_\xi\bar g_{\mu\nu}+2\sigma\bar g_{\mu\nu}
 =2\bar\nabla_{(\mu}\xi_{\nu)}+2\sigma\bar g_{\mu\nu}.
\end{equation}
After quotienting the trace, it becomes the conformal Killing operator
\begin{equation}\label{eq:K0}
 (K_0\xi)_{\mu\nu}
 =2\bar\nabla_{(\mu}\xi_{\nu)}
 -\frac12\bar g_{\mu\nu}\bar\nabla\!\cdot\!\xi.
\end{equation}
Let
\begin{equation}
 C_1h=\left.\frac{d}{d\varepsilon}
 C[\bar g+\varepsilon h]\right|_{\varepsilon=0}
\end{equation}
be the linearized Weyl-curvature map, and let $B_{\rm lin}$ denote the
action-normalized linearized Bach operator.  Formal adjoints use the
spacetime action pairing on a domain for which temporal and spatial boundary
terms vanish.

The conformal cylinder is conformally flat.  Naturality and conformal
covariance therefore give the exact local identity
\begin{equation}\label{eq:C1K}
 \boxed{C_1K=0}
\end{equation}
for arbitrary local gauge parameters, not only for the fifteen residual
conformal Killing parameters.

\begin{proposition}[Action-normalized factorization]\label{prop:factorization}
Use the repository normalization
\begin{equation}\label{eq:Sred}
 S_{\rm red}[g]=\int\sqrt{-g}
 \left(R_{\mu\nu}R^{\mu\nu}-\frac13R^2\right)
 =\frac12\int\sqrt{-g}\,(C^2-E_4).
\end{equation}
On a conformally flat background, modulo the Euler boundary term,
\begin{equation}\label{eq:Bfactor}
 \boxed{B_{\rm lin}=C_1^\sharp C_1.}
\end{equation}
Consequently,
\begin{equation}\label{eq:detour-identities}
 B_{\rm lin}K=0,
 \qquad
 K^\sharp B_{\rm lin}=0,
 \qquad
 B_{\rm lin}^\sharp=B_{\rm lin}.
\end{equation}
\end{proposition}

\begin{proof}
Since $C[\bar g]=0$, the mixed second variation of \eqref{eq:Sred} is
\begin{equation}
 \delta_h\delta_k S_{\rm red}
 =\langle C_1h,C_1k\rangle_{\W}.
\end{equation}
The definition
$\delta S_{\rm red}=\langle h,\cE(g)\rangle$ and
$B_{\rm lin}=D\cE|_{\bar g}$ gives
$\langle h,B_{\rm lin}k\rangle
=\langle C_1h,C_1k\rangle_{\W}$, proving
\eqref{eq:Bfactor}.  Equation \eqref{eq:C1K} then gives the first detour
identity, formal self-adjointness gives the second, and Hessian symmetry gives
the third.
\end{proof}

The Weyl-fibre pairing in Lorentz signature is indefinite.  Thus
$C_1^\sharp C_1$ is a formal-adjoint factorization, not a positivity
factorization.  In trace-free conformal geometry the corresponding local
sequence is
\begin{equation}\label{eq:detour-sequence}
 \Gamma(T)
 \xrightarrow{K_0}
 \Gamma(S^2_0T^*)
 \xrightarrow{B_{\rm lin}}
 \Gamma(S^2_0T^*)
 \xrightarrow{K_0^\sharp}
 \Gamma(T^*).
\end{equation}
On Bach-flat backgrounds this is the conformal deformation detour complex
of \cite{BransonGover}.  Formal self-adjointness holds in all signatures;
ellipticity is a Riemannian statement and is not assumed on the Lorentzian
cylinder.

\subsection{The global four-dimensional deformation complex}

The detour sequence \eqref{eq:detour-sequence} is not the only complex
available on the conformally flat cylinder.  On an oriented
four-dimensional conformally flat pseudo-Riemannian manifold, the adjoint
BGG deformation complex is \cite{GoverPeterson,CapBGG}
\begin{equation}\label{eq:BGG-sequence}
 \Gamma(T)
 \xrightarrow{K_0}
 \Gamma(S^2_0T^*[2])
 \xrightarrow{C_1}
 \Gamma(\mathcal C[2])
 \xrightarrow{C_1^\sharp\star}
 \Gamma(S^2_0T^*[-2])
 \xrightarrow{K_0^\sharp}
 \Gamma(T^*[-2]).
\end{equation}
Here $\mathcal C[2]$ is the algebraic Weyl-curvature bundle and
$C_1^\sharp\star$ means $C_1^\sharp\circ\star$, with the Hodge star acting
on one antisymmetric index pair.  In particular,
\begin{equation}\label{eq:star-complex-identity}
 \boxed{C_1^\sharp\star C_1=0.}
\end{equation}
The repository certificate verifies this identity directly as a
constant-coefficient differential-operator identity in both Euclidean and
Lorentzian conventions; Gover and Peterson supply the invariant global
sequence.

On a locally conformally flat manifold, the flat BGG complex is a fine
resolution of the sheaf of parallel sections of the flat adjoint-tractor
bundle \cite{CapDeformations,CapBGG}.  The cylinder
\begin{equation}
 M=\R\times S^3\simeq S^3
\end{equation}
is simply connected, so that flat local system is the constant system with
fibre $\mathfrak{so}(4,2)_\C$.  Fine-resolution cohomology therefore gives
an all-level global statement without an elliptic argument.

\begin{theorem}[Global smooth deformation cohomology]
\label{thm:global-BGG}
For smooth complexified global sections on the conformal cylinder,
\begin{equation}\label{eq:global-BGG-cohomology}
 H^q_{\rm def}(M)
 \cong H^q(S^3;\C)\otimes\mathfrak{so}(4,2)_\C
 \cong
 \begin{cases}
  \mathfrak{so}(4,2)_\C,&q=0,3,\\
  0,&q=1,2,4.
 \end{cases}
\end{equation}
In particular, the metric and curvature slots are globally exact:
\begin{equation}\label{eq:global-BGG-exactness}
 \boxed{\ker C_1=\im K_0,}
 \qquad
 \boxed{\ker(C_1^\sharp\star)=\im C_1.}
\end{equation}
\end{theorem}

\begin{proof}
The fine resolution computes sheaf cohomology of the flat adjoint local
system.  Simple connectivity trivializes its monodromy, while
$\R\times S^3$ deformation retracts onto $S^3$.  The only nonzero ordinary
cohomology groups are in degrees zero and three.  Vanishing in degrees one
and two is exactly the pair of statements
\eqref{eq:global-BGG-exactness}.
\end{proof}

The theorem is signature-independent and concerns unrestricted smooth
global sections.  It does not by itself establish continuity, closed range,
or bounded inverses on a Hilbert or Krein completion.

\subsection{The Bach quotient as an on-shell curvature space}

The passage from deformation-complex exactness to a detour equation is a
general exact-sequence argument.

\begin{proposition}[Detour--curvature lemma]\label{prop:detour-curvature}
Let
\begin{equation}
 E_0\xrightarrow{K}E_1\xrightarrow{C}E_2
 \xrightarrow{D_2}E_3
\end{equation}
be a complex which is exact at $E_1$ and $E_2$, and let $G:E_2\to E_1'$ be
an operator such that the detour equation is $B=GC$.  Then $[h]\mapsto Ch$
induces an isomorphism
\begin{equation}\label{eq:abstract-detour-curvature}
 \boxed{
 \frac{\ker B}{\im K}
 \cong \ker D_2\cap\ker G.}
\end{equation}
\end{proposition}

\begin{proof}
If $Bh=0$, then $GCh=0$, while $D_2Ch=0$ because the displayed sequence is
a complex.  If $Ch=0$, exactness at $E_1$ gives $h=K\xi$, so the induced map
is injective.  Conversely, if $U\in\ker D_2\cap\ker G$, exactness at $E_2$
gives $U=Ch$; then $Bh=GU=0$, proving surjectivity.
\end{proof}

The middle local cohomology relevant to the free field is
\begin{equation}\label{eq:local-middle}
 \mathcal H_B:=\frac{\ker B_{\rm lin}}{\im K_0}.
\end{equation}
In the four-dimensional deformation complex,
\begin{equation}
 D_2=C_1^\sharp\star,
 \qquad G=C_1^\sharp,
 \qquad B_{\rm lin}=GC_1.
\end{equation}
Define the geometrically on-shell curvature space
\begin{equation}\label{eq:Wgeom}
 \mathcal W_{\rm geom}
 :=\ker(C_1^\sharp\star)\cap\ker C_1^\sharp.
\end{equation}
The smooth global BGG theorem and
Proposition~\ref{prop:detour-curvature} identify the metric quotient with
this space exactly.

\begin{theorem}[Cylinder detour--curvature isomorphism]
\label{thm:curvature-isomorphism}
For smooth complexified fields on the conformal cylinder, the curvature map
$[h]\mapsto C_1h$ induces an $SO(4,2)$-equivariant isomorphism
\begin{equation}\label{eq:curvature-isomorphism}
 \boxed{
 \frac{\ker B_{\rm lin}}{\im K_0}
 \cong\mathcal W_{\rm geom}
 =\ker C_1^\sharp\cap\ker(C_1^\sharp\star).}
\end{equation}
\end{theorem}

\begin{proof}
If $B_{\rm lin}h=0$ and $U=C_1h$, then the factorization
\eqref{eq:Bfactor} gives $C_1^\sharp U=0$, while
\eqref{eq:star-complex-identity} gives $C_1^\sharp\star U=0$.  Thus the
curvature map lands in the displayed intersection.

If $C_1h=0$, exactness at the metric slot in
\eqref{eq:global-BGG-exactness} gives $h=K_0\xi$, proving injectivity.
Conversely, suppose $U$ lies in both kernels.  Exactness at the curvature
slot gives $U=C_1h$ for a smooth global $h$, and then
$B_{\rm lin}h=C_1^\sharp U=0$.  Hence every element of the intersection has
a Bach-flat metric potential, proving surjectivity.
\end{proof}

In Lorentz signature, complexified Weyl tensors split into Hodge
eigenspaces $U=U_++U_-$ with $\star U_\pm=\pm iU_\pm$.  If
$a=C_1^\sharp U_+$ and $b=C_1^\sharp U_-$, the two target equations are
\begin{equation}
 a+b=0,
 \qquad ia-ib=0,
\end{equation}
so $a=b=0$.  Define the smooth on-shell chiral geometric spaces by
\begin{equation}\label{eq:smooth-Weyl-spaces}
 \mathcal W_{{\rm geom},\pm}^{\rm sm}
 :=\{U_\pm:\star U_\pm=\pm iU_\pm,
                \ C_1^\sharp U_\pm=0\}.
\end{equation}
Then Theorem~\ref{thm:curvature-isomorphism} becomes
\begin{equation}\label{eq:smooth-local-target}
 \boxed{
 \frac{\ker B_{\rm lin}}{\im K_0}
 \cong
 \mathcal W_{{\rm geom},+}^{\rm sm}
 \oplus\mathcal W_{{\rm geom},-}^{\rm sm}.}
\end{equation}

All maps commute with $SO(4,2)$.  Hence the quotient-to-curvature map
preserves compact energy and both rotation spins on quotient classes.  It
remains to match this geometric system to the abstract positive-energy
module constructed in Section~\ref{sec:weyl-module}.

\begin{theorem}[Algebraic cylinder realization]\label{thm:algebraic-cylinder}
Let $\mathcal W_{{\rm geom},\pm}^{\rm alg,+}$ denote the $D$-finite,
$SO(4)$-finite positive-energy solutions of
\eqref{eq:smooth-Weyl-spaces}.  The curvature construction restricts to an
isomorphism
\begin{equation}\label{eq:local-target}
 \boxed{
 \mathcal W_{{\rm geom},+}^{\rm alg,+}
 \oplus\mathcal W_{{\rm geom},-}^{\rm alg,+}
 \cong\W_+\oplus\W_-.}
\end{equation}
Every curvature basis vector has an explicit Bach-flat metric representative
in the same $D$-energy and $SO(4)$ type.
\end{theorem}

\begin{proof}
Put $r=\tan(\beta/2)$, $z=1+r^2$, and
$q=d\beta+i\sin\beta\,d\gamma$.  For positive chirality the normalized
highest-weight metric representatives are the radiation-gauge cylinder
modes
\begin{equation}\label{eq:metric-preimage-families}
 \begin{array}{c|c|c}
  \text{tower}&J&\text{nonzero metric component}\\
  \hline
  E_n&n/2&h_{ij}=N_E(J)T[J]_{ij}e^{-int}\\
  A_n&(n-1)/2&h_{0i}=N_A(J)V[J]_ie^{-int}\\
  L_n&(n-2)/2&h_{ij}=N_L(J)T[J]_{ij}e^{-int}.
 \end{array}
\end{equation}
Here
\begin{align}
 V[J]&=-\frac{\sqrt{2J}}{4\pi}
 \cos^{2J-1}\!\frac\beta2\,
 e^{-i[(J+\frac12)\alpha+(J-\frac12)\gamma]}q,\nonumber\\
 T[J]&=\frac{\sqrt{2(2J-1)}}{16\pi}
 \cos^{2J-2}\!\frac\beta2\,
 e^{-i[(J+1)\alpha+(J-1)\gamma]}q\otimes q,
 \label{eq:metric-preimage-harmonics}
\end{align}
and
\begin{equation}
 N_E=\frac1{4\sqrt{J(2J+1)}},\qquad
 N_A=\frac1{2\sqrt{(2J-1)(2J+1)(2J+3)}},\qquad
 N_L=\frac1{4\sqrt{(J+1)(2J+1)}}.
\end{equation}

An exact coordinate calculation retains symbolic integer $n$, constructs
all components of $C_1h$, and applies $C_1^\sharp$ independently.  After
suppressing the common factor
$e^{-i(n t+m_L\alpha+m_R\gamma)}z^{-a}$, nonzero pivots are
\begin{align}
 (C_1h_E)_{0202}
 &=\frac{\sqrt{n(n-1)(n+1)}}{16\pi z^2},\nonumber\\
 (C_1h_A)_{0102}
 &=\frac{\sqrt{n(n-2)(n-1)}}{32\pi\sqrt{n+2}\,z},
 \label{eq:curvature-pivots}\\
 (C_1h_L)_{0202}
 &=\frac{\sqrt{n(n-3)(n-1)}}{16\pi z^2}.\nonumber
\end{align}
They are nonzero at every allowed tower energy.  The full tensors are
algebraically trace-free, have one common Hodge chirality, and obey
$C_1^\sharp C_1h=0$ identically in $n$.  Naturality of $C_1$ extends each
nonzero highest-weight intertwiner to its complete $SO(4)$ irrep.  Defining
$U_{F,n,M}=C_1h_{F,n,M}$ therefore gives the explicit same-block inverse
\begin{equation}\label{eq:cylinder-right-inverse}
 R_nU_{F,n,M}=h_{F,n,M},\qquad C_1R_n=\id,
 \qquad F\in\{E,A,L\}.
\end{equation}
The orientation-reversing cylinder isometry $\alpha\leftrightarrow\gamma$
exchanges the two $SU(2)$ factors and supplies the other chirality.

For exhaustion, the flat positive-chirality BGG data consist of the
lowest-weight curvature primary $(2;2,0)$, the equation module $(4;1,1)$,
and its identity $(5;\tfrac12,\tfrac12)$.  Their alternating character is
\begin{equation}
 \chi_{(2;2,0)}-\chi_{(4;1,1)}
 +\chi_{(5;\frac12,\frac12)}
 =\frac{5q^2-9q^4+4q^5}{(1-q)^4}.
\end{equation}
The cylinder module $\W_+$ is cyclic from the same primary and has this
character by \eqref{eq:weyl-character}; parity supplies the other chirality.
The nonzero graded intertwiner and equality of every graded $SO(4)$
multiplicity prove exhaustion.  Equations
\eqref{eq:metric-preimage-families}--\eqref{eq:cylinder-right-inverse}
supply finite harmonic metric representatives directly, excluding secular
or infinite-harmonic preimages in this algebraic category.
The exact coordinate tensor calculation, a machine-readable certificate,
and the generated pivot table are produced by
\path{symbolic/verify_conformal_cylinder_preimages.py}.
\end{proof}

\subsection{The finite-jet calculation as an independent convention audit}

The earlier finite-jet calculation remains useful, but it is no longer a
substitute for an unknown all-level theorem.  On homogeneous Euclidean
polynomial jets of degree $n=2,\ldots,6$, the matrices for $K$, $C_1$, and
$B_1$ obey
\begin{equation}
 C_1K=0,
 \qquad B_1K=0,
 \qquad \ker C_1=\im K.
\end{equation}
The exact ranks are shown in Table~\ref{tab:detour-ranks}.

\begin{table}[ht]
\centering
\small
\begin{tabular}{c@{\quad}rrrrrrrr}
\toprule
$n$ & $\dim G_n$ & $\dim h_n$ & $\dim W_{n-2}$ & $\dim B_{n-4}$
& $\rank K$ & $\rank C_1$ & $\rank B_1$
& $\dim(\ker B_1/\im K)$\\
\midrule
2 & 90  & 100 & 10  & 0   & 90  & 10  & 0  & 10\\
3 & 160 & 200 & 40  & 0   & 160 & 40  & 0  & 40\\
4 & 259 & 350 & 100 & 10  & 259 & 91  & 9  & 82\\
5 & 392 & 560 & 200 & 40  & 392 & 168 & 32 & 136\\
6 & 564 & 840 & 350 & 100 & 564 & 276 & 74 & 202\\
\bottomrule
\end{tabular}
\caption{Exact homogeneous-polynomial detour ranks.  The finite Euclidean
jet calculation independently audits the conventions and the first five
levels of the global BGG theorem.}
\label{tab:detour-ranks}
\end{table}

The quotient dimensions $(10,40,82,136,202)$ agree with the first five
levels of the parity-complete $E/A/L$ module introduced below.  The same
calculation separates a $15$-dimensional kernel of the low-degree gauge map,
generated by four translations, six rotations, one dilation, and four
special conformal transformations.  These finite results are independent
regression rails for \eqref{eq:global-BGG-exactness} and
\eqref{eq:local-target} in the repository normalization.

\subsection{The finite global zero-mode pair}

The same topology that kills the metric and curvature cohomology leaves two
finite-dimensional groups:
\begin{equation}\label{eq:BGG-zero-mode-pair}
 H^0_{\rm def}(M)\cong\mathfrak{so}(4,2)_\C,
 \qquad
 H^3_{\rm def}(M)\cong\mathfrak{so}(4,2)_\C.
\end{equation}
The degree-zero copy is the familiar space of fifteen conformal Killing
fields.  The degree-three copy lies on the adjoint equation side of the
self-dual deformation complex.  Refining both groups by compact weight gives
\begin{equation}
 4_{-1}\oplus7_0\oplus4_{+1}.
\end{equation}
Poincar\'e duality on $S^3$ and the Killing form pair the two copies, with
\begin{equation}
 \kappa(\g_r,\g_s)=0\qquad\text{unless}\qquad r+s=0.
\end{equation}

Its dimension, representation, and position are exactly those expected of
the dual residual constraint sector.  It is therefore a strong candidate
for a geometric realization of the fifteen moment-map or Taub components
constructed in Section~\ref{sec:moment-map}.  The full five-term BGG sequence
must nevertheless be kept distinct from the shorter four-term Bach/BV
detour complex.  Proposition~\ref{prop:dual-endpoint} below derives the
endpoint duality of the latter directly from cyclic finite-dimensional
linear algebra, without identifying it by fiat with BGG $H^3$.  Locating the
degree-three BGG bundle representative remains a separate geometric
comparison, not a premise of the BV endpoint theorem.

Thus the smooth deformation complex contains a controlled finite pair,
not an unspecified collection of extra global classes.  The shorter BV
detour chain has its own exact endpoint quotient, while comparison of the
two realizations remains useful but is not needed for the residual rank
calculation.

\subsection{Energy-mode Krein completion and its boundary}
\label{sec:energy-mode-completion}

The smooth and algebraic theorems above now admit a natural completion in
the positive-frequency cylinder-energy variables used by the residual
calculation.  This is deliberately narrower than a covariant completion of
the metric Bach/BV complex.

Let
\begin{equation}
 \W_{\rm alg}=\bigoplus_{n\ge2}^{\rm alg}\cH_n
\end{equation}
be the parity-complete $E/A/L$ module.  Declare its normalized cylinder
basis orthonormal for a positive Hilbert majorant and complete the direct
sum to $\mathscr H_1$.  The block involution
\begin{equation}
 \mathfrak J_1=+1_E\oplus(-1_A)\oplus(-1_L)
\end{equation}
is a bounded self-adjoint involution of norm one, and
$[u,v]_1=(u,\mathfrak J_1v)_0$ extends the algebraic invariant form.  Both
the positive and negative eigenspaces are infinite-dimensional.  Thus this
is an infinite-index Krein space, not a Pontryagin space.

The reduced proper-conformal coefficients obey an all-level linear bound
$|c_{\alpha\beta}(n)|\le 2(1+n)$.  With the unit-normalized Clebsch--Gordan
maps, $K^\pm$ are finite-band matrix-weighted shifts of energy order one.
Their finite-support realizations are closable, energy truncations are graph
cores, and their closures equal the maximal square-summability domains.
The equality
\begin{equation}
 (\overline{K^-})^\sharp=\overline{K^+}
\end{equation}
includes equality of those maximal domains.  The conformal commutators are
claimed on the common finite-energy invariant core only; no exponentiation
to a global Krein-unitary $SO(4,2)$ representation is inferred
\cite{JorgensenTian}.

Second quantization gives the Hilbert Fock space
\begin{equation}
 \mathscr F=\Gamma_s(\mathscr H_1),
 \qquad
 \mathfrak J_{\mathscr F}=\Gamma_s(\mathfrak J_1),
\end{equation}
and hence a bosonic Krein--Fock completion.  Give the finite residual ghost
exterior algebra its positive monomial Hilbert norm and set
$\mathscr K=\mathscr F\widehat\otimes\Lambda^\bullet\g^*$.  The centered
ghost insertion remains a bounded complementary-degree cohomological
pairing; it is not promoted here to a nondegenerate Krein metric on every
ghost-number block.

\begin{theorem}[Energy-mode Krein completion]
\label{thm:energy-mode-completion}
For the selected closed-cylinder problem in which all fifteen residual
generators, including $D$, are gauged, the algebraic residual differential
has a closed densely defined maximal extension $\overline Q$ on
$\mathscr K$, and the algebraic state space is a graph core.  Every total
compact-degree eigenspace $\mathscr K_\delta$ is finite-dimensional.  On
the complement of $\mathscr K_0$, the bounded operator
\begin{equation}
 h=\bigoplus_{\delta\ne0}\frac{\iota_D}{\delta}
\end{equation}
preserves $\operatorname{Dom}\overline Q$ and satisfies
\begin{equation}
 \overline Qh+h\overline Q=1-P_0.
\end{equation}
The range of $\overline Q$ is closed, ordinary and reduced cohomology
coincide, and
\begin{equation}
 \boxed{
 H^4(\mathscr K,\overline Q)
 =\operatorname{span}\{[W_+^2],[W_-^2]\},
 \qquad G_{\rm completed}=I_2.}
\end{equation}
\end{theorem}

\begin{proof}
The ghost compact spectrum is contained in $[-4,4]$.  At fixed total degree
$\delta$, matter energy therefore belongs to a finite interval.  Fixed
matter energy bounds particle number by $N\le E/2$, uses only the finite
one-particle blocks with $n\le E$, and has finitely many bosonic partitions.
Thus every $\mathscr K_\delta$ is finite-dimensional.  In particular,
$\mathscr K_0$ has matter energy at most four and is literally the algebraic
centered block.

Write $Q_\delta$ for the finite nilpotent matrix on $\mathscr K_\delta$.
The maximal Hilbert direct sum, with domain
$\sum_\delta\|Q_\delta\Psi_\delta\|^2<\infty$, is closed; total-degree
truncations converge in norm and graph norm.  Blockwise nilpotency also
places $\overline Q\Psi$ back in the domain and gives
$\overline Q^2\Psi=0$.

Exterior contraction $\iota_D$ is bounded of norm one, and nonzero
$\delta$ is integral, so $\|h\|\le1$.  The estimate
\begin{equation}
 \|Q_\delta h_\delta\Psi_\delta\|
 \le \|\Psi_\delta\|+\|h_\delta\|\,
       \|Q_\delta\Psi_\delta\|
\end{equation}
proves domain invariance and extends the algebraic Cartan identity to the
maximal domain.  Hence every off-center closed vector is exact, and the
off-center range equals the closed kernel.  The centered range is
finite-dimensional, so the total range is closed.  The already-certified
finite centered calculation then applies without a limiting argument and
gives the displayed cohomology and Gram matrix.
\end{proof}

This theorem completes modes of the already-selected polarized Weyl state
complex.  It does not by itself identify a metric-field Cauchy topology.
That reduced one-particle comparison is supplied next.  The theorem still
does not make the choice to gauge $D$ universal: if $D$ is retained as a
Hamiltonian or boundary charge, the Cartan contraction is not part of that
physical complex.

\subsection{Lorentzian Cauchy--Sobolev realization}
\label{sec:cauchy-sobolev}

The energy-mode completion can now be reached directly from reduced metric
fields on the Lorentzian cylinder.  Let $D_i$ denote the Levi--Civita
connection of the unit $S^3$.  On transverse-traceless spatial tensors and
transverse one-forms define
\begin{equation}
 (\mathcal C_2h)_{ij}=\epsilon_i{}^{kl}D_kh_{lj},
 \qquad
 (\mathcal C_1v)_i=\epsilon_i{}^{jk}D_jv_k.
\end{equation}
Exact constant-curvature index algebra gives
\begin{equation}\label{eq:curl-squares}
 \boxed{
 \mathcal C_2^\dagger=\mathcal C_2,\qquad
 \mathcal C_2^2=-D^2+3;
 \qquad
 \mathcal C_1^\dagger=\mathcal C_1,\qquad
 \mathcal C_1^2=-D^2+2.}
\end{equation}
The tensor curl preserves symmetry, trace and divergence.  The curvature
terms are fixed by
$[D^n,D_i]h_{nj}=3h_{ij}$ and $[D^n,D_i]v_n=2v_i$ in our convention.

Using \eqref{eq:curl-squares}, the Euclidean Weyl-graviton operator of
Ref.~\cite{BeccariaBekaertTseytlin} continues to the local Lorentzian
factorization
\begin{equation}\label{eq:local-tt-factors}
 \boxed{
 B_{\rm TT}=P_-P_+,qquad
 P_\pm=\partial_t^2+(\mathcal C_2\pm1)^2
       =\partial_t^2-D^2+4\pm2\mathcal C_2.}
\end{equation}
Each factor has wave principal part and a first-order curl term.  A local
extension to spatial two-tensors is normally hyperbolic, its TT subspace is
reducing, and the restriction supplies the reduced TT Green operators.  The
transverse-vector branch is
\begin{equation}\label{eq:vector-wave-factor}
 P_A=\partial_t^2+\mathcal C_1^2
    =\partial_t^2-D^2+2,
\end{equation}
and is normally hyperbolic as well.  Closure of Green-hyperbolic operators
under composition and direct sum then makes the reduced physical TT-plus-vector
system Green hyperbolic \cite{BaerGreenHyperbolic}.  For
$B_{\rm TT}=P_-\circ P_+$, the ordered Green operator is
$G_{P_+}^{\pm}G_{P_-}^{\pm}$.

The local factors interchange the two frequencies between tensor
helicities.  For polarization define the nonlocal spectral operators
\begin{equation}
 A_E=|\mathcal C_2|-1,qquad
 A_L=|\mathcal C_2|+1,qquad
 A_A=|\mathcal C_1|.
\end{equation}
On a TT harmonic labelled by $r\ge0$, $|\mathcal C_2|=r+3$, so the
frequencies are $r+2$ and $r+4$.  On the metric-derived vector branch
$|\mathcal C_1|=r+2$ with $r\ge1$: the $r=0$ band consists of Killing
vectors whose symmetrized gradient vanishes.  Hence the physical $A$ tower
starts at energy three, not two.  At compact energy $N$ the exact
parity-complete multiplicities are
\begin{equation}\label{eq:field-eal-dictionary}
 \boxed{
 d_E(N)=2(N-1)(N+3),\quad
 d_A(N)=2(N-1)(N+1),\quad
 d_L(N)=2(N-3)(N+1),}
\end{equation}
on the respective ranges $N\ge2,3,4$.  Thus $E$, $A$, and $L$ are exactly
the lower TT, non-Killing transverse-vector, and upper TT field branches,
not merely character-matched modules.

The equation factorization alone does not determine the topology.  The
quadratic Weyl action does.  In the convention
$S=-\alpha_g\int C^2$, $\alpha_g>0$, its TT Pais--Uhlenbeck and reduced
vector symplectic forms give the positive residue magnitudes
\begin{equation}\label{eq:branch-residues}
 R_E=R_L=4|\mathcal C_2|,\qquad
 R_A=2(A_A^2-4),
\end{equation}
with signs $(+,-,-)$.  In particular $R_A$ has elliptic order two, not
zero.  The positive-frequency coordinate
\begin{equation}
 z_\alpha=\frac1{\sqrt2}
 \left[(R_\alpha A_\alpha)^{1/2}q_\alpha
 +i(R_\alpha A_\alpha^{-1})^{1/2}p_\alpha\right]
\end{equation}
therefore induces the Cauchy space
\begin{equation}\label{eq:branch-sobolev-spaces}
 \boxed{
 (H^1_{\rm TT}\oplus L^2_{\rm TT})_E
 \oplus(H^{3/2}_{\rm T}\oplus H^{1/2}_{\rm T})_A
 \oplus(H^1_{\rm TT}\oplus L^2_{\rm TT})_L.}
\end{equation}
The normalized metric waves obey
$2R_\alpha(N)N|N_\alpha|^2=1$ at every level.  Together with
\eqref{eq:field-eal-dictionary}, this proves the following result.

\begin{theorem}[Cylinder Cauchy--Sobolev realization]
\label{thm:cauchy-sobolev}
The positive-frequency harmonic transform from the reduced branch Cauchy
space \eqref{eq:branch-sobolev-spaces} extends by density to a Krein-unitary
isomorphism onto the one-particle energy-mode completion $\mathscr H_1$ of
Section~\ref{sec:energy-mode-completion}.  It intertwines the field-induced
signs with
$\mathfrak J_1=+1_E\oplus(-1_A)\oplus(-1_L)$.
The reduced TT Bach and vector operators are Green hyperbolic through
\eqref{eq:local-tt-factors} and \eqref{eq:vector-wave-factor}.
\end{theorem}

For raw TT data $(h_0,h_1,h_2,h_3)$, branch extraction uses
$(4|\mathcal C_2|)^{-1}$ and defines a mixed pullback graph norm.  We do not
replace it by an unsupported product-Sobolev norm.

\subsection{Curved auxiliary identities and the Weyl-curvature route}

\paragraph{Editorial note for the major-revision split.}
The propositions and causal-transport theorem in this subsection are the
mathematical core of the planned causal companion paper.  The intervening
coefficient-search chronology, discarded witness families, rank screens and
implementation ledgers are retained here to keep this archival monolith
reproducible, but are marked for the computational supplement in the split
described by \path{notes/conformal-paper-split-roadmap.md}.
\label{sec:auxiliary-green-witness}

Theorem~\ref{thm:cauchy-sobolev} concerns the reduced physical branches.
The all-degree curved operator statement and local BV equivalence are obtained
through an ordinary-derivative realization.  On the complete trace-free metric bundle
define
\begin{equation}\label{eq:exact-gauge-companion}
 T=\Box\delta-\frac13d\delta^2+\frac R3\delta
 -\operatorname{Ric}\circ\delta
 +\frac13d\langle\operatorname{Ric},h\rangle .
\end{equation}
Exhaustive finite-order jet algebra on the homogeneous cylinder gives the
global natural-operator identities
\begin{equation}\label{eq:exact-ghost-biwave}
 \boxed{TK=\Box(\Box+2),\qquad B_{\rm lin}K=0.}
\end{equation}
The first operator factors as
$(\Box-\operatorname{Ric}+2)(\Box+\operatorname{Ric})$, since
$\nabla\operatorname{Ric}=0$ and $\operatorname{Ric}^2=2\operatorname{Ric}$.
Both factors are normally hyperbolic.  If
$H=B_{\rm lin}+\frac12KT$, then
$HK=\frac12K\Box(\Box+2)$ follows formally.  With the action normalization,
the graded self-adjoint witness blocks are
$(T,2\sharp^{-1},T^\sharp)$, so the actual field block of $QW+WQ$ is
$2H=2B_{\rm lin}+KT$.

A direct product of $H$ on every component of $S^2_0T^*M$ has not been
certified.  The local auxiliary construction instead uses the ordinary-derivative
conformal-gravity action of Ref.~\cite{MetsaevOrdinaryDerivative}, with
fields $(h_{ab},f_{ab},v_a)$.  Its modified de Donder conditions obey
\begin{equation}
 \delta C_{0a}=\Box\xi_{-2,a}-\xi_{0,a},\qquad
 \delta C_{2a}=\Box\xi_{0,a},\qquad
 \delta C_1=\Box\sigma .
\end{equation}
Thus the combined ghost principal symbol is
$g^{\mu\nu}\zeta_\mu\zeta_\nu I_9$.  Direct reconstruction from the exact
covariant action and gauge fermion gives the complete curved Hessian
$E_{\rm aux}$, gauge map $K_{\rm aux}$, companion, and four-row operators.
In a canonical covariant derivative normal form they obey
\begin{equation}
 E_{\rm aux}^{\sharp}=E_{\rm aux},\qquad
 E_{\rm aux}K_{\rm aux}=0,\qquad
 Q_{\rm aux}^2=0,
\end{equation}
and
\begin{equation}
 Q_{\rm aux}W_{\rm aux}+W_{\rm aux}Q_{\rm aux}=P_{\rm aux},
 \qquad W_{\rm aux}^{\sharp}=W_{\rm aux},
 \qquad P_{\rm aux}^{\sharp}=P_{\rm aux}.
\end{equation}
The coefficient reconstruction passes $630/630$ exhaustive order-three and
order-four cylinder jet vectors and globalizes by homogeneity.  These exact
operator identities are what the repository flag
\texttt{curved\_operator\_identity} records.  They must not be confused with
normal hyperbolicity of the field block.

Indeed, the latter scalar-wave target is impossible for this field and gauge
content.  At the null covector $\zeta=(1,1,0,0)$ exact nonzero-minor
certificates give
\begin{equation}\label{eq:null-rank-no-go}
 \boxed{\rank E_2(\zeta)=11,\qquad \rank K_1(\zeta)=9.}
\end{equation}
For every invertible pointwise fibre map $J$, one has
$\rank(J^{-1}E_2)=11$.  For every first-order companion $C$, one has
$\rank(K_1C_1)\leq9$.  At a null covector a scalar normally hyperbolic symbol
would vanish, so the two terms cannot cancel.  This proves a no-go theorem
for the proposed $24$-component pointwise scalar-wave witness, not for Green
hyperbolicity.

The same symbol calculation explains the rank defect.  With
$N_2=J_{\rm act}^{-1}E_2$,
\begin{equation}\label{eq:helicity-two-symbol-quotient}
 \im N_2/\im K_1\cong\cH_{+2}\oplus\cH_{-2},
 \qquad \dim_\C(\im N_2/\im K_1)=2.
\end{equation}
In a real basis this quotient is represented by $f_{22}-f_{33}$ and
$f_{23}$; the transverse $SO(2)$ generator squares to $-4I_2$.  The
linearized Weyl symbol descends to the Hessian-reduced double quotient
\begin{equation}\label{eq:reduced-Weyl-symbol}
 \frac{F/\ker N_2}{N_2^{-1}(\im K_1)/\ker N_2}
 \longrightarrow
 \frac{\im\W_2}{\W_2(N_2^{-1}(\im K_1))}
\end{equation}
and, in the certified quotient bases, is exactly
\begin{equation}
 \boxed{\sigma(\W)_{\rm red}=\frac14I_2.}
\end{equation}
Both sides of \eqref{eq:reduced-Weyl-symbol} have dimension two.  We do not
claim $\ker\W_2=\im K_1$ on the unreduced $24$-field fibre.

The local BV-canonical change
\begin{equation}
 \eta=\xi_0-d\sigma,\qquad
 \widehat f=f-A_g^{-1}G^b(g,b)
\end{equation}
isolates $Qv=-\eta$.  The auxiliary mass map and inverse are pointwise,
\begin{equation}
 A_g(\phi)=-\frac12(\phi-g\operatorname{tr}_g\phi),\qquad
 A_g^{-1}(s)=-2s+\frac23g\operatorname{tr}_g s,
\end{equation}
and the transformation of every antifield row is generated by the same
local BV functional.  The changes, their inverses, and this elimination
preserve support.

\begin{proposition}[Curved auxiliary identities and support-local SDR]
\label{thm:auxiliary-green-witness}
The ordinary-derivative tensor--tensor--vector BV complex on the Lorentzian
conformal cylinder has exact curved four-row operator identities.  The local
BV-canonical change above conjugates the actual curved differential to the
metric complex plus three generalized-auxiliary arrows.  The resulting maps
satisfy
\begin{equation}
 p\circ i=\id,\qquad i\circ p-\id=Qk+kQ
\end{equation}
on every field, ghost, antifield, trace/Weyl, and nonminimal row.  Their
entries are finite differential maps or pointwise inverses, so they preserve
compact, spacelike-compact, and unrestricted smooth support.  The same
canonical split gives the action-level presymplectic comparison
\begin{equation}
 i^*\omega_{\rm aux}-\omega_{\rm met}=d\beta+Q\gamma
\end{equation}
and equality of the induced closed-$S^3$ Cauchy pairings on cohomology.
Thus the repository flags \texttt{curved\_deformation\_retract} and
\texttt{curved\_current\_comparison} are true.  No local
conformal-Killing projector is used.
\end{proposition}

The exact physical symbol quotient selects curvature prolongation as the
current causal route.  For an algebraic Weyl tensor $\Psi$ define its
electric and magnetic STF parts by
\begin{equation}
 E_{ij}=\Psi_{0i0j},\qquad
 B_{ij}=\frac12\epsilon_i{}^{kl}\Psi_{0jkl}.
\end{equation}
The ten-dimensional algebraic bundle
$\operatorname{STF}_2(S^3)\oplus\operatorname{STF}_2(S^3)$ does not close by
itself at first order.  If
\(V_{\mu,\nu,\sigma}=\nabla^\rho\Psi_{\mu\rho\nu\sigma}\), its independent
components form a sixteen-dimensional Cotton slot, naturally two trace-free
$3\times3$ matrices.  Exact contraction of the Cotton definition, Bach row,
and Hodge-dual compatibility row gives a $34\times26$ first-order coefficient
table on $(E,B,V)$ whose temporal matrix has rank $26$ and leaves eight
constraints.  An exhaustive comparison on all
$10\binom{6}{4}=150$ Weyl two-jets proves equality with the curved covariant
Bianchi--Bach equations and globalizes by cylinder homogeneity.  Thus
\texttt{curved\_EB\_equations} and
\texttt{curved\_EB\_first\_order\_closure} are exact.

This result is not yet a hyperbolicity theorem.  The canonical unadjusted
choice of 26 evolution rows has characteristic polynomial
\begin{equation}\label{eq:EB-unadjusted-characteristic}
 \frac1{144}\lambda^2(\lambda^2-1)^8
 (4\lambda^2-1)^2(3\lambda^2+1)^2,
\end{equation}
so the pair $\lambda=\pm i/\sqrt3$ proves that constraint additions are
necessary.  Exact constant-curvature symbol algebra nevertheless gives
\begin{equation}\label{eq:EB-principal-constraint}
 \operatorname{div}\operatorname{curl}_2
 =\frac12\operatorname{curl}_1\operatorname{div}.
\end{equation}
The constraint-adjusted square system is now exact at the level required for
the Cauchy problem.  Its first twenty rows are covariant combinations, while
its final six differ from the vector Bach rows by the two vector constraints
$a,c$.  Consequently the pointwise row modules differ by rank six.  This is
not a solution-space defect: differentiating the primary $q,r$ constraints
and using the exact vector Bach rows generates $a=c=0$.  Conversely, the
adjusted rows together with all fourteen constraints recover every covariant
row.  Thus the two presentations generate the same formally integrable
differential ideal.

The adjusted $26$-state system has a positive symmetrizer and causal speeds
$0,\pm\frac12,\pm1/\sqrt3,\pm1$.  Its complete sourced subsidiary identity
includes the unit-$S^3$ curvature commutator correction, and compatible
sources preserve all fourteen constraints.  This proves the three flags
\begin{center}\small
\path{curved_EB_symmetric_hyperbolicity},\par
\path{curved_sourced_constraint_identity}, and\par
\path{curved_constraint_propagation}.
\end{center}
The unadjusted raw subsidiary
system retains an imaginary longitudinal pair; the certified local
constraint addition, not a hidden spectral projection, removes it.

The global solution-space audit is independent of that analytic repair and
is already exact.  The exhaustive jet equivalence, the smooth cylinder BGG
theorem, and the symbolic all-energy metric preimages give
\begin{equation}
 \ker(C_1^\sharp)\cap\ker(C_1^\sharp\star)
 \cong E\oplus A\oplus L
\end{equation}
with both chiralities.  Symbolically,
$\dim\mathcal C_n-2\rank C_1^\sharp=6n^2-14$, equal to the parity-complete
$E/A/L$ multiplicity, and the exact rational characters agree.  The Cotton
definition has an algebraic identity coefficient, so its graph adds no
second solution copy.  Thus the flag
\begin{center}\small
\path{EAL_curvature_spectrum_match}
\end{center}
is true; the levels $2$ through $6$ are retained only as regressions, not as
the proof.

The target is not a second residual Chevalley--Eilenberg calculation.  That
calculation and its pairing are already complete.  It is a local chain of
pairing-compatible, support-preserving quasi-isomorphisms
\begin{equation}\label{eq:covariant-causal-bridge}
 \Gamma_c(\mathcal C_{\rm met})[1]
 \simeq\Gamma_c(\mathcal C_{\rm aux})[1]
 \simeq\Gamma_c(\mathcal C_{\rm prol})[1]
 \simeq\Gamma_{sc}(\mathcal C_{\rm prol})
 \simeq\mathcal C_\Sigma
 \simeq\mathcal W_+\oplus\mathcal W_-
 \quad\text{with residual endpoints}.
\end{equation}
The curvature variable must be adjoined as a contractible prolongation,
for example through $\widehat\Psi=\Psi-C_1h$, rather than replacing the
metric potential by curvature.  Its inclusion, projection, and homotopy
must contain every field, equation, identity, antifield, trace/Weyl,
nonminimal, and first-order-reduction row and use only finite differential
operators and pointwise inverses.

Two further support-local subtheorems are now exact.  First, adjoining the
graphs $\Psi=C_1h$ and $c=\operatorname{div}\Psi$ with their cotangent lifts
gives a BV-canonical contractible SDR.  Second, the curvature rows form the
autonomous compatibility complex
\begin{equation}
 U_{26}\xrightarrow{(L,K)}F_{26}\oplus C_{14}
 \xrightarrow{(-R,S)}I_{14},
 \qquad SK=RL,
\end{equation}
together with its formal cotangent-adjoint complex.  A zeroth-order backward
witness makes $QW+WQ$ diagonal with the symmetric-hyperbolic blocks $L$ and
$S$ and algebraic graph identities.

These subtheorems cannot be direct-summed: the autonomous curvature complex
carries the nonzero $E/A/L$ module and would duplicate physical cohomology.
The state/equation part of the required mapping-cylinder chain map has been
constructed.  It supplies the explicit local operator $A_{\rm eq}$ satisfying
\begin{equation}
 (L,K)T=A_{\rm eq}E_{\rm aux},\qquad
 T=(C_1,\operatorname{div}C_1).
\end{equation}
The first square is verified on all $10\binom{8}{4}=700$ metric four-jets.
The formerly reported order-zero identity-square defect was a coordinate
artifact.  The common-Rees diagnostic had composed the curved attachment with
the flat-Fourier equation projection, whereas the fibre-identified curved
cotangent row uses
\begin{equation}
 p_E=(D^{-1}S_h^\sharp D,1,0)
\end{equation}
and the corresponding identity projection $p_I$.  Reconstructing these maps
coefficientwise gives
\begin{equation}
 A_{\rm eq}=A_{\rm core}p_E,
 \qquad B_{\rm id}=B_{\rm core}p_I,
 \qquad N_{\rm curv}A_{\rm eq}=B_{\rm id}C_{\rm aux}.
\end{equation}
The last equality follows from the exact metric-core square together with the
equation-to-identity block $p_IC_{\rm aux}=C_{\rm core}p_E$ of the certified
all-row curved retract.  Thus the old rank-four matrix is retained only as a
scoped mixed-coordinate diagnostic and is not an operator obstruction.  The
preceding square $TK_{\rm aux}=0$ remains independently exhaustive on all
diffeomorphism and Weyl-scalar jets and all four boost ghosts.

Substitution into the complete sixteen-block cotangent mapping cylinder now
proves coefficientwise
\begin{equation}
 Q_{\rm prol}^2=0,
 \qquad PI=1,
 \qquad IP-1=Q_{\rm prol}H+HQ_{\rm prol},
\end{equation}
as well as odd BV cyclicity.  Every field, equation, identity, cotangent and
contractible cone row is enumerated, and $I,P,H$ contain only finite-order
differential operators and pointwise inverses.  Hence
\texttt{support\_local\_prolongation\_retract} and
\texttt{prolonged\_BV\_operator\_identity} are true.  The prolonged current
certificate now binds the corrected curved core-chain provenance and proves
$I^*\omega_{\rm prol}-\omega_{\rm aux}=d\beta+Q\gamma$ off shell, with
$\beta=\delta Y_I$ and $\gamma=0$.  Hence
\texttt{prolonged\_current\_comparison} is true.  Cyclicity of the direct
causal homotopies makes
$\Delta_\Lambda=\Lambda_{{\rm prol},+}-\Lambda_{{\rm prol},-}$ graded
antisymmetric; the common algebraic and trace contractions cancel from this
difference.  Combining the resulting parent Green/current identity with the
prolonged-to-auxiliary and auxiliary-to-metric $d+Q$ improvements proves that
the causal and Cauchy-current pairings agree on cohomology, with all-energy
normalization $+I_E\oplus(-I_A)\oplus(-I_L)$.

The two legacy analytic exceptions
are $E_{\rm aux}+KC$ on the auxiliary field row and its cotangent copy; they
retain the scalar-wave null-rank obstruction.  Thus mere direct-sum
conjugation does not produce the superseded canonical prolonged witness.
They do not obstruct the direct tractor causal homotopy or its pairing
transport.

Two exact fail-closed diagnostics further delimit this open construction.
For $q=g^{-1}(\zeta,\zeta)$, the generic prenormal principal symbol
$P_2=J_{\rm act}^{-1}E_2+K_1C_1$ obeys
\begin{equation}\label{eq:prenormal-minimal-polynomial}
 (P_2-qI_{24})^2=0,
 \qquad
 (2qI_{24}-P_2)P_2=P_2(2qI_{24}-P_2)=q^2I_{24}.
\end{equation}
Its Smith multiplicities are $6/12/6$, corresponding respectively to
algebraic, wave and biwave factors.  This is a principal-symbol identity, not
a Green theorem.  Indeed, the naive frozen lower-order completion
$2qI-P$ has nonzero remainders in orders zero, one and two.  This rejects only
that literal completion; it does not rule out corrected first-order factors
or a full covariant composition.

The complete invariant lower-order ansatz makes the next gate finite.  The
$SO(3)$ cylinder-holonomy decomposition gives
\begin{equation}\label{eq:lower-factor-ansatz-dimensions}
 \dim D_0=38,
 \qquad
 \dim D_1=93,
 \qquad
 D=D_{\rm naive}+X_1^\mu\nabla_\mu+X_0.
\end{equation}
The exact cubic divisibility equations for both $DP$ and $PD$ have a
simultaneous $45$-parameter solution family.  There is therefore no cubic
obstruction.  Restoring the two independent factor splittings and all
algebraic coefficients gives
\begin{equation}\label{eq:lower-factor-full-variable-count}
 45+2\cdot93+5\cdot38=421
\end{equation}
nonlinear unknowns after the cubic gate.  The scalar principal normalization
$qI_{24}$ loses no invariant solutions, since a parallel invertible
$qH/qH^{-1}$ pair may be redistributed between the two factors.
The complete $421$-variable problem is now assembled as an exact sparse
quadratic symmetrized-PBW system.  Its order-zero, order-one and order-two
row counts are respectively $240$, $960$ and $2484$.  At order two, the
$2484\times190$ matrix multiplying the algebraic variables has rank $100$,
so its cokernel has dimension $2384$ and $90$ algebraic variables remain
free.  Exact Schur projection gives $2130$ nonzero polynomial constraints,
$365$ up to scale, and no constant polynomial obstruction.  These projected
equations and the lower orders have not been solved.

A convention audit corrected the covector commutator: its raised curvature
index uses the spatial projector, not a four-dimensional Kronecker delta, so
all mixed time--space curvature components vanish.  Independent
coordinate-jet commutator, Weitzenb\"ock and focused $\Box^2$ composition
tests now pass.  Exact triangular symmetrized-jet PBW inversion exhausts all
$1680$ four-jet basis elements, passes $504$ ordered-word round trips and
certifies associativity, so the quadratic-factor composition backend is
ready.  A naive transpose/reversal of the already sorted
$\Box^2$ component table nevertheless leaves a $48$-entry defect.  This is
not an independent primal-composition counterexample: the sorted coefficient
matrices no longer individually carry their derivative-index slots.  It does
show that the naive table adjoint is invalid.  The required pairing-aware
adjoint has now been constructed on the symmetrized-PBW coefficient tensors.
With the action pairing it verifies

\begin{equation}\label{eq:general-factor-sharp-reduction}
 P^\sharp=P,\qquad D_{\rm naive}^\sharp=D_{\rm naive},\qquad
 (DP)^\sharp=PD.
\end{equation}

Requiring $D^\sharp=D$ reduces the invariant first-order family from $93$ to
$44$ dimensions and its algebraic family from $38$ to $24$.  The resulting
cubic matrix has size $11520\times137$, rank $116$, and a $21$-dimensional
kernel.  Taking $R_-=L_+^\sharp$ and $R_+=L_-^\sharp$ makes the right
factorization the adjoint of the left rather than an independent system.
Restoring the left splitting and algebraic coefficients leaves a
$214$-parameter nonlinear problem.  Its complete order-two gate has $1242$
rows; the $1242\times100$ algebraic matrix has rank $52$, leaving $48$ free
algebraic directions.  Exact cokernel projection gives $1050$ nonzero
constraints, $179$ up to scale, with no constant polynomial obstruction.
The projected quadratic constraints and orders one and zero remain unsolved
and unpromoted.  All $179$ normalized constraints have maximal degree two.
Their $1497$ quadratic monomials have coefficient rank $124$; the resulting
$55$ left-kernel combinations have zero affine part.  Hence exact
vector-space reduction eliminates no variable, and the next step is a
genuinely nonlinear ideal calculation rather than further linear algebra.

The complete bare-wave subfamily can nevertheless be decided exactly.  If
one factor in each of $DP$ and $PD$ is the literal rough $Box$, the four
inner/outer orientations give rational systems with $159$ unknowns and

\begin{equation}\label{eq:bare-box-factor-obstruction}
 \rank A=159,
 \qquad
 \rank[A\mid b]=160.
\end{equation}

All four have the same one-row left-null certificate: the coefficient of
$\nabla_{(0}\nabla_{1)}$ mapping $f_{01}$ to $f_{00}$ is $-8$ in the
required right-hand side and vanishes in every correction column.  This is
a no-go only for the bare-$\Box$ subfamily.  When both factors have nonzero
first-order coefficients, the quadratic product $A_-A_+$ changes the
order-two equations.  In fact, a support-minimal rational zero-sum split
using two invariant directions repairs this row together with its full
$SO(3)$ orbit $f_{0i}\to f_{00}$.  The old bare-wave obstruction therefore
does not extend to the general family.  Fixing that minimal split nevertheless
produces a new scoped order-two obstruction after all $190$ invariant
algebraic variables are admitted:

\begin{equation}\label{eq:fixed-minimal-split-rank-no-go}
 \rank A=100,
 \qquad
 \rank[A\mid b]=101.
\end{equation}

Its one-row left-null witness is the coefficient of
$\nabla_{(0}\nabla_{1)}:h_{22}\to f_{01}$, whose required value is $16$ and
whose $190$ variable coefficients vanish.  This rules out only that fixed
minimal split, not other first-order splits, the complete $421$-variable
system, or the $214$-parameter sharp subfamily.
Thus the cubic gate proves neither a factorization nor a Green theorem.

As a further exact screen, multiplying the $179$ reduced sharp quadrics by
constants and all $114$ degree-one variables gives a sparse Macaulay matrix
with $136585$ monomial rows and $20585$ multiplier columns.  A deterministic
minor over a prime avoiding every input denominator yields the rigorous
characteristic-zero degree-three bound
\begin{equation}\label{eq:sharp-macaulay-rank-bounds}
 12861\leq \rank_{\mathbb Q} M_3\leq14136.
\end{equation}
The upper bound is $124\cdot114$, from the exact dimension of the quadratic
coefficient space.  The full rational ranks have not been computed, so this
screen decides neither membership of a nonzero constant in the ideal nor the
low-degree quotient dimensions.  It promotes no no-go or Green flag.

The exhaustive relative-incidence ledger contains nine possible odd-adjoint
pairs.  No single pair yields reciprocal coupling, and the smallest physical
two-way candidate is the pair $4+5$.  On the ordered core
$(M_{\rm aux},X_U,Y_U^\sharp)$ it gives
\begin{equation}\label{eq:minimal-relative-saddle}
 \begin{pmatrix}
 E_{\rm aux}+KC & R & S\\
 S^\sharp E_{\rm aux} & L_{26} & 0\\
 R^\sharp E_{\rm aux} & 0 & L_{26}^\sharp
 \end{pmatrix}.
\end{equation}
The finite-block Schur formula is exact, but its Schur complement contains
the Green operators of $L_{26}$ and $L_{26}^\sharp$ and is therefore
nonlocal.  Conversely, the unreduced saddle is order two, so the certified
first-order symmetric-hyperbolic theorem does not apply to it.  Moreover, its
natural local instantiation
\begin{equation}\label{eq:minimal-saddle-instantiation}
 A_F=p_F A_{\rm eq},
 \qquad S=A_F^\sharp,
 \qquad R=A_F^\sharp J_U
\end{equation}
has an exact balanced Douglis temporal rank obstruction:
\begin{equation}\label{eq:minimal-saddle-rank-no-go}
 \rank\sigma_{\rm DN}(dt)\leq107<116,
 \qquad \operatorname{defect}\sigma_{\rm DN}(dt)\geq9.
\end{equation}
The temporal leading coefficient of the characteristic determinant therefore
vanishes and no positive temporal symmetrizer exists for this realization.
This is a no-go only for the smallest pair-$4+5$ ansatz with
\eqref{eq:minimal-saddle-instantiation}; it does not exclude larger relative
witnesses or an additional support-local first-order prolongation.

The expanded-relative incidence audit exhausts all nine odd-adjoint pairs
and the $162$-dimensional Hom family determined by the declared cylinder
$SO(3)$ multiplicities.  It finds exactly
three minimal reciprocal two-pair candidates, labelled $1+6$, $1+7$ and
$2+7$.  Even using all reciprocal curvature partners, the maximum cross-rank
on each central $24$-component auxiliary block is $21$.  The three missing
directions are rotation scalars.  Hence any ansatz which makes the entire
central auxiliary diagonal subprincipal has temporal rank defect at least
three.  A viable expanded witness must retain or locally prolong these scalar
directions.  This is a scoped design theorem; coefficient tables,
characteristics and a symmetrizer remain open, and no flag is promoted.
An independent coefficientwise calculation constructs the three rotation
generators in the actual component bases of all sixteen blocks.  The nine
relative Hom equations have nullities
$4,18,4,36,14,14,22,36,14$, whose sum is $162$.  The explicit pair-$1+6$
maps, the temporal $K$ and $N_{\rm curv}^{\sharp}$ coefficients, the induced
vector projector, and the retained scalar diagonal all have zero
rotation-generator intertwining defect.

The three uncovered scalar coordinates are exactly $h_{00}$, $f_{00}$ and
$v_0$.  The restriction of the existing $KC$ coefficient to these directions
is triangular of rank three and determinant $-1$.  For the pair-$1+6$
candidate, explicit finite-order coefficient maps give the exact numerator
identity

\begin{equation}\label{eq:relative-pair16-numerator}
 K R_1 N_{\rm curv}^{\sharp}R_6^{\sharp}=-\Pi_{\rm vector}.
\end{equation}

To pass from this numerator to the saddle Schur term $BD^{-1}C$, the actual
$92$-component curvature temporal diagonal is needed.  In the normalized curvature
equations and their formal-adjoint rows,
\begin{equation}\label{eq:relative-curvature-temporal-diagonal}
 D(dt)=\operatorname{diag}(I_{26},-I_{40},-I_{26}),
 \qquad D(dt)^{-1}=D(dt).
\end{equation}
Here the first-order compact-support convention
$N_{{\rm curv},\mu}^{\sharp}=-N_{{\rm curv},\mu}^{T}$ is used both in the
middle diagonal and in $C$.  Consequently the exact pair-$1+6$ Schur term is
$BD^{-1}C=+\Pi_{\rm vector}$, and the field Schur block is
$E_{{\rm aux},2}+D_{\rm scalar}-\Pi_{\rm vector}$.  Retaining the rank-three scalar diagonal gives
an assembled $116\times116$ temporal Douglis symbol of exact rank $116$ and
determinant one.  This closes the temporal-invertibility gate only.  The
scalar diagonal has not been lifted to a cyclic all-row witness, and the
three spatial first-order coefficients of $R_6^{\sharp}$ have not yet been
constructed.  Hence the arbitrary-covector relative symbol itself, its
characteristic polynomial, positive symmetrizer,
lower-order completion and every-degree Green identities remain open.  Thus
no Green-theoretic flag follows from
\eqref{eq:relative-pair16-numerator}--\eqref{eq:relative-curvature-temporal-diagonal}.

The subsequent complete $SO(3)$-equivariant audit makes this boundary
sharper.  The actual $R_6^\sharp$ coefficient equations have a
$22$-dimensional temporal commutant and a $46$-dimensional spatial-vector
commutant.  After fixing the temporal normalization above, all $46$ spatial
parameters were retained.  Nevertheless, for the intrinsic aligned pencil
$Q(z)=P_{\rm weighted}((-z,1,0,0))$ the vectors
$a_0=2f_{23}$ and $a_1=h_{23}$ obey
\[
 Q(1)a_0=0,\qquad Q(1)a_1+Q'(1)a_0=0,
\]
and the two parameter-sensitivity maps have exact rank zero.  Hence every
regular member of this fixed-temporal pair-$1+6$, cyclic $-2\Pi$ family
retains a length-two elementary divisor.  This rules out a semisimple
faithful strong linearization, and therefore a positive symmetric-hyperbolic
symmetrizer, within this family.  It is not a no-go for another relative
incidence, a changed temporal normalization, a larger witness, or generalized
Green hyperbolicity.  This screen by itself promotes no causal or final
covariant $H^4$ flag; the later curved tractor transfer supplies the causal
homotopy by a different construction.

This non-semisimplicity is not itself a Green obstruction.  An independent
BV-row audit identifies $2f_{23}$ with the already contractible shifted
auxiliary summand, while $h_{23}$ has nonzero Weyl helicity two; the
extension splits after the certified support-local auxiliary shift and dies
as an extension on $Q$-cohomology.  Moreover, on the aligned physical
principal quotient the exact block is
\[
 \begin{pmatrix}qI_2&0\\4\rho^2I_2&qI_2\end{pmatrix},
 \qquad q=\rho^2-\tau^2,
\]
whose abstract operator form $\left(\begin{smallmatrix}D&0\\R&D\end{smallmatrix}\right)$
has the two-sided same-sided Green inverse
$\left(\begin{smallmatrix}G&0\\-GRG&G\end{smallmatrix}\right)$.  Thus the
remaining task is a projector-free local exact-extension construction, not
removal of the physical Jordan chain.  Refining the bundle support graph by
the evolution/subsidiary equation-dual split leaves a reciprocal rank-$34$
$(h,f,C^\sharp)$ component; the other curvature tails are triangular
singletons.  This aligned principal calculation does not yet promote the
full physical biwave or prolonged Green flags, since the local curvature
quotient and lower-order intertwining remain open.
On this aligned channel the certified BV shift is explicit as well: writing
$L=1-z^2$, it conjugates the complex block to
$\operatorname{diag}(L^2,-1)$, while the witness remains the closed block
$\left(\begin{smallmatrix}L&0\\4&L\end{smallmatrix}\right)$ with exact
same-sided inverse
$\left(\begin{smallmatrix}G&0\\-4G^2&G\end{smallmatrix}\right)$.  No inverse
curl, Laplacian, TT projector or helicity projector occurs in this formula;
what remains open is its support-local realization in the complete
$116$-row exact extension.
On the exact TT-plus-shifted-auxiliary operator subcomplex one can therefore
promote the scoped statements that the physical biwave block is Green
hyperbolic and its Jordan extension is causal: the restricted witness is
$\operatorname{diag}(B_{\rm TT},1,B_{\rm TT},1)$ and the corresponding
restricted $Q\Lambda_\pm+\Lambda_\pm Q=1$ identity is exact.  These are not
the complete prolonged Green statements.  The rank-$4$ vector sector is
contracted below, but the rank-$14$ equation-cone extension and a
support-local map from arbitrary sources into its physical biwave carrier
remain open.
The first alternative-incidence screen is also exact: pairs $1+7$ and
$2+7$ have a nonzero joint intrinsic-sensitivity image of rank $16$.
Their smallest temporally regular slices have determinant $8$ and entirely
real causal roots, but are non-semisimple at speeds $0$ and $\pm1$; the first
directly sensitive spatial correction does not repair either slice.  This is
not a no-go for the complete alternative families.
The canonical $16$-parameter subfamily spanning the full intrinsic
sensitivity image is nevertheless uniformly non-semisimple: the zero-root
valuation and kernel-dimension bounds are respectively $40$ versus at most
$33$ for pair $1+7$, and $48$ versus at most $47$ for pair $2+7$.  This
still does not rule out the raw $122$-parameter families or generalized
Green extensions.
A projector-free differential-module analysis reduces the reciprocal
rank-$34$ component further.  The six vector-gauge columns together with the
$q^\sharp/r^\sharp$ subsidiary coordinates present a rank-$12$ invariant
submodule with an exact recursive Green inverse.  The rank-$22$ quotient
contains a closed rank-$8$ symmetric-hyperbolic constraint quotient and an
unresolved rank-$14$ field cokernel.  Since
$(C_1,\operatorname{div}C_1)K=0$, curvature descends locally to this
cokernel; constructing its induced biwave intertwiner and inverse is the
remaining central block problem.  The raw off-diagonal ideal is not
nilpotent, so the inverse does not follow from a naive finite Neumann series.
The remaining rank-$4$ vector field singleton, together with its mandatory
primal/cotangent closure of rank $16$, is already the shifted
$(\eta,v,v^\sharp,\eta^\sharp)$ contractible summand.  Replacing its witness
locally gives $P_{\rm vec}=I_{16}$, $G_+=G_-=I_{16}$ and
$\Lambda_+=\Lambda_-=W$, with both homotopy defects zero and vanishing causal
propagator.  Thus the unresolved analytic content of this filtration is the
rank-$14$ curvature field cokernel and its all-row insertion.
For reproducibility, the rank-$14$ handoff is persisted as a compact
content-addressed manifest: it records the ordered quotient and filtration
maps, full $26$-state evolution and lower-order tables, all fourteen source
constraints and subsidiary matrices, symmetrizers, BV row layout, formal
adjoints and current conventions by authoritative path and hash rather than
duplicating the large matrices.
The rank-$14$ quotient has a projector-free local presentation.  Raw
$T=(C_1,\operatorname{div}C_1)$ has rank $5$ and kernel rank $9$, while the
compatible-source kernel has rank $12$, but $\operatorname{im}T$ is not
contained in it: the defect ranks are $3$ generically and $1$ at the aligned
null covector.  Thus $K_{12}/I_5$ is undefined.  This is distinct from the
exact null Weyl split into rank-three gauge preimages and two helicities, on
which the reduced Weyl map is $\frac14I_2$.

The exact replacement is the equation cone
\begin{equation}
 \widehat T=(T,E_{\rm aux}),\qquad
 \widehat K=(K_{\rm state},-A_C),\qquad
 \widehat K\widehat T
 =K_{\rm state}T-A_CE_{\rm aux}=0,
\end{equation}
the constraint component of the certified chain square for
$E_{\rm curv}=(L_{\rm WC},K_{\rm state})$.  Hence the bridge is an
equation-level deformation retract using the auxiliary $f$ equation and all
forty curvature rows, not a direct state retraction.

The full cone, including its incoming gauge row, has degree ranks
$9\to24\to50\to49\to14$.  The actual curved projections $p_E,p_I$ make its
operator squares exact; the old rank-four defect used mixed coordinates.
Moreover the fifteen-entry curved attachment supplies the degree-$(-2)$
differential on the residual null PBW page
$\C^2\xrightarrow{d_{12}}\C^4\xrightarrow{d_{23}}\C^2$.
Both maps have rank two and an explicit homotopy gives
$h_{12}d_{12}=1$, $d_{12}h_{12}+h_{23}d_{23}=1$, and $d_{23}h_{23}=1$.
Thus the null filtered cone is algebraically exact.

This finite-page result is not a total local inverse: $DH+HD=1$ would require
$H_1K=I_9$ with $\rank K=5$ and, at the correctly typed upper endpoint,
$[N,B]H_4=I_{14}$ with $\rank[N,B]=13$.  The certified composite SDR instead
gives $P_{\rm alg}=DH_{\rm alg}+H_{\rm alg}D$ as a support-local cyclic
idempotent chain projector, with $P_{\rm end}=1-P_{\rm alg}$.  It contracts
$356$ prolonged components and retains the $30$-component metric--curvature
graph.  The endpoint has now also been reconstructed coefficientwise in the
ordered rows
\begin{equation}\label{eq:exact-thirty-row-endpoint}
 G_{\rm met}[5]\xrightarrow{K_{\rm met}}h[10]
 \xrightarrow{\bar B} \bar E_{\rm met}[10]
 \xrightarrow{C_{\rm met}} I_{\rm met}[5],
 \qquad \operatorname{ord}(K_{\rm met},\bar B,C_{\rm met})=(1,4,1).
\end{equation}
All curved lower-order coefficients are emitted.  Direct coefficient checks
give $Q_{\rm end}^2=0$, the two formal-adjoint identities and odd cyclicity;
the local graph maps satisfy $p_{\rm end}j_{\rm end}=1$ and strictly
intertwine $Q_{\rm prol}$ and $Q_{\rm end}$.  Thus the earlier $66\to30$
flat-Fourier hashes remain only a dimension/order/sign regression and are no
longer the source of the curved endpoint claim.

The endpoint equation coordinates already contain the fibre identification
$\bar B=J_{\rm met}^{-1}E_{\rm Bach}$.  Consequently the abstract minimal
witness coefficient $2\sharp^{-1}$ must not be inserted a second time.  The
coefficientwise canonical backward map is
\begin{equation}\label{eq:endpoint-canonical-backward-witness}
 W_0=(T_{\rm gf},1,T_{\rm gf}^{\sharp}),
 \qquad
 D_{\rm end}=Q_{\rm end}W_0+W_0Q_{\rm end}.
\end{equation}
With the superseded middle coefficient $2I$, the trace-free principal matrix
at $\zeta=dt$ has spectrum $1^4\oplus2^5$; this is the exact normalization
counterexample missed by the first textual guard.  With the identity in
\eqref{eq:endpoint-canonical-backward-witness}, exhaustive fourth-order
coefficients instead give
\begin{equation}\label{eq:endpoint-scalar-tracefree-symbol}
 \sigma_4(D_{\rm TF})(\zeta)
 =\bigl(g^{\mu\nu}\zeta_\mu\zeta_\nu\bigr)^2I_9.
\end{equation}
The ghost block is a triangular extension of
$\Box(\Box+2)I_4$ by the Weyl-scalar identity; the metric trace is likewise a
triangular identity extension.  Thus the sole primal analytic block is the
rank-nine $D_{\rm TF}=P_{\rm TF}D_MP_{\rm TF}$, together with its formal-adjoint
copy in the equation row.

The canonical curvature backward-map audit identifies a precise boundary.
The upper row lifts locally: there is a rank-four algebraic map $R$ with
\begin{equation}
 A_{\rm core}R=i_C B_{\rm core}.
\end{equation}
The analogous canonical middle lift would require
$T_{\rm core}S=A_F$.  But the exact factorization
$T_{\rm core}=J_{\rm WC}W_{EB}$ obeys
$\pi_{EB}J_{\rm WC}=1$, whereas $\pi_{EB}A_F=0$ and
$\rank A_F=5$.  Hence such an $S$ is impossible on the leading Douglis
page.  Equivalently, the two rank-five symbol images have zero intersection
on the certified generic, timelike, spacelike, null and temporal
representatives.  This is a no-go only for restricting the canonical
zeroth-order middle map to the endpoint graph, not for a relative witness.

Indeed, the rank-five $A_F$ incidence has an exact support-local cyclic lift
through the algebraic complement of the full $386$-component complex.  Its
two off-diagonal legs are nonzero and fixed by the $386\to66\to30$
projectors.  However the minimal projected saddle has the stronger exact
identity
\begin{equation}\label{eq:endpoint-minimal-saddle-nilpotence}
 P_{\rm end}L_{A_F}P_{\rm alg}L_{A_F}P_{\rm end}=0.
\end{equation}
Thus its Schur endpoint is unchanged: $S_{\rm end}=D-CB=D$.  The relative
incidence introduces no new endpoint obstruction, but it also cannot improve
Green invertibility of the exact thirty-row diagonal.  Larger relative
witnesses are not excluded.

The most direct same-bundle factorization of $D_{\rm TF}$ is now excluded
exactly.  Write
\begin{equation}\label{eq:endpoint-wave-factor-ansatz}
 L_\pm=P_{\rm tr}+\Box P_{\rm TF}
       +A_\pm^\mu\nabla_\mu+B_\pm .
\end{equation}
On $S^2_0T^*M\cong V_0\oplus V_1\oplus V_2$, the complete parallel
$SO(3)$-invariant trace-free-preserving families have dimensions nine at
first order and three at order zero.  A general nondegenerate wave-type
leading pair $qH_-,qH_+$ loses no solutions under the displayed scalar
normalization: the scalar target forces $H_-H_+=1$, and replacing the factors
by $L_-H_+$ and $H_+^{-1}L_+$ preserves their product and the complete
lower-order families.  Exact symmetrized PBW composition of
$L_-L_+-D_{\rm TF}$ has no order-four defect.  Order three has coefficient and
augmented rank nine and forces $A_+=-A_-$.  At order two the three sums
$B_-+B_+$ are fixed; the remaining $42$ quadratics admit twelve explicit
degree-two rational multipliers $w_i$ satisfying
\begin{equation}\label{eq:endpoint-factor-nullstellensatz}
 \boxed{\sum_i w_i f_i=1.}
\end{equation}
Hence no factorization of the form \eqref{eq:endpoint-wave-factor-ansatz}
exists, and orders one and zero are unreachable.  This is a no-go for the
complete parallel invariant two-wave-factor family, not for mixed-order or
enlarged systems, triangular Green extensions, or the equation-cone
curvature-to-metric route.  Its exact receipt is
\texttt{curved\_endpoint\_metric\_factorization\_no\_go.json}; the nonzero
multipliers are printed in
Appendix~\ref{app:endpoint-factorization-receipt}.

The parent-detour route is now complete at the causal-homotopy level.  For
the flat adjoint-tractor connection, the Yang--Mills detour complex of
Ref.~\cite{GoverSombergSoucek} has ranks $15\to60\to60\to15$ and a
degreewise normally hyperbolic parent witness.  Exact Kostant contraction
and finite homological perturbation terminate after two terms and compress
it to ranks $4\to9\to9\to4$.  In curved PBW normal form the former
$48$-entry commuting-derivative defect decomposes into a complete
ten-dimensional $SO(3)$-invariant Hom basis: two time-derivative, five
spatial-derivative, and three algebraic channels.  The curvature action on
the derivative suffix, form slot and graded adjoint-tractor slot reconstructs
all ten channels from the Kostant data, with no fitted coefficient, and
cancels all $48$ entries.  The one-form convention is fixed independently by
$-R_M^T$ on $\Lambda^1T^*$ together with $\operatorname{ad}_M=[M,\cdot]$;
the other three sign/transpose choices fail the exact harmonic-carrier gate.

The resulting curved differential maps $i_0,i_1,p_0,p_1,h_0,h_1$ are
finite-order and support local.  Their inclusion, projection, homotopy and
cyclic-adjoint defects all vanish.  Moreover the complete compressed parent
middle operator, not merely its principal symbol, satisfies
\begin{equation}\label{eq:curved-tractor-Bach-match}
 P_{\rm mid}^{\rm parent,comp}
 =-2S^T J_{\rm met}\,\bar B\,S
\end{equation}
coefficientwise at derivative orders zero through four.  Thus the parent
advanced and retarded homotopies transfer directly:
\begin{equation}\label{eq:tractor-green-transfer}
 \Lambda_{{\rm TF},\pm}=p\Lambda_{{\rm parent},\pm}i,
 \qquad
 q_{\rm TF}\Lambda_{{\rm TF},\pm}
 +\Lambda_{{\rm TF},\pm}q_{\rm TF}=1,
 \qquad
 \Lambda_{{\rm TF},+}^{\sharp}=\Lambda_{{\rm TF},-}.
\end{equation}
Finite-order pre- and postcomposition preserve the parent causal support.

The four trace/Weyl rows are attached without a nonlocal projector.  Put
$d=\frac12\operatorname{div}$ and let the local shear from split to
unshifted endpoint coordinates have
\begin{equation}\label{eq:trace-Weyl-shear}
 U_G=\begin{pmatrix}1&0\\-d&1\end{pmatrix},
 \qquad
 U_I=\begin{pmatrix}1&d^\sharp\\0&1\end{pmatrix},
 \qquad U_M=U_E=1.
\end{equation}
Its inverse is obtained by reversing the off-diagonal signs,
$U^\sharp=U^{-1}$, and $qU=U(q_{\rm TF}\oplus q_{\rm tr})$.  The trace
complex is two pointwise identity doublets with contraction $h_{\rm tr}$,
so the complete thirty-row endpoint homotopy is
\begin{equation}\label{eq:complete-endpoint-causal-homotopy}
 \Lambda_{{\rm end},\pm}
 =U(\Lambda_{{\rm TF},\pm}\oplus h_{\rm tr})U^{-1}.
\end{equation}
Combining this with the certified $356/30$ hybrid projector gives the full
all-row formula
\begin{equation}\label{eq:full-prolonged-causal-homotopy}
 \boxed{\Lambda_{{\rm prol},\pm}
 =H_{\rm alg}+i_{\rm end}\Lambda_{{\rm end},\pm}p_{\rm end}},
 \qquad
 Q_{\rm prol}\Lambda_{{\rm prol},\pm}
 +\Lambda_{{\rm prol},\pm}Q_{\rm prol}=1.
\end{equation}
Every algebraic term is finite-order local, hence
$\operatorname{supp}\Lambda_{{\rm prol},\pm}f\subseteq
J^\pm(\operatorname{supp}f)$; cyclicity gives
$\Lambda_{{\rm prol},+}^\sharp=\Lambda_{{\rm prol},-}$.  This promotes
\texttt{causal\_green\_homotopy}.  It does not construct a canonical Green
inverse for $D_{\rm TF}$, a global identity
$\Lambda_{{\rm end},\pm}=W_0G_{{\rm end},\pm}$, or a prolonged witness
$QW+WQ$ with degreewise Green inverses.  Those stronger objects remain
unproved but are unnecessary for the direct transferred homotopy
\eqref{eq:full-prolonged-causal-homotopy}.  The coefficient and assembly
receipts are respectively
\begin{center}\footnotesize
\path{adjoint_tractor_bgg_curved_pbw.json},\par
\path{curved_full_prolonged_green_homotopy_assembly.json}.
\end{center}

The selected causal route therefore has the following explicit flags:
\begin{equation}\label{eq:open-curvature-flags}
 \begin{aligned}
 &\texttt{curved\_EB\_equations},
 &&\texttt{curved\_EB\_first\_order\_closure},\\
 &\texttt{curved\_EB\_symmetric\_hyperbolicity},
 &&\texttt{curved\_sourced\_constraint\_identity},\\
 &\texttt{curved\_constraint\_propagation},
 &&\texttt{EAL\_curvature\_spectrum\_match},\\
 &\texttt{support\_local\_prolongation\_retract},
 &&\texttt{prolonged\_BV\_operator\_identity},\\
 &\texttt{prolonged\_green\_witness},
 &&\texttt{curvature\_causal\_green\_operators},\\
 &\texttt{causal\_green\_homotopy},
 &&\texttt{causal\_quasi\_isomorphism},\\
 &\texttt{residual\_endpoint\_recovery},
 &&\texttt{SO42\_equivariant\_transport},\\
 &\texttt{prolonged\_current\_comparison}.&&
 \end{aligned}
\end{equation}
Thirteen flags in \eqref{eq:open-curvature-flags} are true: the curved
equations, first-order closure, symmetric hyperbolicity, sourced identity,
constraint propagation, all-level spectrum, support-local prolongation
retract, prolonged BV operator identity, off-shell prolonged current
comparison, all-row causal Green homotopy, causal quasi-isomorphism, residual
endpoint recovery, and equivariant transport.  The two false flags are the
implementation-specific \texttt{prolonged\_green\_witness} and
\texttt{curvature\_causal\_green\_operators}.  They refer to a canonical
$D_{\rm TF}$ same-sided inverse and a degreewise $QW+WQ$ realization, neither
of which is asserted or required by the direct tractor transfer.  The paper
constructs degree-$(-1)$ maps satisfying
\begin{equation}\label{eq:causal-green-homotopy}
 Q_{\rm prol}\Lambda_\pm+\Lambda_\pm Q_{\rm prol}=1,
 \qquad
 \operatorname{supp}\Lambda_\pm f\subseteq J^\pm(\operatorname{supp}f).
\end{equation}
This is the transferred formula \eqref{eq:full-prolonged-causal-homotopy},
not a Green inverse for the canonical $D_{\rm TF}$ block.  Homogeneous
constraint propagation and symmetric hyperbolicity by themselves would not
have proved \eqref{eq:causal-green-homotopy}; the parent Green theorem and
the exact curved cyclic retract supply the missing sourced causal operation.
Nor does a nonequivariant vector-space identification suffice for the next
step.  Here the causal difference
\begin{equation}\label{eq:causal-quasi-map}
 \Lambda=\Lambda_+-\Lambda_-:
 \Gamma_c(\mathcal C_{\rm prol})[1]
 \longrightarrow\Gamma_{sc}(\mathcal C_{\rm prol})
\end{equation}
is a quasi-isomorphism.  Its cutoff inverse is represented by
$[u]\mapsto[Q(\chi_+u)]$; the two cohomology-inverse identities follow from
the past-compact and future-compact contractions and the exact support
sequence.  Since $S^3$ is compact,
$\Gamma_{sc}(\mathcal C_{\rm prol})=\Gamma^\infty(\mathcal C_{\rm prol})$.

For each of the fifteen conformal-Killing generators $\xi_a$, the compact
source $j_a=Q(\chi\xi_a)$ satisfies
$[\Lambda j_a]=[\xi_a]$.  The dual quotient sections obey the complementary
identity and the endpoint pairing is perfect.  Their compact grading is
$4_{-1}\oplus7_0\oplus4_{+1}$, the suspension scalar is $+1$, and the
curvature mapping-cylinder and auxiliary pairs are contractible; hence they
add no second ghost or momentum copy.  The causal/Cauchy bridge preserves the
residual $SO(4,2)$ action up to a certified chain homotopy.  If $\rho$ is any
local conformal chain generator and $\kappa=[Q,\chi]$, then
\begin{equation}\label{eq:causal-cutoff-equivariance}
 [\kappa,\rho]=[Q,[\chi,\rho]].
\end{equation}
The homotopy has compact support on every spacelike-compact section because
$d\chi$ lies in a compact time slab and $S^3$ is compact.  Together with the
natural curvature mapping cylinder, smooth global BGG equivariance, and the
all-level $E/A/L$ exhaustion, this induces the strict $SO(4,2)$ action on
cohomology.  The chosen Cauchy split itself need not be strictly equivariant.

The certified odd cyclic mapping cylinder gives the
quadratic parent
\begin{equation}
 S_{\rm prol}=\frac12\langle\Phi,D\Omega Q_{\rm prol}\Phi\rangle,
 \qquad I^*\omega_{\rm prol}-\omega_{\rm aux}=d\beta+Q\gamma,
 \qquad \beta=\delta Y_I,\quad\gamma=0,
\end{equation}
where $Y_I$ is the finite Lagrange concomitant of the $T$ and $A$ shears.
The current certificate now binds the corrected core-chain digest and checks
this pullback explicitly.  Hence the flag
\begin{center}\small
\path{prolonged_current_comparison}
\end{center}
is true.  Moreover
$\Lambda_+^\sharp=\Lambda_-$ makes
$\Delta_\Lambda=\Lambda_+-\Lambda_-$ graded antisymmetric.  The common
algebraic and trace contractions cancel in $\Delta_\Lambda$, and cyclic
transfer reduces its pairing to the parent Green/current identity.  The two
$d+Q$ current improvements then identify the prolonged, auxiliary, metric,
Cauchy and causal pairings on cohomology.  The normalized result is
$+I_E\oplus(-I_A)\oplus(-I_L)$ before residual reduction.

\begin{theorem}[Covariant causal transport]\label{thm:covariant-causal-H4}
For the selected closed-cylinder BV--BFV polarization, the chain
\eqref{eq:covariant-causal-bridge} consists of pairing-compatible
quasi-isomorphisms.  It recovers exactly the fifteen conformal-Killing ghost
classes and their fifteen dual endpoint classes, introduces no duplicate
residual sector, and is $SO(4,2)$-equivariant on cohomology.  Consequently
\begin{equation}\label{eq:final-covariant-H4}
 \boxed{
 H^4_{\rm cov}
 =\operatorname{span}\{[W_+^2],[W_-^2]\},
 \qquad \operatorname{sig}G_{\rm cov}=(2,0),
 \qquad G_{\rm cov}=I_2.}
\end{equation}
These are residual vertex/deformation classes.  Their Hermitian form is
positive-definite; $I_2$ is its value after the certified basis,
orientation, and unit-overlap normalization.  The theorem transports the
certified residual result; it does not recompute Chevalley--Eilenberg
cohomology in auxiliary variables or assert a positive graviton Fock norm.
\end{theorem}

\begin{proof}
The support exact sequence turns the same-sided homotopies in
\eqref{eq:causal-green-homotopy} into the map
\eqref{eq:causal-quasi-map}, which is a quasi-isomorphism.  Compactness of
$S^3$ identifies its target with
all smooth global solutions.  The cutoff identities above recover the
fifteen endpoint pairs, while the support-local SDRs contract every auxiliary
and curvature copy.  Equation~\eqref{eq:causal-cutoff-equivariance} gives
homotopy-equivariant residual transport.  Cyclicity and the two current
improvements identify the causal pairing with the normalized energy and
residual pairings.  Theorem~\ref{thm:residual-main} then gives
\eqref{eq:final-covariant-H4}.
\end{proof}

The receipts are
\begin{center}\footnotesize
\path{curved_causal_transport_recognition.json},\par
\path{curved_SO42_causal_transport_recognition.json},\par
\path{curved_direct_causal_pairing_transport.json},\par
\path{covariant_H4_transport.json},\quad
\path{covariant_gram_transport.json}.
\end{center}
The known
obstruction to factorizing the gauge-invariant spin-two conformal operator
away from Einstein backgrounds is relevant context \cite{NutmaTaronna}, but
is not used as the no-go theorem \eqref{eq:null-rank-no-go}.

\section{The on-shell Weyl module}\label{sec:weyl-module}

\subsection{Character and cylinder towers}

We now define the coefficient module used in the exact residual
calculation.  A positive-chirality Weyl curvature is a conformal primary
with labels
\begin{equation}
 (\Delta;j_L,j_R)=(2;2,0),
\end{equation}
and the opposite chirality is its parity conjugate.  The Bach equation is in
$(4;1,1)$ and its differential identity is in
$(5;\tfrac12,\tfrac12)$.  The resulting one-chirality character is
\begin{equation}\label{eq:weyl-character}
 \chi_+(q)
 =\chi_{(2;2,0)}-\chi_{(4;1,1)}
 +\chi_{(5;\frac12,\frac12)}
 =\frac{5q^2-9q^4+4q^5}{(1-q)^4}.
\end{equation}
This defines the abstract algebraic positive-energy module $\W_+$.  The
character and tower calculation below, together with the explicit
intertwiner and metric preimages of
Theorem~\ref{thm:algebraic-cylinder}, identify it with the geometric chiral
system.

On $\R\times S^3$, $\W_+$ decomposes into three multiplicity-free towers:
\begin{align}
 E_n^+&:\quad
 \left(\frac{n+2}{2},\frac{n-2}{2}\right),
 &\dim E_n&=(n+3)(n-1), &&n\ge2,\label{eq:E-tower}\\
 A_n^+&:\quad
 \left(\frac n2,\frac{n-2}{2}\right),
 &\dim A_n&=(n+1)(n-1), &&n\ge3,\label{eq:A-tower}\\
 L_n^+&:\quad
 \left(\frac n2,\frac{n-4}{2}\right),
 &\dim L_n&=(n+1)(n-3), &&n\ge4.\label{eq:L-tower}
\end{align}
The negative-chirality towers exchange $j_L$ and $j_R$.  Summing
\eqref{eq:E-tower}--\eqref{eq:L-tower} reproduces
\eqref{eq:weyl-character}.  For both chiralities the first levels have
dimensions
\begin{equation}
 \dim\cH_2=10,
 \quad \dim\cH_3=40,
 \quad \dim\cH_4=82,
 \quad \dim\cH_5=136,
 \quad \dim\cH_6=202.
\end{equation}

The labels $E,A,L$ are cylinder oscillator labels only.  In particular,
they should not be interpreted as three independently signable conformal
representations.  Proper conformal transformations mix the towers.

\subsection{Invariant form and exact generator action}

In normalized cylinder coordinates the unique standard-real,
parity-compatible invariant form is, up to overall convention,
\begin{equation}\label{eq:Jconf}
 J_{\rm conf}=+I_E\oplus(-I_A)\oplus(-I_L).
\end{equation}
It is nondegenerate and indefinite.  The proper conformal lowering and
raising generators satisfy
\begin{equation}\label{eq:J-adjoint}
 (K^-_q)^\sharp=K^+_q,
 \qquad X^\sharp=X\quad(X=D,SU(2)_L,SU(2)_R),
\end{equation}
where $X^\sharp=J_{\rm conf}^{-1}X^\dagger J_{\rm conf}$.

Multiplicity one reduces each allowed lowering block to a scalar reduced
coefficient.  With source energy $n$, the six stable families are
\begin{align}
 c_{EE}(n)&=\sqrt{\frac{2(n-1)(n+1)(n+3)}{n+2}},
 &c_{AE}(n)&=\sqrt{\frac{8(n-1)}{(n-2)(n+2)}},\nonumber\\
 c_{AA}(n)&=\sqrt{\frac{2(n-3)(n-1)(n+2)}{n-2}},
 &c_{LE}(n)&=-\sqrt{\frac{2(n-3)}{n-2}},\label{eq:reduced-coefficients}\\
 c_{LA}(n)&=\sqrt{\frac8{n-2}},
 &c_{LL}(n)&=\sqrt{2(n-2)(n+1)}.\nonumber
\end{align}
Together with Clebsch--Gordan coefficients and \eqref{eq:J-adjoint}, these
determine the all-level one-particle $\mathfrak{so}(4,2)$ action.
The exact reconstruction verifies the conformal commutators on every
interior level of a finite energy buffer; the top shell is used only as a
buffer and is never mistaken for a finite conformal representation.

The free Fock coefficient module is the bosonic second quantization
\begin{equation}\label{eq:fock-module}
 \F_{\W}=\operatorname{Sym}(\W_+\oplus\W_-),
\end{equation}
with multiplicative lift of \eqref{eq:Jconf}.  The residual generators
preserve oscillator particle number.

\section{Hamiltonian moment-map normalization}\label{sec:moment-map}

The residual action is not introduced as an arbitrary matrix
representation.  In oscillator coordinates $z$, write the symplectic form
as
\begin{equation}\label{eq:symplectic-form}
 \Omega=i\,d\bar z\,J\wedge dz.
\end{equation}
If a conformal generator $X$ acts linearly by $\delta_Xz=K_Xz$, its
quadratic Hamiltonian moment map is
\begin{equation}\label{eq:moment-map}
 \mu_X(z,\bar z)=\bar z\,M_Xz,
 \qquad M_X=JK_X,
 \qquad d\mu_X=\iota_{X_X}\Omega.
\end{equation}
The $J$-adjoint identities ensure the correct reality of $\mu_X$.

\subsection{The conformal Taub constraint}

The same residual parameters arise as reducibilities of the local
Diff$\times$Weyl complex.  Let
\begin{equation}
 \epsilon=(\xi,\sigma),
 \qquad K(\xi,\sigma)=0,
 \qquad \sigma=-\frac14\bar\nabla_\mu\xi^\mu,
\end{equation}
and let $h$ satisfy $B^{(1)}[h]=0$.  With the Taylor convention
\begin{equation}
 g(\lambda)=\bar g+\lambda h+\frac{\lambda^2}{2}k+O(\lambda^3),
\end{equation}
the second-order equation is
\begin{equation}\label{eq:second-order-Bach}
 B^{(1)}[k]+B^{(2)}[h,h]=0.
\end{equation}

\begin{proposition}[Conformal Taub solvability constraint]
\label{prop:Taub-constraint}
For every residual conformal reducibility $\epsilon$, the current
\begin{equation}
 J_\epsilon^\mu[h]
 :=\xi_\nu B^{(2)\mu\nu}[h,h]
\end{equation}
is conserved.  If $h$ is tangent to a second-order family of Bach-flat
metrics on the compact cylinder, then
\begin{equation}\label{eq:Taub-charge}
 \boxed{
 \mathcal T_\epsilon[h]
 :=\int_{S^3}n_\mu\xi_\nu
 B^{(2)\mu\nu}[h,h],d\Sigma=0.}
\end{equation}
\end{proposition}

\begin{proof}
At second order around a Bach-flat background, the exact divergence and
trace identities imply
$\bar\nabla_\mu B^{(2)\mu\nu}[h,h]=0$ and
$B^{(2)\mu}{}_{\mu}[h,h]=0$ whenever $B^{(1)}[h]=0$.  Hence
\begin{equation}
 \bar\nabla_\mu J_\epsilon^\mu
 =B^{(2)\mu\nu}\bar\nabla_{(\mu}\xi_{\nu)}
 =-\sigma B^{(2)\mu}{}_{\mu}=0.
\end{equation}
If \eqref{eq:second-order-Bach} has a solution, the charge equals minus the
integral of $\xi_\nu B^{(1)\mu\nu}[k]$.  For a reducibility parameter, the
linearized Noether current is a spatial divergence.  Since
$\partial S^3=\varnothing$, its integral vanishes.  This is the standard
compact-Cauchy linearization-stability argument
\cite{AltasTekinLinearization,AltasTekinTaub}.
\end{proof}

The covariant phase-space Noether identity identifies this same quadratic
constraint with the Hamiltonian of the residual linear action:
\begin{equation}\label{eq:moment-Taub-target}
 \mu_\epsilon(h)
 =\frac12\Omega_\Sigma(h,R_\epsilon h)
 =\mathcal T_\epsilon[h]
\end{equation}
in the Taylor and orientation conventions fixed above.  Changing the
polarization convention moves the corresponding factor of two between the
bilinear kernel and its quadratic polynomial.  Both sides form equivariant
adjoint/coadjoint $\mathfrak{so}(4,2)$ modules.  Since the complexified
adjoint module is simple, a nonzero equivariant comparison is unique up to
one scalar.  This is also the local-BRST relation between reducibilities and
conserved-current classes
\cite{BarnichAntifield,BarnichBrandtHenneaux}; the covariant Hamiltonian
identity is the higher-derivative counterpart of the standard Noether-charge
construction \cite{IyerWald}.

The equality in \eqref{eq:moment-Taub-target} has an abstract Hessian proof.
Let $H_\epsilon[g]$ be the covariant constraint Hamiltonian.  Because
$\epsilon$ fixes the background,
$H_\epsilon[\bar g]=dH_\epsilon|_{\bar g}=0$.  Its Hessian on linearized
solutions is simultaneously
\begin{equation}\label{eq:Taub-Hessian}
 d^2H_\epsilon|_{\bar g}(h,h)
 =\Omega_\Sigma(h,R_\epsilon h)
 =2\mathcal T_\epsilon[h].
\end{equation}
The first equality is the Hamiltonian definition of the residual linear
action; the second is the quadratic term in the same nonlinear constraint
with the Taylor convention of \eqref{eq:second-order-Bach}.  Thus the direct
component calculations below fix conventions and the common scalar rather
than supplying fifteen logically independent proofs.

\begin{cproposition}[All-energy Taub/moment-map normalization]
\label{prop:all-energy-Taub}
For $S_{\rm red}=-S_{\rm HH}/2$ modulo Euler, the canonical
fifteen-component quadratic kernel is
\begin{equation}
 M_X^{\rm can}=-\frac12 J K_X.
\end{equation}
The independent quadratic Bach-source calculation uses a different
normalization for the cylinder waves and conformal-Killing vectors.  In that
convention the proper-conformal kernel is
\begin{equation}\label{eq:taub-normalization}
 M_{\rm Taub}=-\frac{\sqrt2}{4\pi}\,J K^-.
\end{equation}
All six stable lowering families are therefore fixed at arbitrary source
energy by \eqref{eq:reduced-coefficients}.  The $D$ component agrees
directly with the quadratic Ostrogradsky Noether Hamiltonian.  Two
independent integrations of
$n_\mu\xi_\nu B^{(2)\mu\nu}[h_1,h_2]$ reproduce the $A_3\to E_2$ and
$L_4\to A_3$ reduced coefficients.  Exact conformal equivariance then fixes
the remaining magnetic components and compact charges.
\end{cproposition}

The proof is constructive.  The all-level matrices obey the complete
interior conformal algebra and the $J$-adjoint identities.  Multiplication
by $-J/2$ therefore gives the canonical equivariant moment map.  Direct
evaluation of the fourth-order and vector quadratic Hamiltonians gives
$+\omega$, $-\omega$, and $-\omega$ on the normalized $E,A,L$ oscillators
for $S_{\rm HH}$, and hence the displayed $D$ kernel for $S_{\rm red}$.
The two curvature-derived proper-conformal components fix the single
conversion
\begin{equation}
 M_{K^-}^{\rm raw}=\frac{\sqrt2}{2\pi}M_{K^-}^{\rm can}.
\end{equation}
The exact all-$n$ formulas and their energy-six regression buffer are
generated by
\path{symbolic/verify_conformal_taub_moment_map_all_levels.py}.

After Proposition~\ref{prop:dual-endpoint}, the quadratic Bach source also
defines a canonical Kuranishi class
\begin{equation}
 \kappa_2(h,h)=[B^{(2)}[h,h]]\in\operatorname{coker}K^\sharp.
\end{equation}
Its image under endpoint duality evaluates on a reducibility as
\begin{equation}\label{eq:endpoint-Taub-map}
 \Theta(\kappa_2(h,h))(\epsilon)
 =\langle\epsilon,B^{(2)}[h,h]\rangle
 =\mathcal T_\epsilon[h]=\mu_\epsilon(h).
\end{equation}
The exact quotient invariance, direct $D$ normalization, two independent
Bach-source seeds, and equivariant fifteen-component completion are composed
by \path{symbolic/verify_conformal_taub_obstruction_map.py}.  This does not
compute the separate time-slice transgression from the bulk endpoint degree
to BFV ghost momentum.

This section supplies the Hamiltonian normalization of the residual action
and its Taub interpretation.  It reconstructs all components by equivariance
after two direct Bach-source normalizations; it does not claim that every
magnetic block was redundantly recomputed as a separate curvature integral.
The cohomology theorem below depends only on the exact module action.

\section{Residual BRST and Cartan homotopy reduction}\label{sec:cartan}

\subsection{The absolute residual complex}

Let
\begin{equation}
 \g=\mathfrak{so}(4,2)
\end{equation}
with compact grading
\begin{equation}\label{eq:g-grading}
 \g=\g_{-1}\oplus\g_0\oplus\g_{+1},
 \qquad \dim(\g_{-1},\g_0,\g_{+1})=(4,7,4).
\end{equation}
Here $D$ and $\mathfrak{so}(4)$ lie in $\g_0$, while proper conformal
lowering and raising generators lie in $\g_{-1}$ and $\g_{+1}$.
The dual ghosts carry the opposite compact degree.

For a coefficient module $V\subset\F_{\W}$, the absolute residual complex
is
\begin{equation}\label{eq:CE-complex}
 C^q(V)=V\otimes\Lambda^q\g^*,
\end{equation}
with Chevalley--Eilenberg differential
\begin{equation}\label{eq:CE-differential}
 d=c^a\rho(G_a)
 -\frac12 f^a{}_{bc}c^bc^c\iota_a.
\end{equation}
It raises ghost number by one and preserves the total compact degree
\begin{equation}\label{eq:total-degree}
 \delta=D_{\rm matter}+D_{\rm ghost}.
\end{equation}

The standard centered polarization uses the product $v$ of the four ghosts
dual to the raising generators.  It has
\begin{equation}\label{eq:ghost-vacuum-degree}
 \operatorname{gh}(v)=4,
 \qquad D_{\rm ghost}(v)=-4.
\end{equation}
This is the residual ghost convention used in the cylinder construction of
\cite{HamadaBRST}.  Its compact degree and uniqueness as the bottom
rotational scalar follow from the $(4,7,4)$ exterior-algebra polarization,
and Lemma~\ref{lem:ghost-pairing} shows that its complementary-degree pairing
is canonical up to one overall volume.  The selected closed-universe BFV
certificate and raw cyclic form fix that volume in the algebraic cylinder
model.  Theorem~\ref{thm:pure-weyl-polarized} identifies it with the pairing
of the selected algebraic gauge-fixed BV--BFV state domain; the broader
Riegert sector of \cite{HamadaBRST} is not included.

\subsection{Full residual gauging is a physical assumption}\label{sec:D-gauge}

Hamada's background-independence programme identifies the transformations
which preserve the cylinder gauge with residual diffeomorphisms and imposes
the resulting conformal charges through BRST
\cite{HamadaCFT,HamadaBRST}.  This supplies a direct precedent for the
complex \eqref{eq:CE-complex}.  It does not prove that every quantization of
strict pure-Weyl gravity must use the same quotient.

There are at least two distinct cylinder problems.
\begin{enumerate}
\item In the \emph{fully gauged residual problem}, the cylinder is a gauge
representative of a closed generally covariant system.  All fifteen
conformal-Killing transformations, including compact time translation, are
generated by residual constraints.  The ghost for $D$ exists and the Cartan
homotopy below is an allowed operation.
\item In the \emph{fixed-background cylinder problem}, $D$ is the
Hamiltonian and $SO(4,2)$ is a global symmetry acting on states or
observables.  Its eigenspaces are not quotiented.  The $D$ ghost and its
contraction cannot be used to erase nonzero energy.
\end{enumerate}
Proposition~\ref{prop:Taub-constraint} gives the first choice a direct
linearization-stability interpretation.  Within the closed-universe Dirac
boundary problem, the slice $S^3$ has no spatial boundary supporting an
independent asymptotic charge, and the residual conformal Killing parameters
instead generate quadratic integrability constraints, including the $D$
constraint.  Subject to the all-level normalization in
\eqref{eq:moment-Taub-target}, full residual gauging is therefore a coherent
zero-boundary-charge sector of pure conformal gravity; it is not a device
introduced merely to erase positive cylinder energy.  Compactness and the
Taub identity support this interpretation but do not make it the unique
physical quantization.

Adding a boundary or asymptotic region, coupling an external clock, or
deparametrizing the cylinder can convert $D$ into a physical global charge.
That is the second problem above, not a contradiction of the first.
Because $[K^-,K^+]$ contains $D$, the second problem is not obtained by
simply deleting one ghost while leaving the rest of
\eqref{eq:CE-differential} unchanged.  It requires a different relative or
boundary-charge complex.  Theorem~\ref{thm:residual-main} answers only the
first problem.  A strict pure-Weyl derivation must state the temporal
boundary conditions, the residual gauge group, and the treatment of its
charges before calling the result a physical state space.

The precise derived-reduction statement, including the endpoint-duality and
BV--BFV suspension proof, is Theorem~\ref{thm:derived-residual-reduction}
after the local-to-residual transfer has been constructed.  Its accompanying
remark records the fixed-cylinder and relative alternatives.

\subsection{Cartan homotopy}

Let $\iota_D$ be contraction with the generator $D$ in the residual ghost
exterior algebra.  The usual Cartan identity is
\begin{equation}\label{eq:cartan-identity}
 d\iota_D+\iota_Dd=\cL_D.
\end{equation}

\begin{theorem}[Cartan homotopy reduction to compact degree zero]\label{thm:cartan}
Decompose the residual complex into total compact-degree eigenspaces,
\begin{equation}
 C=\bigoplus_{\delta}C_\delta,
 \qquad \cL_D|_{C_\delta}=\delta\,\id.
\end{equation}
Every $C_\delta$ with $\delta\ne0$ is contractible.  Hence inclusion of the
centered subcomplex is a quasi-isomorphism:
\begin{equation}\label{eq:cartan-quasiiso}
 \boxed{H^\bullet(C,d)\cong H^\bullet(C_{\delta=0},d).}
\end{equation}
\end{theorem}

\begin{proof}
On $C_\delta$ with $\delta\ne0$, define
\begin{equation}
 h_\delta=\frac{\iota_D}{\delta}.
\end{equation}
Equation \eqref{eq:cartan-identity} gives
\begin{equation}
 dh_\delta+h_\delta d=\id.
\end{equation}
Thus every closed vector of nonzero total degree has the explicit primitive
$\delta^{-1}\iota_D\Psi$.  If $\cL_D$ has a nilpotent part, the same proof
works wherever $\cL_D$ is invertible by using
$h=\iota_D\cL_D^{-1}$.
\end{proof}

The theorem depends on treating $D$ as a residual gauge generator.  If
cylinder time translation is retained as a physical Hamiltonian or a
boundary charge, contraction by its ghost is not an allowed quotient and
\eqref{eq:cartan-quasiiso} does not apply.

\subsection{Finite centered inventory}

At ghost number four, the minimal residual exterior algebra contains at
most the four negative-degree ghosts, so
\begin{equation}
 D_{\rm ghost}\ge-4.
\end{equation}
At $\delta=0$ this implies
\begin{equation}\label{eq:matter-bound}
 D_{\rm matter}\le4.
\end{equation}
Every nonvacuum Weyl quantum has weight at least two.  Therefore only
particle numbers $N=0,1,2$ can contribute to the centered ghost degree
selected by the Hamada polarization.  All $N\ge3$ sectors have matter weight at least six and are
contractible before any matrix-rank computation.

\section{Complete centered cohomology}\label{sec:centered}

\subsection{The relative weight-four kernel}

The complete matter Fock space at weight four is
\begin{equation}\label{eq:F4}
 \F_4=\cH_4\oplus\operatorname{Sym}^2\cH_2,
 \qquad \dim\F_4=82+55=137.
\end{equation}
The induced matter form has signature
\begin{equation}\label{eq:F4-signature}
 \operatorname{sig}(J_{\F_4})=(97,40).
\end{equation}
The relative conformal-primary conditions are
\begin{equation}\label{eq:relative-conditions}
 (D-4)|\Psi\rangle=0,
 \qquad R|\Psi\rangle=0,
 \qquad K^-_q|\Psi\rangle=0.
\end{equation}
Their exact kernel on the full space \eqref{eq:F4} is two-dimensional.  In a
spin-two magnetic basis its normalized chiral generators are
\begin{equation}\label{eq:W-square-states}
 |W_\pm^2\rangle
 =\sqrt{\frac25}|2,-2\rangle_\pm
 -\sqrt{\frac25}|1,-1\rangle_\pm
 +\frac1{\sqrt5}|0,0\rangle_\pm.
\end{equation}
They obey
\begin{equation}\label{eq:relative-I2}
 \langle W_r^2|J_{\F}|W_s^2\rangle=\delta_{rs}.
\end{equation}
This relative computation is not by itself an absolute cohomology theorem;
incoming exact states must also be excluded.

There is a simple representation-theoretic reason for the number two.  The
weight-two one-particle space is
\begin{equation}
 E_2^+\oplus E_2^-=(2,0)\oplus(0,2).
\end{equation}
The symmetric square decomposes as
\begin{equation}
 \operatorname{Sym}^2(2,0)
 \oplus\bigl[(2,0)\otimes(0,2)\bigr]
 \oplus\operatorname{Sym}^2(0,2).
\end{equation}
There is one compact scalar in each same-chirality symmetric square and no
scalar in the mixed product.  Since $E_2^\pm$ are lowest-weight spaces,
every $K^-_q$ annihilates both factors, so these two scalar singlets are
relative primaries.  By contrast, the one-particle weight-four space
contains
\begin{equation}
 (3,1)\oplus(2,1)\oplus(2,0)
\end{equation}
and its parity conjugate, with no compact scalar.  Thus representation
theory predicts the two relative candidates before the matrix calculation.
The absolute computation below is still necessary to exclude exactness and
to prove the one-particle acyclicity.  A complete Hochschild--Serre or
Kostant-style derivation of the latter ranks would be a useful conceptual
replacement for the large matrices \cite{CapBGG}.

\subsection{Absolute kernel and image}

The cutoff-complete absolute calculation supplies that missing test.  For
one chirality in the one-particle coefficient module, the centered complex
around ghost number four has exact dimensions
\begin{equation}\label{eq:one-particle-complex}
 C^3_0(290)\xrightarrow{d_3}
 C^4_0(1311)\xrightarrow{d_4}
 C^5_0(3657).
\end{equation}
The entries lie in
$\mathbb Q(i,\sqrt2,\sqrt3,\sqrt5)$.  Good-prime reduction modulo $241$,
combined with exact nilpotency, fixes the characteristic-zero ranks to
\begin{equation}
 \rank d_3=260,
 \qquad \rank d_4=1051.
\end{equation}
Since the ranks sum to $1311$,
\begin{equation}\label{eq:one-particle-zero}
 H^4_{\delta=0,N=1}=0.
\end{equation}
Parity supplies the identical result for the opposite chirality.

The lowest two-particle block begins at matter weight four.  Exterior
saturation makes its beginning
\begin{equation}\label{eq:two-particle-complex}
 0\longrightarrow C^4_0(55)
 \xrightarrow{d_4}C^5_0(385)
 \xrightarrow{d_5}C^6_0(1155).
\end{equation}
There is no incoming $C^3_0$ space: three residual ghosts have compact
degree at least $-3$ and cannot center a two-particle coefficient of minimum
weight four.  Exact nilpotency gives $d_5d_4=0$.
The two vectors \eqref{eq:W-square-states} lie in $\ker d_4$, while a
good-prime minor gives $\rank d_4\ge53$.  The displayed independent kernel
vectors give $\rank d_4\le53$.  Therefore
\begin{equation}\label{eq:two-particle-H}
 H^4_{\delta=0,N=2}
 =\operatorname{span}\{[W_+^2],[W_-^2]\}.
\end{equation}

There is also a deliberately independent small cross-check of this central
rank.  The standard-library-only executable
\path{symbolic/verify_conformal_residual_rank53_independent.py} reconstructs
the integral spin-two action on
$\operatorname{Sym}^2((2,0)\oplus(0,2))$ without importing the conformal
generator, Fock or BRST implementations used above.  It finds the block ranks
$14+25+14=53$ at two good primes and exhibits one exact singlet in each
chirality, so the exact characteristic-zero rank is $53$.  This check does
not replace the one-particle absolute window, full CE nilpotency, BV--BFV
transfer or covariant audits.

For the vacuum coefficient module the complex is ordinary Lie-algebra
cohomology.  Since
\begin{equation}
 H^\bullet(\mathfrak{so}(4,2);\C)
 \cong\Lambda(u_3,u_5,u_7),
\end{equation}
one has
\begin{equation}\label{eq:vacuum-zero}
 H^4_{\delta=0,N=0}=0.
\end{equation}

\subsection{The canonical residual CE pairing}

The residual ghost normalization is fixed almost completely by the compact
grading.  Write
\begin{equation}
 \g=\g_{-1}\oplus\g_0\oplus\g_{+1},
 \qquad \dim(\g_{-1},\g_0,\g_{+1})=(4,7,4),
\end{equation}
and choose nonzero determinant vectors
\begin{equation}
 v_-\in\det(\g_{+1}^*),
 \qquad
 \Theta_0\in\det(\g_0^*),
 \qquad
 v_+\in\det(\g_{-1}^*),
\end{equation}
with $v_+=v_-^\dagger$.  Their product
\begin{equation}
 \Omega_\g=v_+\wedge\Theta_0\wedge v_-
\end{equation}
is the top Berezin volume on all fifteen residual ghosts.  Define the
polarized pairing by coefficient extraction,
\begin{equation}\label{eq:canonical-ghost-pairing}
 (\alpha,\beta)_{\rm gh}
 :=[\alpha^\dagger\wedge\Theta_0\wedge\beta]_{\Omega_\g}.
\end{equation}

\begin{lemma}[Canonical centered ghost pairing]\label{lem:ghost-pairing}
After fixing the orientation and the overall residual volume,
\eqref{eq:canonical-ghost-pairing} is the unique determinant-saturating
pairing compatible with the $(4,7,4)$ polarization.  It obeys
\begin{equation}\label{eq:ghost-norm}
 (v_-,v_-)_{\rm gh}=1.
\end{equation}
Because $\mathfrak{so}(4,2)$ is semisimple and hence unimodular, the
Chevalley--Eilenberg differential is graded skew for the associated
complementary-degree Poincar\'e pairing.
\end{lemma}

\begin{proof}
Each of the three determinant lines is one-dimensional.  Once their product
is oriented and normalized, coefficient extraction gives
\eqref{eq:ghost-norm} and leaves no matrix freedom.  Unimodularity makes the
top CE volume invariant, so integration by parts in the exterior algebra
gives graded skew-adjointness of the CE differential.
\end{proof}

Consequently exact cochains are orthogonal to closed complementary-degree
cochains, and the pairing descends to residual CE cohomology.

In Hamada's oscillator notation, $v_-$ is the product of the four
proper-conformal ghosts, $\Theta_0$ contains the dilation/time ghost and the
six rotation ghosts, and the bra supplies the four conjugate ghosts
\cite{HamadaBRST}.  The factor of $i$ is a real-structure convention.  The
repository exterior-algebra certificate verifies the saturation,
Hermiticity, compact degree $-4$, and unit overlap directly.

The fixed insertion is not a nondegenerate inner product on every individual
ghost-number block.  It is the polarized form of the canonical
complementary-degree CE pairing.  It is sufficient here because the incoming
two-particle space is empty and the two normalized representatives are
fixed.  Multiplying \eqref{eq:ghost-norm} by the matter restriction
\eqref{eq:relative-I2} gives their exact residual cohomological Gram matrix.
It is Hermitian and sesquilinear on the complex oscillator module; the
parity-compatible real structure gives its real restriction.  This is an
intrinsic statement about the chosen residual CE complex.  We reserve the
phrase ``transferred pure-Weyl BV cohomological pairing'' for the
cyclic-transfer statement in Section~\ref{sec:transfer}.  The selected
algebraic BV--BFV reduction induces this residual orientation and overall
volume by Theorem~\ref{thm:pure-weyl-polarized}.

\begin{theorem}[Complete cohomology under full residual conformal gauging]\label{thm:residual-main}
Let the coefficient module be the polarized algebraic free
symmetric-algebra (Fock) realization $\F(\W_+\oplus\W_-)$ of the
linearized classical theory, let the residual algebra be
$\mathfrak{so}(4,2)$, treat $D$ as a gauge generator, and use the centered
ghost polarization \eqref{eq:ghost-vacuum-degree}.  Then the complete
minimal absolute residual cohomology at the selected ghost number four is
\begin{equation}\label{eq:main-result}
 \boxed{
 H^4_{\rm res,min}
 =\operatorname{span}\{[W_+^2],[W_-^2]\},
 \qquad \operatorname{sig}G_{\rm res}=(2,0),
 \qquad G_{\rm res}=I_2.}
\end{equation}
The two survivors are residual vertex/deformation classes.  Their Hermitian
form is positive-definite, and the displayed $I_2$ uses the fixed chiral
basis, ghost orientation, and unit-overlap normalization.  There is no
one-particle class and no higher-matter-weight class in this minimal centered
polarization.
\end{theorem}

\begin{proof}
The Cartan homotopy, Theorem~\ref{thm:cartan}, eliminates every nonzero
total compact degree.  The ghost-degree floor then reduces the centered
inventory to $N=0,1,2$.  Equations \eqref{eq:vacuum-zero},
\eqref{eq:one-particle-zero}, and \eqref{eq:two-particle-H} exhaust those
sectors.  Equations \eqref{eq:relative-I2} and \eqref{eq:ghost-norm} give the
pairing.
\end{proof}

\begin{remark}[What is positive]
The theorem does not turn the indefinite one-particle module
\eqref{eq:Jconf} into a positive particle space.  It removes the
one-particle sector from this absolute residual cohomology.  The
positive-definite $(2,0)$ form, normalized to $I_2$, belongs to two ghost-dressed
scalar vertex/deformation classes.
\end{remark}

\section{Descent and the dynamical/topological split}\label{sec:descent}

\subsection{Why weight four is selected}

Let $V_h$ be a scalar conformal primary of weight $h$, and let
\begin{equation}\label{eq:ghost-volume}
 \omega=\frac1{4!}\epsilon_{\mu\nu\rho\sigma}
 c^\mu c^\nu c^\rho c^\sigma
\end{equation}
be the top lowering-ghost product.  In the residual descent conventions,
\begin{equation}\label{eq:descent-transformations}
 sV_h=c^\mu\nabla_\mu V_h
 +\frac h4(\nabla_\mu c^\mu)V_h,
 \qquad
 s\omega=-\omega\nabla_\mu c^\mu,
 \qquad \omega c^\mu=0.
\end{equation}
It follows that both the integrated density and the unintegrated
ghost-dressed operator have the same closure defect,
proportional to $h/4-1$.  Hence
\begin{equation}\label{eq:descent-map}
 \boxed{[\omega V_4]\longleftrightarrow
 \left[\int_M d^4x\,\sqrt{-g}\,V_4\right].}
\end{equation}
The $-4$ compact degree of the ghost factor is therefore the conformal
counterpart of the four-dimensional volume form, rather than an accidental
normal-ordering shift.  This weight-four state--operator mechanism and a
Weyl-square example were already explicit in \cite{HamadaBRST}; the new use
here is the absolute exhaustion and pairing calculation of
Section~\ref{sec:centered}.

Three cohomological quotients appear and must not be identified without the
transfer theorem:
\begin{equation}\label{eq:three-cohomologies}
 H^4\bigl(\mathfrak{so}(4,2);\F_{\W}\bigr),
 \qquad
 H^{0,4}_{\rm loc}(s\mid d),
 \qquad
 \frac{\{\text{integrated local functionals}\}}
 {\{\text{allowed boundary functionals}\}}.
\end{equation}
The first is the residual Chevalley--Eilenberg group computed exactly in
this paper.  The second is local BRST cohomology modulo spacetime descent
\cite{BarnichBrandtHenneaux}.  The third depends on topology and boundary
conditions.  Equation \eqref{eq:descent-map} identifies representatives in
the residual construction; its promotion to an isomorphism with the full
local BRST group modulo all allowed spacetime boundary functionals is an
additional local-descent question, not part of the polarized state theorem.

\subsection{Parity basis}

Parity exchanges $W_+^2$ and $W_-^2$.  Define
\begin{equation}\label{eq:parity-basis}
 e=\frac{W_+^2+W_-^2}{\sqrt2},
 \qquad
 o=\frac{W_+^2-W_-^2}{\sqrt2}.
\end{equation}
Up to the orientation-dependent factor of $i$ in Lorentzian self-duality
conventions, descent identifies
\begin{equation}
 e\sim C_{\mu\nu\rho\sigma}C^{\mu\nu\rho\sigma},
 \qquad
 o\sim C_{\mu\nu\rho\sigma}\widetilde
 C^{\mu\nu\rho\sigma}.
\end{equation}
Because \eqref{eq:parity-basis} is orthogonal, its Gram matrix remains
$I_2$.

\subsection{Euler--Lagrange quotient}

Let
\begin{equation}
 \cE:H^4_{\rm res,min}
 \longrightarrow
 \{\text{local equation and gauge deformations}\}
\end{equation}
be the Euler--Lagrange map obtained by varying the integrated density.  In
the conventions of Section~\ref{sec:local-detour},
\begin{equation}\label{eq:EL-map}
 \cE(e)=\lambda_B B_{\mu\nu},
 \qquad \lambda_B\ne0,
 \qquad
 \cE(o)=0.
\end{equation}
The first identity is the Weyl-square variation.  Under locality,
smoothness, Lorentz invariance, and a four-derivative bound, it is the unique
nontrivial self-interaction of a single linear conformal graviton modulo
equivalence \cite{BoulangerHenneaux}.

The second identity is Chern--Weil transgression.  In a local frame,
\begin{equation}\label{eq:chern-weil}
 P_4=\Tr(R\wedge R)=dQ_3(\Gamma),
\end{equation}
where
\begin{equation}\label{eq:CS3}
 Q_3(\Gamma)=\Tr\left(
 \Gamma\wedge d\Gamma+\frac23\Gamma\wedge\Gamma\wedge\Gamma
 \right).
\end{equation}
Its variation is
\begin{equation}\label{eq:Pvariation}
 \delta P_4=2d\Tr(\delta\Gamma\wedge R),
\end{equation}
using the Bianchi identity.  On a finite cylinder, in a chosen global frame,
\begin{equation}\label{eq:finite-cylinder-P}
 \int_{[t_1,t_2]\times S^3}P_4
 =\CS[\Gamma(t_2)]-\CS[\Gamma(t_1)].
\end{equation}
Thus fixed-endpoint or compactly supported variations have no bulk
Euler--Lagrange derivative.

Equations \eqref{eq:EL-map}--\eqref{eq:finite-cylinder-P} give the exact
sequence
\begin{equation}\label{eq:EL-exact-sequence}
 0\longrightarrow\operatorname{span}\{o\}
 \longrightarrow H^4_{\rm res,min}
 \xrightarrow{\cE}\operatorname{span}\{B_{\mu\nu}\}
 \longrightarrow0.
\end{equation}
Quotienting variationally trivial local bulk functionals yields
\begin{equation}\label{eq:dynamic-quotient}
 H^4_{\rm dyn}
 =H^4_{\rm res,min}/\ker\cE
 \cong\operatorname{span}\{[C^2]\},
 \qquad
 \boxed{G_{\rm dyn,res}=I_1.}
\end{equation}
Equivalently,
\begin{equation}\label{eq:I2-split}
 \boxed{I_2=I_1^{\rm dynamical}\oplus I_1^{\rm topological}.}
\end{equation}

The odd class is only locally and variationally trivial.  A gravitational
theta parameter can affect boundaries, horizons, contact terms, large frame
transformations, and nontrivial topology \cite{FischlerKundu}.  We do not
quotient those global data without specifying the boundary problem.

\section{The local-to-residual transfer}\label{sec:transfer}

Theorem~\ref{thm:residual-main} begins with the locally reduced Weyl module.
This section states exactly what is needed to identify it with the complete
free pure-Weyl BV answer.

\subsection{The free Fock lift is formal}

\begin{proposition}[Algebraic symmetric-algebra lift]
\label{prop:Fock-lift}
Let $(C,q)$ be the algebraic one-particle free complex over characteristic
zero, with $q$ extended as a derivation to $\operatorname{Sym}C$.  Then
\begin{equation}\label{eq:Fock-lift}
 \boxed{
 H(\operatorname{Sym}C,q)
 \cong\operatorname{Sym}H(C,q).}
\end{equation}
In particular, once the one-particle coefficient module is
$\W_+\oplus\W_-$, its free many-particle cohomology is
$\F(\W_+\oplus\W_-)$.
\end{proposition}

\begin{proof}
Degree by degree, $C^{\otimes N}$ obeys the K\"unneth theorem.  Taking
symmetric coinvariants under the finite group $S_N$ is exact in
characteristic zero, so
$H(\operatorname{Sym}^NC,q)\cong\operatorname{Sym}^NH(C,q)$.  Taking the
algebraic direct sum over $N$ proves \eqref{eq:Fock-lift}.
\end{proof}

Thus the Fock-space coefficient module is not an independent field-theory
input and is not the unpolarized classical BV complex.
Proposition~\ref{prop:Fock-lift} is the symmetric-algebra lift of the free
linearized BRST differential.  Its selected
Hilbert/Krein symmetric-Fock realization and the closed residual operator
are supplied separately by Theorem~\ref{thm:energy-mode-completion}; that
theorem does not claim an arbitrary topological tensor-product completion.

\subsection{Equivariant Cartan transfer}

Let $(C,q)$ be the local coefficient complex after the fifteen
conformal-Killing zero modes have been separated, and let $H=H(q)$.  In the
field-theoretic application to states, $C$ denotes the time-slice-reduced,
polarized complex of \eqref{eq:BV-BFV-polarization}, not the entire
unpolarized bulk BV tangent complex.  Suppose
there is a strong deformation retract
\begin{equation}\label{eq:SDR}
 (H,0)\mathrel{\mathop{\rightleftarrows}^{\jmath}_{p}}(C,q),
 \qquad
 \jmath p=\id-qs-sq,
 \qquad p\jmath=\id_H.
\end{equation}
Let $\Delta$ contain the residual zero-mode differential and mixed terms, so
$Q=q+\Delta$.  With the sign convention in \eqref{eq:SDR}, homological
perturbation gives
\begin{align}\label{eq:HPL-data}
 I&=(\id+s\Delta)^{-1}\jmath,
 &P&=p(\id+\Delta s)^{-1},
 &Q_H&=p\Delta(\id+s\Delta)^{-1}\jmath.
\end{align}
If
\begin{equation}\label{eq:D-equivariance}
 [D,q]=[D,\Delta]=[D,s]=0,
 \qquad pD=D_Hp,
 \qquad D\jmath=\jmath D_H,
\end{equation}
then the transferred operators
\begin{equation}
 \iota_D^H=P\iota_DI,
 \qquad \cL_D^H=P\cL_DI
\end{equation}
obey
\begin{equation}\label{eq:transferred-Cartan}
 [Q_H,\iota_D^H]_+=\cL_D^H.
\end{equation}
Thus every transferred subspace on which $\cL_D^H$ is invertible is again
contractible.

\subsection{Cyclic transfer and the pairing}

Let $\sharp$ be the graded adjoint for the full BV/Krein pairing.  Suppose
the unperturbed retract is cyclic,
\begin{equation}\label{eq:cyclic-SDR}
 \jmath^\sharp=p,
 \qquad (s\Delta)^\sharp=\Delta s.
\end{equation}
Then taking the adjoint of \eqref{eq:HPL-data} gives
\begin{equation}
 I^\sharp=P.
\end{equation}
Since $PI=\id_H$,
\begin{equation}\label{eq:HPL-isometry}
 \boxed{I^\sharp I=\id_H.}
\end{equation}
The dressed inclusion is therefore an exact isometry: contractible local
dressing changes representatives without changing the reduced Gram matrix.

\begin{proposition}[Blockwise cyclic retract]\label{prop:cyclic-retract}
Let a finite compact-energy and $SO(4)$ block carry a nondegenerate graded
pairing, let $q^\sharp=-q$, and suppose the induced pairing on $H(q)$ is
nondegenerate.  Then the block admits a cyclic strong deformation retract.
The algebraic direct sum of these retracts may be chosen compact-degree and
$SO(4)$ equivariant.
\end{proposition}

\begin{proof}
Write $B=\im q$ and choose a nondegenerate representative space
$\cH\subset\ker q$.  In $\cH^\perp$, choose an isotropic complement $L$ to
$B$.  Then $q:L\to B$ is an isomorphism.  Define $s|_B=(q|_L)^{-1}$ and set
$s=0$ on $\cH\oplus L$.  The decomposition
$C=\cH\oplus B\oplus L$ gives $qs+sq=1-\jmath p$, while the hyperbolic
pairing of $B$ and $L$ gives the cyclic adjoint relation.  Choosing the
splitting in each finite symmetry block and taking the algebraic direct sum
preserves compact degree and $SO(4)$ type.
\end{proof}

This proposition does not provide one splitting which is simultaneously
equivariant under the noncompact generators $K^\pm$.  Such a splitting
cannot be obtained by averaging over the nonunitary conformal module, which
need not be semisimple.  Compatibility with the full $SO(4,2)$ action, the
global zero-mode split, and the BV pairing remains a field-theory
obligation.  Boundedness on a completed cylinder space is a separate
analytic question.

\subsection{Raw polynomial instantiation and measured noncompact defects}

\paragraph{Computational-supplement material.}
This subsection records the coordinate instantiation, measured cutoff
defects, and basis-level chain dictionaries behind the transfer theorem.  It
is retained in the archival monolith, but its record-by-record implementation
content is assigned to the computational supplement in the planned
major-revision split; the extracted residual paper will retain only the
resulting propositions and their precise claim boundaries.

The abstract discussion can be instantiated before choosing an analytic
completion.  In the finite polynomial conformal category, fixed total
compact energy gives the exact rational chain
\begin{equation}\label{eq:raw-polynomial-chain}
 G_n\xrightarrow{A_n}M_n\xrightarrow{B_n}E_n
 \xrightarrow{C_n}I_n,
\end{equation}
where the four slots contain the Diff$\times$Weyl ghosts, trace-free metric
and trace doublet, Bach equation/metric-antifield row, and Noether
identity/ghost-antifield row.  The Weyl scalar is written as
\begin{equation}
 \omega=\sigma+\frac14\partial\!\cdot\!\xi,
\end{equation}
so the vector arrow is the trace-free conformal-Killing operator and the
trace is an explicit doublet.  The raw matrices commute with all four
coordinate translations and special-conformal vector fields, not only with
$D\times SO(4)$.

The interpretation of these rows can be derived directly from the minimal
master action rather than assigned after the calculation.  With conventional
BV variables $(h_{\mu\nu},c^\mu,\omega)$ and antifields
$(h^{*\mu\nu},c^*_\mu,\omega^*)$, take
\begin{align}\label{eq:minimal-master-action}
 S_{\rm BV}={}&S_W[g]
 +\int g^{*\mu\nu}\bigl(\cL_cg_{\mu\nu}+2\omega g_{\mu\nu}\bigr)
 +\int c^*_\mu\frac12[c,c]^\mu
 +\int\omega^*\cL_c\omega .
\end{align}
At the conformally flat background, the tangent differential of its
quadratic part has the four-row form
\begin{equation}\label{eq:field-BV-tangent-chain}
 (c,\omega)\xrightarrow{K}(h)
 \xrightarrow{B_{\rm lin}}(h^*)
 \xrightarrow{K^\sharp}(c^*,\omega^*).
\end{equation}
Here tangent-complex degree is minus conventional BV ghost number: the four
rows have degrees $-1,0,1,2$, whereas the underlying BV coordinates have
ghost numbers $1,0,-1,-2$.  In a flat conformal chart, one consistent
antibracket convention gives
\begin{align}
 Qh_{\mu\nu}&=2\partial_{(\mu}c_{\nu)}+2\omega\eta_{\mu\nu},
 &Qh^*_{\mu\nu}&=B_{\rm lin}[h]_{\mu\nu},\nonumber\\
 Qc^*_\nu&=-2\partial^\mu h^*_{\mu\nu},
 &Q\omega^*&=2h^{*\mu}{}_{\mu}.
 \label{eq:quadratic-field-BV-differential}
\end{align}
The nonlinear terms in \eqref{eq:minimal-master-action}, including
$c^*[c,c]$, $\omega^*\cL_c\omega$, and the field-dependent part of
$h^*\cL_ch$, do not belong to this free tangent differential.

Define the invertible component redefinitions
\begin{align}\label{eq:field-raw-redefinitions}
 \Omega&=\omega+\frac14\partial\!\cdot c,
 &\tau&=\frac18\operatorname{tr}h,
 &h_0&=h-\frac14\eta\operatorname{tr}h,\nonumber\\
 \tau^*&=2\operatorname{tr}h^*,
 &h_0^*&=h^*-\frac14\eta\operatorname{tr}h^*,
 &i^*_\mu&=-\frac12\left(c^*_\mu+\frac14\partial_\mu\omega^*\right),
 \qquad \Omega^*=\omega^*.
\end{align}
Then \eqref{eq:field-BV-tangent-chain} becomes exactly
\begin{equation}
 (c,\Omega)\longmapsto(K_0c,\Omega),\qquad
 (h_0,\tau)\longmapsto(B_{\rm lin}h_0,0),\qquad
 (h_0^*,\tau^*)\longmapsto(\partial\!\cdot h_0^*,\tau^*),
\end{equation}
which is the raw chain \eqref{eq:raw-polynomial-chain}.  In particular,
$\Omega\to\tau$ and $\tau^*\to\Omega^*$ are displayed unit contractible
pairs.

\begin{cproposition}[Minimal field-theoretic realization of the raw detour complex]
\label{prop:field-BV-dictionary}
For each $D$-finite, $SO(4)$-finite polynomial block in the complete centered
buffer $n=2,3,4,5$, the redefinitions
\eqref{eq:field-raw-redefinitions} give exact rational inverse maps
\begin{equation}
 F_n:C^{\rm min}_{\rm BV,n}\longrightarrow C^{\rm raw}_n,
 \qquad
 G_n:C^{\rm raw}_n\longrightarrow C^{\rm min}_{\rm BV,n}
\end{equation}
satisfying
\begin{equation}\label{eq:minimal-field-raw-chain-maps}
 F_nQ_{\rm BV}=Q_{\rm raw}F_n,qquad
 G_nQ_{\rm raw}=Q_{\rm BV}G_n,qquad
 F_nG_n=G_nF_n=1.
\end{equation}
Thus the comparison homotopies vanish: this phase is a chain isomorphism,
not merely a rank-preserving quasi-isomorphism.  If $P_{\rm tr}$ selects
$(\Omega,\tau,\tau^*,\Omega^*)$ and $s_{\rm tr}$ reverses the two unit
arrows, the same matrices verify
\begin{equation}\label{eq:trace-contraction}
 [Q_{\rm raw},P_{\rm tr}]=0,
 \qquad
 Q_{\rm raw}s_{\rm tr}+s_{\rm tr}Q_{\rm raw}=P_{\rm tr}.
\end{equation}
Consequently, in the selected algebraic blocks,
\begin{equation}\label{eq:minimal-field-direct-sum}
 \boxed{
 C^{\rm min,alg}_{\rm BV,n}
 \cong C^{\rm tf,min}_{\rm raw,n}\oplus C_{{\rm tr},n},
 \qquad H(C_{{\rm tr},n})=0.}
\end{equation}
The block dimensions are
\begin{equation}
\begin{array}{c|rrrr|r}
n&\dim G_n&\dim M_n&\dim E_n&\dim I_n&\dim C_n\\ \hline
2&90&100&0&0&190\\
3&160&200&0&0&360\\
4&259&350&10&1&620\\
5&392&560&40&8&1000.
\end{array}
\end{equation}
The same triangular trace change is invertible on the inhomogeneous
low-mode gauge space and maps the full fifteen-dimensional conformal-Killing
kernel to $\Omega=0$ without changing its rank.  More precisely, the exact
low-mode matrices give
\begin{equation}\label{eq:minimal-zero-mode-identity}
 T\bigl(\ker K_{\rm BV}\bigr)=\ker K_{\rm raw}
 \cong\mathfrak{so}(4,2).
\end{equation}
Hence no residual conformal zero mode is created, removed, or absorbed into
the trace contraction.
\end{cproposition}

The executable
\begin{center}
\path{symbolic/verify_conformal_field_bv_dictionary.py}
\end{center}
constructs both sides independently, verifies every matrix identity in
\eqref{eq:minimal-field-raw-chain-maps}, and emits a 2170-record table with
one entry per raw basis vector, including conventional ghost number,
antifield number, compact degree, $SO(4)$ content, differential image, and
an exact equality flag.  This proposition identifies the \emph{minimal
quadratic chain}.  By itself it does not derive a chosen gauge-fixed nonminimal
domain, replace both sides of the global reducibility sector, inventory all
additional BV rows, or transfer the field-theoretic cyclic pairing.

The first of these qualifications can be closed without choosing an
analytic completion, but it is an implementation corollary rather than a
premise of the gauge-invariant classical theorem.  Adjoin vector and scalar antighost--multiplier pairs
and their antifields,
\begin{equation}
 (\bar c_\mu,b_\mu;\bar c^{*\mu},b^{*\mu}),\qquad
 (\bar\omega,b_\omega;\bar\omega^*,b_\omega^*),
\end{equation}
with the standard coordinate-BRST rules $s\bar c=b$, $sb=0$ and their Weyl
counterparts.  The suspended tangent complex used in
\eqref{eq:field-BV-tangent-chain} has the transposed arrows
$b\to\bar c$ and $\bar c^*\to b^*$.  Both conventions are retained in the
machine dictionary.

In trace-split variables choose the conformal Landau fermion
\begin{equation}\label{eq:cylinder-gauge-fermion}
 \Psi_{\rm gf}=\int_M\!\sqrt{|\bar g|}\,
 \left(\bar c_\perp^\mu\bar\nabla^\nu h^0_{\mu\nu}
       +\bar\omega\tau\right),
 \qquad \bar c_\perp=P_\perp\bar c .
\end{equation}
It has ghost number $-1$, is a weight-four scalar density, and commutes with
$D\times SO(4)$.  No multiplier-square term is used: all multipliers remain
in the displayed complex and are removed only by an explicit homotopy.  The
canonical transformation is
\begin{equation}
 T_\Psi:\quad
 \phi_A^*\longmapsto
 \phi_A^*+\frac{\delta\Psi_{\rm gf}}{\delta\phi^A},
 \qquad
 Q_{\rm gf}=T_\Psi Q_{\rm ext}T_\Psi^{-1}.
\end{equation}
For the vector condition the two Euler derivatives are
$\delta\Psi/\delta\bar c=\bar\nabla\!\cdot h_0$ and
$\delta\Psi/\delta h_0=-K_0\bar c/2$ in the conventions of
\eqref{eq:quadratic-field-BV-differential}; the trace derivatives are the
two unit maps $\tau\leftrightarrow\bar\omega$.

\begin{cproposition}[Gauge-fixed/nonminimal equivalence]
\label{prop:gauge-fixed-equivalence}
Let $C^{\rm ext}_n=C^{\rm min}_{\rm BV,n}\oplus C_{{\rm nm},n}$, where
$C_{{\rm nm},n}$ contains the complete Diff and Weyl nonminimal pairs and
their antifield duals.  In every centered block $n=2,3,4,5$, the Hessian of
\eqref{eq:cylinder-gauge-fermion} gives an exact unipotent map
$T_{\Psi,n}=1+N_n$, with $N_n^2=0$, such that
\begin{equation}
 Q_{{\rm gf},n}=T_{\Psi,n}Q_{{\rm ext},n}T_{\Psi,n}^{-1},
 \qquad Q_{{\rm gf},n}^2=0.
\end{equation}
If $p_0,\jmath_0,s_0$ are the direct-sum contraction of the nonminimal
pairs, then
\begin{equation}
 p_{\rm gf}=p_0T_\Psi^{-1},\qquad
 \jmath_{\rm gf}=T_\Psi\jmath_0,\qquad
 s_{\rm gf}=T_\Psi s_0T_\Psi^{-1}
\end{equation}
obey
\begin{equation}\label{eq:gauge-fixed-contraction}
 p_{\rm gf}\jmath_{\rm gf}=1,
 \qquad
 \jmath_{\rm gf}p_{\rm gf}
 =1-Q_{\rm gf}s_{\rm gf}-s_{\rm gf}Q_{\rm gf}.
\end{equation}
Consequently,
\begin{equation}\label{eq:gauge-fixed-equivalence}
 \boxed{
 C^{\rm gf,alg}_{\rm BV,n}
 \simeq C^{\rm min,alg}_{\rm BV,n}\oplus C_{{\rm nm,triv},n},
 \qquad H(C_{{\rm nm,triv},n})=0.}
\end{equation}
The certified dimensions are
\begin{equation}
\begin{array}{c|rrr}
n&\dim C^{\rm min}_n&\dim C^{\rm nm}_n&\dim C^{\rm gf}_n\\ \hline
2&190&52&242\\
3&360&128&488\\
4&620&264&884\\
5&1000&480&1480.
\end{array}
\end{equation}
On the inhomogeneous low-mode block, the exact projector $P_Z$ and its
fifty-dimensional complement satisfy
\begin{equation}\label{eq:gauge-fixed-zero-modes}
 P_ZQ_{\rm gf}=Q_{\rm gf}P_Z,
 \qquad Q_{\rm gf}P_Z=0,
 \qquad \rank P_Z=15,
 \qquad \ker(K|_{\operatorname{im}P_\perp})=0.
\end{equation}
Thus gauge fixing neither absorbs nor duplicates any conformal-Killing
mode.  No generalized auxiliary field is silently eliminated.
\end{cproposition}

The executable
\begin{center}
\path{symbolic/verify_conformal_gauge_fixed_equivalence.py}
\end{center}
constructs the canonical shear and the transported contraction independently
in every block.  Its 3094-record ledger classifies every minimal field, nonminimal
field, multiplier, and antifield coordinate.  Proposition
\ref{prop:gauge-fixed-equivalence} does not compute the degree-shifting
time-slice transgression \eqref{eq:zero-mode-transgression}, establish the
post-polarization row ledger, or transfer the field-theoretic pairing;
those constructions are supplied below.  The bulk dual endpoint itself is
identified independently in Proposition~\ref{prop:dual-endpoint}.

\subsection{The dual endpoint and the BFV degree shift}

There is a short exact argument for the dual bulk zero modes.  Write the
minimal detour chain in the form
\begin{equation}
 G\xrightarrow{A}M\xrightarrow{B}E\xrightarrow{C}I,
 \qquad C=\pm A^\sharp,
\end{equation}
where the tangent BV structure gives perfect pairings between $G$ and $I$
and between $M$ and $E$.  Let $Z=\ker A$.

\begin{proposition}[Canonical dual endpoint]\label{prop:dual-endpoint}
The endpoint quotient pairs perfectly with the residual gauge kernel:
\begin{equation}\label{eq:dual-endpoint}
 \boxed{
 \Theta:I/\operatorname{im}C\xrightarrow{\ \sim\ }Z^*,
 \qquad \Theta([u])(z)=\langle z,u\rangle_{G,I}.}
\end{equation}
In the exact conformal-Killing block,
\begin{equation}
 \dim Z=\dim\operatorname{coker}C=15,
 \qquad
 Z=4_{-1}\oplus7_0\oplus4_{+1},
 \qquad
 Z^*=4_{+1}\oplus7_0\oplus4_{-1}.
\end{equation}
\end{proposition}

\begin{proof}
For $z\in\ker A$,
$\langle z,Ce\rangle=\pm\langle Az,e\rangle=0$, so $\Theta$ is
well defined.  If $\Theta([u])=0$, then
$u\in(\ker A)^\perp$.  Finite-dimensional adjoint linear algebra gives
$(\ker A)^\perp=\operatorname{im}A^\sharp=\operatorname{im}C$, proving
injectivity; equality of dimensions gives surjectivity.
\end{proof}

The machine certificate does more than compare dimensions.  If $Z$ is the
$65\times15$ matrix of the already certified conformal-Killing basis and
$J_{GI}$ is the endpoint pairing, it constructs
\begin{equation}
 \Theta=Z^\dagger J_{GI},
 \qquad
 J_{GI}C=A^\dagger J_{ME},
\end{equation}
verifies $\Theta C=0$, and produces an exact section $S$ with
$\Theta S=I_{15}$.  It also verifies
\begin{equation}
 I=\operatorname{im}C\oplus S(Z^*)
\end{equation}
by an exact rank-$65$ decomposition.  These identities and the dual compact
grading are generated by
\path{symbolic/verify_conformal_dual_zero_modes.py}.

The interpretation requires one further distinction.  The bulk endpoint
class, the residual BFV momentum, and the constraint value all transform in
dual modules, but they are not the same coordinate:
\begin{equation}\label{eq:residual-role-table}
 \begin{array}{c|c|c|c}
 \text{object}&\text{space}&\text{degree}&\text{role}\\ \hline
 c^a&Z[1]&+1&\text{residual ghost}\\
 b_a&Z^*[-1]&-1&\text{BFV ghost momentum}\\
 \mu_a&Z^*\text{-valued functions}&0&\text{constraint}\\
 \lbrack u\rbrack&\operatorname{coker}C\cong Z^*&\text{bulk endpoint}&\text{obstruction codomain}
 \end{array}
\end{equation}
The time-slice BV--BFV construction must supply a degree-shifting
transgression
\begin{equation}\label{eq:zero-mode-transgression}
 \tau:H^{\rm bulk}_{\rm endpoint}\longrightarrow Z^*[-1]_{\rm BFV},
\end{equation}
in the sense of the bulk-to-boundary BV--BFV correspondence
\cite{CattaneoMnevReshetikhin}.
Equivariance reduces it to one scalar,
$\tau=\lambda\Theta$, after the generator bases are fixed.  In the selected
algebraic BFV convention this scalar is no longer free.  Normalize endpoint
representatives by $\Theta([u_a])=e^a$.  The endpoint/Taub theorem and the
canonical BFV charge give, on the ghost-free slice,
\begin{equation}\label{eq:lambda-sign-test}
 Q_{\rm bulk}u_D=\mu_D,
 \qquad
 Q_{\rm BFV}b_D=\mu_D,
 \qquad b_D=\lambda u_D,
\end{equation}
so compatibility forces $\lambda=+1$ for
$\Omega_{\rm gh}=\delta b_a\wedge\delta c^a$ and the sign of
$\Omega_{\rm res}$ displayed below.  Equivariance extends the result to all
fifteen components.  Reversing either canonical sign reverses $\lambda$ but
cannot leave an arbitrary magnitude.  The exact role audit verifies one
residual ghost copy, one endpoint quotient, one $15+15$ BFV cotangent
sector, and refuses to identify $[u]$ with $b_a$ without the degree
suspension \eqref{eq:zero-mode-transgression}; see
\path{symbolic/verify_conformal_residual_bfv_roles.py}.
The normalized suspension, homogeneous generator order, cotangent
orientation, and centered overlap are certified independently by
\path{symbolic/verify_conformal_zero_mode_transgression.py}.

Schematically, the no-duplication map is a transfer, not a deletion:
\begin{equation}\label{eq:zero-mode-replacement}
 \begin{array}{ccc}
 Z_{\rm local}\oplus H^{\rm bulk}_{\rm endpoint}
 &\subset&C^{\rm min}_{\rm BV}\\
 \downarrow\scriptstyle{\mathrm{id}\oplus\tau}&&\downarrow\\
 Z[1]\oplus Z^*[-1]&\subset&C^{\rm res}_{\rm BFV}.
 \end{array}
\end{equation}
The moment map is a function into $Z^*$ and is not a third set of zero-mode
coordinates.

\subsection{The free symmetric-algebra state complex is obtained after polarization}

The full cyclic bulk BV tangent complex should not be asserted to have only
one cohomological row: its ghost, physical, and antifield sectors occur with
their cyclic duals.  The single-row criterion used by the residual
calculation belongs after the sequence
\begin{equation}\label{eq:BV-BFV-polarization}
 C^{\rm bulk}_{\rm BV}
 \longrightarrow C^\Sigma_{\rm BFV}
 \longrightarrow C^{+}_{\rm state},
\end{equation}
where the first arrow is time-slice BV--BFV transgression and the second is
canonical phase-space reduction followed by the positive-frequency
Lagrangian polarization.  Nonminimal quartets contract in this passage;
negative-frequency modes supply symplectic duals rather than additional
positive-energy coefficient states.  The final coefficient complex is the
polarized free symmetric-algebra (Fock) realization
\begin{equation}
 \F(\W_+\oplus\W_-)\otimes\Lambda^\bullet\g^*,
\end{equation}
not the unpolarized bulk BV antifield complex.  Here ``Fock'' names
the bosonic symmetric-algebra construction associated with the chosen
positive-frequency Lagrangian subspace; it does not assert a positive
one-particle Hilbert metric.

Correspondingly, the even form on the positive-energy module is induced
from the real phase-space symplectic form,
\begin{equation}\label{eq:pairing-pipeline}
 \Omega_{\rm BV}\longrightarrow\Omega_{\rm BFV}
 \longrightarrow\Omega_\Sigma
 \longrightarrow
 J_+(u,v)=-i\Omega_\Sigma(\bar u,v),
\end{equation}
in the convention \eqref{eq:symplectic-form}; reversing the convention
reverses both displayed signs.  This is distinct from
restricting the odd BV form to one positive-frequency half.  A complete row
ledger must therefore be generated after the two arrows in
\eqref{eq:BV-BFV-polarization}.  The selected algebraic ledger and pairing
are constructed below.

\begin{cproposition}[Raw metric-BV retract and conformal transfer]
\label{prop:raw-BV-transfer}
In the complete energy buffer $n=2,3,4,5$, exact rational changes of basis
give maps $p_n,\jmath_n,s_n$ satisfying
\begin{equation}
 p_n\jmath_n=1,
 \qquad
 \jmath_np_n=1-q_ns_n-s_nq_n,
\end{equation}
with cohomology dimensions $10,40,82,136$.  The chosen retract is not
strictly equivariant under the noncompact generators.  For one Cartesian
component, the defect ranks are
\begin{equation}\label{eq:raw-defect-ranks}
\begin{array}{c|rrrr}
 &\rank\rho_H&\rank\delta_{\jmath}&\rank\delta_p&\rank\delta_s\\
\hline
 P_0&40&29&5&35\\
 K_0&82&80&9&97.
\end{array}
\end{equation}
Every inclusion defect is explicitly $q$-exact.  Each projection defect has
an explicit right $q$-homotopy, while the homotopy defect obeys the induced
projector identity.  The transferred maps obey all sixteen $[K_b,P_a]$
brackets exactly on cohomology.  Moreover,
\begin{equation}\label{eq:physical-HPL-zero}
 p\rho_a s\rho_b\jmath=0
\end{equation}
on the physical metric row for every pair of proper-conformal components;
all higher terms vanish because a second $s$ acts on the gauge row.
\end{cproposition}

Thus full equivariance of an arbitrarily chosen representative inclusion is
not needed and, for this natural exact splitting, is false.  What transfers
strictly is the conformal action on cohomology.  The calculation therefore
realizes the homotopy-equivariant outcome anticipated above rather than
silently replacing it by a strict SDR.  The exact matrices, defect
homotopies, and generated table are produced by
\path{symbolic/verify_conformal_raw_bv_transfer.py}.

\begin{cproposition}[Raw cross-energy cyclic form]
\label{prop:raw-cross-energy}
Let $J_2$ be the positive energy-two form pulled back from the electric and
magnetic Weyl curvature.  In the raw polynomial cohomology basis the four
raising blocks have joint full row rank, and the recursion
\begin{equation}\label{eq:cross-energy-recursion}
 J_nK^+_{a,n-1}=-(K^-_{a,n})^\dagger J_{n-1},
 \qquad a=1,\ldots,4,
\end{equation}
determines a unique nondegenerate $J_n$ at every subsequent energy.  Through
the complete centered buffer its exact inertias are
\begin{equation}
\begin{array}{c|rrrr}
n&2&3&4&5\\ \hline
(n_+,n_-)&(10,0)&(24,16)&(42,40)&(64,72).
\end{array}
\end{equation}
These are precisely the parity-complete $+E,-A,-L$ signatures.

Transporting the form through the exact raw adapted basis makes
$q,s,\jmath,p$ cyclic in the ghost, metric, equation-antifield and
identity-antifield rows.  Every elementary representative dressing
$s\rho_A\jmath$ lies in their common isotropic gauge subspace.  Consequently
all mutual quadratic dressing terms vanish and
\begin{equation}
 I^\sharp I=1
\end{equation}
for the transferred inclusion in the centered window.
\end{cproposition}

The form is obtained from \eqref{eq:cross-energy-recursion}, not from a
postulated $E/A/L$ change of basis.  Exact rational congruence reduction
computes the displayed inertias without numerical eigenvalues.  The
certificate is
\path{symbolic/verify_conformal_cross_energy_pairing.py}.

\begin{cproposition}[Complete centered HPL correction test]
\label{prop:full-HPL-test}
For every allowed ordered pair of the fifteen residual generators on source
energies two, three, and four which remains in the coefficient buffer,
\begin{equation}
 p\Delta s\Delta\jmath=0,
 \qquad
 s\Delta s\Delta\jmath=0.
\end{equation}
The second identity holds before projection and kills every higher physical
row term.  The universal ghost-structure term acts trivially on the local
rows and $s\jmath=0$.  The only path beyond the buffer is double raising from
energy four; it vanishes by the four-raising-ghost exterior saturation of
the centered top shell.  Hence
\begin{equation}
 \boxed{Q_H=p\Delta\jmath=d_{\rm CE}}
\end{equation}
in the complete centered coefficient window.
\end{cproposition}

This check contains $555$ compact, lowering, raising, and mixed ordered
pairs; it is stronger than the proper-conformal pair test in
\eqref{eq:physical-HPL-zero}.  Its machine report is generated by
\path{symbolic/verify_conformal_full_hpl_transfer.py}.

\subsection{The filtered local-to-residual complex}

The possible failure modes can be displayed without compressing them into
the phrase ``no other row.''  This filtration is formed \emph{after} the
time-slice transfer, canonical phase-space reduction, positive-frequency
polarization, and nonminimal-quartet contraction in
\eqref{eq:BV-BFV-polarization}; it is not the cohomology spectral sequence
of the full unpolarized cyclic bulk tangent complex.  After the
conformal-Killing zero modes have been transferred from the local ghosts,
introduce the residual-ghost filtration on the polarized state-resolution
bicomplex
\begin{equation}\label{eq:bicomplex}
 C^{p,q}=\Lambda^p\g^*\widehat\otimes C^q_{\rm local},
 \qquad
 F^pC=\bigoplus_{r\ge p}C^{r,*}.
\end{equation}
Here $p$ is residual ghost number and $q$ is the degree in the remaining
polarized local resolution.  Antifields and negative-frequency conjugates
which supplied the bulk cyclic duals are not silently deleted: their role
must already have been accounted for by the two arrows in
\eqref{eq:BV-BFV-polarization}.  A transferred degree-one differential has
the filtered expansion
\begin{equation}\label{eq:transferred-expansion}
 Q_H=q_{\rm local}+d_{\rm res}
 +\sum_{k\ge2}Q_{[k]},
\end{equation}
with bidegrees
\begin{equation}
 |q_{\rm local}|=(0,1),
 \qquad |d_{\rm res}|=(1,0),
 \qquad |Q_{[k]}|=(k,1-k).
\end{equation}
The terms $Q_{[k]}$ encode higher transferred ghost operations; assuming
that they vanish is a substantive strictness assertion, not automatic
homological bookkeeping.

The resulting spectral sequence begins
\begin{align}\label{eq:spectral-pages}
 E_1^{p,q}
 &=\Lambda^p\g^*\otimes H^q(q_{\rm local}),
 &E_2^{p,q}
 &=H^p\!\left(\g;H^q(q_{\rm local})\right),
\end{align}
and
\begin{equation}\label{eq:spectral-bidegree}
 d_r:E_r^{p,q}\longrightarrow E_r^{p+r,q-r+1}.
\end{equation}
For the desired term $E_r^{4,0}$, a higher differential can enter from
$(4-r,r-1)$ for $r=2,3,4$, and can leave toward $(4+r,1-r)$ for
$2\le r\le11$.  Thus local classes in degrees $q=1,2,3$ and antifield or
negative-degree classes down to $q=-10$ are potential modifiers.  They must
be inventoried or excluded; it is not enough to compute only the local
degree-zero row.  Since $0\le p\le15$, this filtration is bounded.  On each
finite compact-energy block the cochain spaces used here are finite, so the
spectral sequence converges.  An infinite completed state space additionally
requires a complete, separated filtration.  Even after collapse, a cyclic
splitting is needed to control extension data and the Hermitian form.

This is the standard local-BRST setting in which antifield and gauge-fixed
cohomologies must be compared carefully
\cite{BarnichBrandtHenneaux,BarnichHenneauxHurthSkenderis}.

\begin{proposition}[Single-row strictness and collapse]
\label{prop:single-row}
For the polarized state-resolution bicomplex just described, suppose the
global zero modes have been transferred exactly,
$H^q(q_{\rm local})=0$ for $q\ne0$, and the free residual perturbation
$\Delta$ has bidegree $(1,0)$.  Suppose in addition that the chosen SDR is
compatible with the local grading, so that $s$ has bidegree $(0,-1)$ and
$\jmath,p$ are supported in local degree zero.  Then
\begin{equation}\label{eq:strict-transfer}
 \boxed{Q_H=p\Delta\jmath,}
\end{equation}
all higher HPL terms vanish, and the local-to-residual spectral sequence
collapses to
\begin{equation}
 H^p\!\left(\g;H^0(q_{\rm local})\right).
\end{equation}
\end{proposition}

\begin{proof}
The homotopy $s$ lowers local degree by one.  Since both $\jmath(H)$ and the
image retained by $p$ lie in local degree zero,
\begin{equation}
 p\Delta(s\Delta)^r\jmath=0\qquad(r\ge1).
\end{equation}
The HPL series therefore reduces to \eqref{eq:strict-transfer}.  The
$E_1$ page has only the row $q=0$, so no higher differential can enter or
leave it.
\end{proof}

The proposition concerns the perturbation obtained from the \emph{free
metric BV bicomplex}.  It must not be inferred by naively restricting a
nonlinear minimal $Q$-manifold.  In the Dneprov--Grigoriev variables,
curvature-dependent terms schematically of the form $\xi\xi W$ and
$\xi\xi C$ occur in ghost transformations and can have total compact degree
zero \cite{DneprovGrigoriev}.  They belong to the nonlinear comparison and
must either be identified with the interacting deformation or removed by
the equivalence back to the free metric presentation before
Proposition~\ref{prop:single-row} is applied.

\begin{proposition}[Formal filtered-transfer criterion]\label{prop:formal-transfer}
Suppose the zero-mode split produces \eqref{eq:bicomplex}; the filtration is
convergent; $H^q(q_{\rm local})$ vanishes in every row that can enter or
leave $(4,0)$; and the transferred differential on $H^0(q_{\rm local})$ is
the strict residual Chevalley--Eilenberg differential, with all relevant
$Q_{[k]}$ absent.  Then the full cohomology in the selected filtered degree
is isomorphic, as a vector space, to the residual group
\begin{equation}
 H^4\!\left(\g;H^0(q_{\rm local})\right).
\end{equation}
If in addition the retract is compact-degree equivariant and cyclic, the
Cartan homotopy and the cohomological Hermitian pairing transfer, and the
dressed inclusion is the isometry \eqref{eq:HPL-isometry}.
\end{proposition}

\begin{proof}
Under the stated row and strictness assumptions no differential
\eqref{eq:spectral-bidegree} enters or leaves $(4,0)$, so the bounded
spectral sequence collapses there to its $E_2$ value.  The equivariant and
cyclic statements are \eqref{eq:transferred-Cartan} and
\eqref{eq:HPL-isometry}.  The cyclic splitting fixes the extension of the
pairing from the associated graded object.
\end{proof}

\subsection{End-to-end centered integration in the polynomial category}

The transferred coefficient matrices of
Proposition~\ref{prop:raw-BV-transfer} can be coupled directly to the
independently generated residual BFV/CE ghosts.  This removes the
hand-specified $E/A/L$ matrices from the final centered rank calculation.

\begin{cproposition}[Metric-to-residual integration]
\label{prop:metric-residual-integration}
In the parity-complete algebraic polynomial complex, the pipeline
\begin{equation}
 \text{metric BV rows}\longrightarrow
 H(q)\longrightarrow
 C^\bullet(\mathfrak{so}(4,2);H(q))
\end{equation}
has the exact centered dimensions and ranks
\begin{equation}\label{eq:integrated-rank-table}
\begin{array}{c|rrr|rr|r}
 &\dim C^3&\dim C^4&\dim C^5&\rank d_3&\rank d_4&\dim H^4\\
\hline
\text{vacuum}&147&407&833&116&291&0\\
\text{one particle}&580&2622&7314&520&2102&0\\
\text{two particles}&0&55&385&0&53&2.
\end{array}
\end{equation}
For the two-particle row, the next space has dimension $1155$ and
$d_5d_4=0$ exactly.  Orientation reversal acts on the two kernel vectors by
\begin{equation}
 \operatorname{diag}(-1,+1).
\end{equation}
Pulling the positive energy-two electric/magnetic curvature form back to
the metric representatives gives the raw kernel Gram matrix
\begin{equation}
 \operatorname{diag}\!\left(\frac5{64},5\right).
\end{equation}
After exact normalization and multiplication by the intrinsic unit residual
ghost norm, the representative Gram matrix is $I_2$.
\end{cproposition}

The odd and even eigenvectors are respectively the Pontryagin and
Weyl-square directions.  The calculation verifies the vector-space transfer
and a positive-definite representative form on this two-dimensional
cohomology, normalized to $I_2$, in the algebraic model.
Proposition~\ref{prop:raw-cross-energy} independently fixes its
cross-energy cyclic normalization and dressed isometry.  The certificate and
machine-readable ranks are produced by
\begin{center}
\path{symbolic/verify_conformal_metric_to_residual_integration.py}.
\end{center}

\subsection{The selected closed-universe BFV problem}

The residual quotient is now specified as part of the theorem's input.  The
Cauchy surface is the closed oriented $S^3$, no boundary-charge variables or
external clock are adjoined, and all fifteen conformal reducibilities are
represented by BFV constraints.  In particular,
\begin{equation}\label{eq:closed-BFV-choice}
 \partial S^3=\varnothing,\qquad
 \rank Q_{\partial}=0,\qquad
 D\in\{G_A\}_{A=1}^{15}.
\end{equation}
The complementary-degree CE pairing then gives
Lemma~\ref{lem:ghost-pairing}.  This makes the centered polarization and the
use of $\iota_D$ internally consistent with the chosen boundary problem.
These are precisely the closed-universe assumptions entering the derived
zero-locus and stabilizer quotient of
Theorem~\ref{thm:derived-residual-reduction}.

Equation~\eqref{eq:closed-BFV-choice} is a selection, not a universal
theorem.  Adding a boundary, clock, asymptotic charge, or deparametrized
Hamiltonian removes $D$ from the constraint list and disables the Cartan
contraction.  The two alternatives are represented explicitly by
\path{symbolic/verify_conformal_closed_universe_bfv.py}.

\subsection{The selected algebraic BV--BFV completion}

The smooth BGG theorem and Theorem~\ref{thm:algebraic-cylinder} remove the
global PDE-surjectivity and finite-mode preimage gaps.
Propositions~\ref{prop:raw-BV-transfer} and
\ref{prop:metric-residual-integration}, together with
Propositions~\ref{prop:raw-cross-energy} and \ref{prop:full-HPL-test},
additionally close the kernel, image, noncompact action, cyclic form,
transferred differential, and centered-rank calculation in an explicit
polynomial metric-BV model.  Proposition~\ref{prop:field-BV-dictionary}
also identifies its four minimal rows exactly with the tangent complex
derived from the quadratic master action.
Proposition~\ref{prop:gauge-fixed-equivalence} further derives a concrete
Landau-gauge canonical transformation, retains both nonminimal antifield
duals, contracts the entire nonminimal extension, and proves that the
fifteen conformal-Killing modes are unchanged.  Proposition
\ref{prop:dual-endpoint} now also identifies the bulk endpoint quotient
canonically with $Z^*$, while \eqref{eq:endpoint-Taub-map} identifies the
quadratic endpoint obstruction with the certified Taub/moment map.  In the
selected algebraic BV--BFV polarization these ingredients now close.

\begin{theorem}[Selected algebraic BV--BFV polarization]
\label{thm:pure-weyl-polarized}
Work on the $D$-finite, $SO(4)$-finite cylinder.  Choose the closed-universe
BFV problem \eqref{eq:closed-BFV-choice} and the positive-frequency
Lagrangian polarization.  Then:
\begin{enumerate}
\item the normalized endpoint suspension is
\begin{equation}\label{eq:tau-equals-theta}
 \boxed{\tau=\Theta,\qquad\lambda_D=\lambda_{\rm all}=+1,}
\end{equation}
and transfers, rather than copies, the local zero-mode cotangent pair into
one residual pair $(c^a,b_a)$;
\item the polarized free symmetric-algebra state complex is exactly
\begin{equation}\label{eq:polarized-state-complex}
 \boxed{
 \cH_{\rm state}
 =\operatorname{Sym}(\W_+\oplus\W_-)
 \otimes\Lambda^\bullet\mathfrak{so}(4,2)^*,}
\end{equation}
with local nonzero-mode cohomology concentrated in polarized degree zero;
\item the induced Hermitian field pairing on the two centered residual
classes is positive-definite and, in the displayed normalized basis, is
\begin{equation}\label{eq:field-I2}
 \boxed{
 \left\langle W_r^2v_-,W_s^2v_-\right\rangle_{\rm field}
 =\delta_{rs},\qquad G_{\rm field}=I_2.}
\end{equation}
\end{enumerate}
No single-row or Fock-space assertion is made about the unpolarized classical
bulk BV tangent complex.
\end{theorem}

\begin{proof}
Choose endpoint representatives $u_a$ by
$\Theta([u_a])=e^a$.  Equation~\eqref{eq:endpoint-Taub-map} gives
$Q_{\rm bulk}u_a=\mu_a$ in the quotient.  For
$\Omega_{\rm gh}=\delta b_a\wedge\delta c^a$ and
\begin{equation}\label{eq:residual-BFV-charge}
 \Omega_{\rm res}=c^a\mu_a-\frac12f^a{}_{bc}c^bc^cb_a,
\end{equation}
the ghost-free Hamiltonian vector field gives $Q_{\rm BFV}b_a=\mu_a$.
The $D$ component of $b_a=\lambda u_a$ therefore forces $\lambda=1$;
coadjoint equivariance gives all fifteen components.  In the ordered
$4_{-1}\oplus7_0\oplus4_{+1}$ basis, the induced determinant orientation
agrees with $v_+\wedge\Theta_0\wedge v_-$ and fixes
$(v_-,v_-)_{\rm gh}=1$.

The complexified physical phase space splits as
$L_+\oplus L_-$, where $L_+=\W_+\oplus\W_-$ and $L_-$ is its symplectic
conjugate.  The exact cross-energy form makes both halves Lagrangian and
induces $J_+(u,v)=-i\Omega_\Sigma(\bar u,v)$.  The all-energy conformal
action preserves $L_+$.  Every local, trace, and nonminimal complement is an
explicit differential doublet.  On its polynomial state algebra the standard
number-operator contraction leaves only the vacuum.  Finally $c^a$ acts by
exterior multiplication and $b_a$ by $\iota_a$, so no second ghost exterior
algebra or endpoint state row is present.  This proves
\eqref{eq:polarized-state-complex}.

The action-normalized energy-two $E$ primary has unit symplectic norm.  An
exact energy-two isometry identifies this form with the raw curvature
pullback $J_2$; the contravariant recursion
\eqref{eq:cross-energy-recursion} fixes every higher energy.  Multiplying the
resulting normalized two-class matter Gram matrix by the unit oriented ghost
overlap yields \eqref{eq:field-I2}.  The exact constructions and sector
ledger are generated by the three executables
\begin{center}\small
\path{symbolic/verify_conformal_zero_mode_transgression.py},\\
\path{symbolic/verify_conformal_polarized_state_complex.py},\\
\path{symbolic/verify_conformal_polarized_pairing_transfer.py}.
\end{center}
\end{proof}

\begin{theorem}[Derived residual reduction for the selected closed universe]
\label{thm:derived-residual-reduction}
Let $\mathcal P_{\rm lin}$ denote the formal linearized pure-Weyl phase
space after quotienting local Diff$\times$Weyl transformations with
vanishing Cauchy data.  Let
\begin{equation}
 \g=\operatorname{Lie}\operatorname{Stab}_{\rm gauge}(\bar g)
   \cong\mathfrak{so}(4,2)
\end{equation}
be the infinitesimal residual stabilizer of the cylinder representative,
and let $\mu:\mathcal P_{\rm lin}\to\g^*$ be the quadratic moment map
normalized by \eqref{eq:moment-Taub-target}.  Assume that
\begin{enumerate}
\item $\Sigma=S^3$ is closed and carries no boundary or asymptotic charge;
\item no external clock is adjoined and the cylinder is not deparametrized;
\item all fifteen residual reducibilities, including $D$, are gauged;
\item one works formally near $\bar g$ in the zero-charge, second-order
integrable sector; and
\item the residual action is treated infinitesimally, without assuming a
free group action or an integrated $SO(4,2)$ representation.
\end{enumerate}
Then, in the selected formal algebraic category constructed in this paper,
the residual classical problem is represented by the derived Hamiltonian
reduction
\begin{equation}\label{eq:derived-residual-reduction}
 \boxed{
 \mathcal P_{\rm der}
 =\left[\mathcal P_{\rm lin}
   \mathop{\times}^{\mathbf R}_{\g^*}\{0\}\,/\!/^{\mathbf R}\g\right].}
\end{equation}
Equivalently, it is represented by the residual BFV complex with coordinates
$c^a$ and $b_a$ of degrees $+1$ and $-1$, whose differential satisfies
\begin{equation}\label{eq:derived-residual-BFV}
 Qf=c^a\{\mu_a,f\},\qquad
 Qc^a=-\frac12f^a{}_{bc}c^bc^c,\qquad
 Qb_a\big|_{c=0}=\mu_a.
\end{equation}
After the BV--BFV time-slice map and the selected positive-frequency
polarization, this becomes the absolute complex
\eqref{eq:CE-differential} on the free symmetric-algebra coefficient module.
\end{theorem}

\begin{proof}
Proposition~\ref{prop:Taub-constraint} proves that every second-order
integrable perturbation lies in the zero locus of the fifteen quadratic Taub
charges.  Equation~\eqref{eq:moment-Taub-target}, with the all-energy
normalization of Proposition~\ref{prop:all-energy-Taub}, identifies those
charges componentwise with the equivariant Hamiltonian moment map $\mu$.
Thus the equation imposed by each residual BFV momentum is the actual
linearization-stability constraint, rather than a new algebraic restriction.

The residual kernel is $Z=\ker K=\ker A\cong\g$ in the notation of the
minimal detour chain.  Proposition~\ref{prop:dual-endpoint}
and \eqref{eq:dual-endpoint} give the perfect pairing
$\operatorname{coker}K^\sharp\cong Z^*\cong\g^*$; the role separation in
\eqref{eq:residual-role-table} distinguishes this endpoint quotient, the
constraint value $\mu$, and the BFV momentum.  The time-slice map
\eqref{eq:zero-mode-transgression} then supplies the required degree
suspension.  The normalized identity \eqref{eq:tau-equals-theta}, proved in
Theorem~\ref{thm:pure-weyl-polarized}, shows that it produces exactly one
cotangent pair $(c^a,b_a)$ with degrees $(+1,-1)$ and no duplicate endpoint
sector.

With the canonical ghost symplectic form, the residual BFV charge
\eqref{eq:residual-BFV-charge} generates \eqref{eq:derived-residual-BFV}.
The restriction $c=0$ in its last equation isolates the Koszul term; the
full $Qb_a$ also contains the coadjoint ghost term.  Equivariance of $\mu$
and the Jacobi identity give $Q^2=0$.  Consequently the Koszul part is the
formal derived zero fibre of $\mu$, while the ghost part is the
Chevalley--Eilenberg derived quotient by the infinitesimal stabilizer.  This
is precisely \eqref{eq:derived-residual-reduction}; nonregular stabilizer
directions are retained rather than replaced by an ordinary manifold
quotient \cite{CattaneoMnevReshetikhin}.

Finally, item~2 of Theorem~\ref{thm:pure-weyl-polarized} identifies the
polarized coefficient algebra with
\eqref{eq:polarized-state-complex}.  There $c^a$ acts by exterior
multiplication and $b_a$ by $\iota_a$, so the BFV vector field generated by
\eqref{eq:residual-BFV-charge} is exactly the absolute differential
\eqref{eq:CE-differential}.  This proves the last assertion.  The argument
is formal and infinitesimal in the selected algebraic category; it neither
constructs a global singular derived stack nor identifies the unpolarized
classical BV tangent complex with a graviton Fock space.
\end{proof}

\begin{remark}[Fixed cylinder and relative alternatives]
In the fixed-background cylinder problem the moment-map values are global
charges rather than constraints: one neither restricts to $\mu^{-1}(0)$ nor
quotients by the full stabilizer, and the coefficient module remains an
$SO(4,2)$ module with $D$ as Hamiltonian.  The relative primary kernel
\eqref{eq:relative-conditions} is useful representation-theoretic data, but
is not the derived reduction \eqref{eq:derived-residual-reduction}.  Gauging
only a subgroup, retaining boundary charges, or deparametrizing time requires
the corresponding relative or boundary BFV complex; it is not obtained by
deleting the $D$ ghost from the absolute complex.
\end{remark}

\begin{corollary}[Algebraic strict-pure-Weyl result]
\label{cor:algebraic-pure-weyl}
For the selected algebraic closed-cylinder BV--BFV data, the complete
free/linearized polarized residual state cohomology in the selected degree is
\eqref{eq:main-result}.  Its field-induced Hermitian form is positive-definite
and normalized to $I_2$, and its
local Euler--Lagrange quotient is \eqref{eq:dynamic-quotient} with form
$I_1$.
\end{corollary}

\begin{proof}
Theorems~\ref{thm:curvature-isomorphism} and
\ref{thm:algebraic-cylinder} identify the one-particle module, while
Proposition~\ref{prop:Fock-lift} gives its free Fock lift.  The minimal
field/raw dictionary and gauge-fixed contraction exhaust the local algebraic
variables.  Theorem~\ref{thm:pure-weyl-polarized} supplies the one-copy
residual BFV sector and its pairing.  Propositions
\ref{prop:raw-BV-transfer}, \ref{prop:full-HPL-test}, and
\ref{prop:metric-residual-integration} then give the strict residual
differential, the centered kernel and image, and the two normalized classes.
\end{proof}

This theorem is deliberately category-limited.  It constructs the selected
finite algebraic degree suspension characterized by endpoint duality and
compatibility of the two cohomological vector fields; it is not an
independent continuity theorem for a temporal boundary one-form.  It also
does not prove closed range or uniqueness among other boundary
polarizations.  A Hilbert, Krein, distributional, or completed Fock
statement requires separate analysis.

\section{Scope, interpretation, and open questions}\label{sec:scope}

\paragraph{Residual gauge choice.}
The result is specific to quotienting the full residual $SO(4,2)$ action on
the compact cylinder.  If some generators are retained as global or
asymptotic charges, the physical complex is different.  In particular,
Cartan contraction cannot be used to erase nonzero $D$ energy when $D$ is
the physical Hamiltonian.

\paragraph{No positive graviton Fock claim.}
The unreduced module remains indefinite and the one-particle residual
cohomology is zero.  The positive-definite pairing belongs to two
weight-four residual vertex/deformation classes.
Calling $I_2$ a positive graviton Hilbert space would therefore misidentify
both the cohomological degree and the state--operator map.

\paragraph{Classical BV data and the free state representation.}
The local BV--BFV differential, current, and residual moment map are
classical linearized data.  Choosing a positive-frequency Lagrangian
polarization and completing the symmetric algebra produces the selected free
BRST state representation; it is already a free quantization choice, not an
unpolarized classical phase space.  The result does not show that interactions
preserve the reduced space or pairing, and it does not establish the quantum
master equation.  A local Diff$\times$Weyl anomaly can obstruct quantum
nilpotency even though Corollary~\ref{cor:algebraic-pure-weyl} and
Theorem~\ref{thm:energy-mode-completion} hold in the stated free/linearized
category.

\paragraph{Separate quantum local-algebra status.}
A separate \texttt{LOCAL-ALGEBRAIC} programme classifies curvature
counterterm and anomaly candidates.  It currently supplies invariant-basis
reductions at derivative orders four and six, Weyl/Schouten and Hodge-sign
checks, the universal diffeomorphism descent, the complete intrinsic Euler
tower, and fail-closed interfaces for importing antifields and forming
relative cohomology.  Its antifield-zero lower-form grading enumeration now
has 720 refined signatures and a complete factored tensor-slot realization;
the Bianchi, integration-by-parts, dimension-specific, and sparse-matrix
quotients remain open.  These results do not enter the Lorentzian causal
theorem proved here.  In particular, they do not yet establish the production
$H^{0,4}(s\mid d)$ or $H^{1,4}(s\mid d)$ quotients, anomaly coefficients,
antifield completion, restoration of the quantum master equation, or a
Lorentzian quantum theory.  Detailed ranks, lifecycle states, and certificate
paths are maintained in the companion computational supplement,
\path{quantum-weyl/README.md}, and \path{quantum-weyl/reports/}, rather than
repeated in this paper.
\paragraph{Vertex versus dynamics.}
The two residual classes do not define two independent local gravitational
dynamics.  The parity-even class is the Weyl action direction; the
parity-odd class is a theta direction, locally related by a canonical
boundary transformation.  The positive dynamical quotient is therefore
one-dimensional.

\paragraph{Bridge, completion, and covariant boundary.}
The all-level smooth metric-to-$\mathcal W_{\rm geom}$ identification and
the algebraic $E/A/L$ curvature intertwiner are exact.  The raw polynomial
metric-BV complex now supplies its explicit ghost and antifield detour rows,
explicit
noncompact homotopies, strict induced conformal action, cross-energy cyclic
form, vanishing centered HPL corrections, canonical residual CE pairing,
and end-to-end centered ranks.  Its minimal four-row differential has also
been derived from the quadratic field-theoretic master action and identified
by exact inverse chain maps.  A concrete conformal-Landau gauge fermion and
the full vector and scalar nonminimal sectors have now also been canonically
adjoined and explicitly contracted, with all fifteen conformal-Killing modes
preserved.  The selected algebraic endpoint-to-BFV suspension now has
$\lambda=+1$; the post-BFV positive-frequency ledger contracts every
nonphysical doublet to its vacuum and leaves
\eqref{eq:polarized-state-complex}; and the pairing transfer through
$\Omega_{\rm BV}\to\Omega_{\rm BFV}\to\Omega_\Sigma\to J_+$ gives
$G_{\rm field}=I_2$.
The direct Taub comparison does not require first identifying a
degree-three BGG representative, although that geometric comparison remains
valuable.  The normalized energy modes now have the Krein--Fock completion
of Theorem~\ref{thm:energy-mode-completion}, with closed residual operator,
closed range, and unchanged $H^4$.  The reduced Lorentzian metric fields
have the Cauchy--Sobolev realization of
Theorem~\ref{thm:cauchy-sobolev}, while
Proposition~\ref{thm:auxiliary-green-witness} supplies the exact curved
operator identities, support-local BV deformation retract, and covariant
current comparison.  The null-symbol theorem
\eqref{eq:null-rank-no-go} rules out a scalar-wave witness on the original
$24$-field bundle and identifies the uncancelled quotient with the two
physical helicities; the reduced Weyl symbol \eqref{eq:reduced-Weyl-symbol}
is an isomorphism.  The exact $26$-state Weyl--Cotton equations, their
constraint-adjusted symmetric-hyperbolic closure, sourced subsidiary
identity and all-level $E/A/L$ spectrum are now proved.  The all-row local
prolongation, BV differential, and off-shell current comparison are exact.
The transferred all-row causal homotopy
\eqref{eq:full-prolonged-causal-homotopy} is exact.  Its causal difference is
a compact-to-global quasi-isomorphism on the cylinder; cutoff representatives
recover all fifteen residual endpoints and their duals without duplication,
and the cutoff commutator transfers the $SO(4,2)$ action up to chain homotopy.
The direct cyclic Green/current comparison completes the pairing transport.
Thus Theorem~\ref{thm:covariant-causal-H4} proves the final covariant $H^4$
statement by transport through \eqref{eq:covariant-causal-bridge}, not by
another residual CE calculation.  The two false flags in
\eqref{eq:open-curvature-flags} remain scoped to the canonical endpoint Green
inverse and prolonged-witness implementations and are not dependencies of
that theorem.  Hadamard completion and uniqueness among boundary or
deparametrized polarizations remain separate analytic questions.  The exact ranks in
Table~\ref{tab:detour-ranks} remain independent convention and regression
checks.

\section{Conclusion}\label{sec:conclusion}

Residual conformal reduction changes the state-space question rather than
re-signing the indefinite oscillators of free Weyl gravity.  For the selected
closed-cylinder polarization, Cartan contraction reduces the absolute
residual complex to three matter sectors.  The vacuum and one-particle
sectors are acyclic, while the two-particle sector gives
\begin{equation*}
 H^4_{\rm res,min}
 =\operatorname{span}\{[W_+^2],[W_-^2]\},
 \qquad \operatorname{sig}G_{\rm res}=(2,0),
 \qquad G_{\rm res}=I_2.
\end{equation*}
They are Weyl-square and Pontryagin deformation directions, not one-particle
gravitons.  The residual Hermitian form is positive-definite and is
normalized to $I_2$ only after the selected basis, orientation, and
unit-overlap choices; the variational quotient leaves one dynamical class.

The covariant theorem follows by transport, not by a second residual
calculation.  The auxiliary and curvature mapping-cylinder retracts are
support local, the Weyl--Cotton system carries the complete $E/A/L$ module,
and the curved adjoint-tractor plus hybrid contraction gives causal
homotopies on all prolonged rows.  Their causal difference is a
compact-to-global quasi-isomorphism, recovers all fifteen residual endpoints
without duplication, transfers $SO(4,2)$ up to the cutoff chain homotopy, and
identifies the Green and current pairings.  Consequently
Theorem~\ref{thm:covariant-causal-H4} gives
$H^4_{\rm cov}=\operatorname{span}\{[W_+^2],[W_-^2]\}$ with signature
$(2,0)$, normalized to $G_{\rm cov}=I_2$ after the same choices.

The theorem remains classical, free, and polarization dependent.  The
rank-$11$ versus rank-$9$ null-symbol result rules out only the stated
$24$-field scalar-wave witness, and the invariant two-wave factorization is
excluded only in its certified family.  A canonical endpoint Green inverse,
a prolonged degreewise Green witness, alternative boundary polarizations,
distributional/Hadamard completion, interaction descent, quantum-master
restoration, and quantum unitarity are not claimed.

\appendix

\section{Conventions}\label{app:conventions}

The cylinder radius is one and the Lorentzian signature is $(-,+,+,+)$.
Compact energy is the eigenvalue of $D=i\partial_t$ in the oscillator
convention.  The compact algebra is
$SO(4)\cong SU(2)_L\times SU(2)_R$.  A label $(j_L,j_R)$ has dimension
$(2j_L+1)(2j_R+1)$.  Positive chirality means the lowest Weyl primary
$(2;2,0)$; parity exchanges $j_L$ and $j_R$.

The normalized module convention is \eqref{eq:Jconf}.  The reduced curvature
action \eqref{eq:Sred} differs by an overall factor and sign from the common
$-\int C^2$ oscillator convention.  All moment-map comparisons derive the
symplectic form and Taub kernel from the same action, so the overall action
factor cancels.  Lorentzian self-duality may insert a factor of $i$ in the
odd parity combination; no conclusion depends on that convention.

The residual generator degrees are $(0^7,-1^4,+1^4)$ and dual ghost degrees
are their negatives.  Ghost number and compact degree are independent.
The selected absolute residual degree in this paper is ghost number four,
fixed by \eqref{eq:ghost-vacuum-degree}; it is not ordinary ghost number
zero.  We call it a candidate physical degree only under the full residual
gauging specified in Section~\ref{sec:D-gauge}.

\section{Exact rank certification}\label{app:ranks}

This implementation-facing appendix is assigned to the computational
supplement in the planned paper split.  It remains here until the extracted
articles and supplement build and verify independently.

The one-particle matrices contain radicals.  They are represented over
\begin{equation}
 K=\mathbb Q(i,\sqrt2,\sqrt3,\sqrt5).
\end{equation}
For a lower bound on rank, the calculation maps this field to
$\mathbb F_{241}$ by
\begin{equation}
 i\mapsto64,
 \qquad \sqrt2\mapsto22,
 \qquad \sqrt3\mapsto56,
 \qquad \sqrt5\mapsto103.
\end{equation}
Each image obeys its defining square relation.  A nonzero minor after this
good-prime reduction is a nonzero characteristic-zero minor.  Exact
nilpotency provides the complementary upper bound
\begin{equation}
 \rank d_{q-1}+\rank d_q\le\dim C^q.
\end{equation}
For the one-particle complex the modular lower bounds saturate this inequality.
For the two-particle complex the two explicit independent kernel vectors
provide the upper bound which the modular rank saturates.  Thus the ranks
quoted in Section~\ref{sec:centered} are exact characteristic-zero ranks,
not probabilistic modular estimates.

\section{Endpoint factorization no-go receipt}
\label{app:endpoint-factorization-receipt}

This generated algebraic receipt is supplement material.  It remains an
active, tracked input of the archival monolith during the split.

For auditability, this appendix records the nonzero part of the exact
Nullstellensatz certificate used in
\eqref{eq:endpoint-factor-nullstellensatz}.  The variables
$a_0,\ldots,a_8$ are the nine coefficients of the invariant first-order
term $A_-$ after the cubic equation has imposed $A_+=-A_-$.  The $f_i$ are
the indexed residual quadratic coefficient equations after the three
zeroth-order sums have been eliminated; all unlisted multipliers vanish.
The table is generated mechanically from
\path{curved_endpoint_metric_factorization_no_go.json}, and its
generator re-expands the displayed identity over $\mathbb Q$ before writing
the TeX file.

% Generated by symbolic/generate_endpoint_factorization_nullstellensatz_tex.py.
% Source certificate SHA-256: 1e435ada0c08822f5e3bbbd7c8ed43ad17c4252ab0aaed9d3d76cfddd399231d
% Do not edit this receipt by hand.
\begingroup
\scriptsize
\noindent\(\begin{aligned} f_{0}&=-2 a_{0} a_{4}\mathbin{-}a_{0} a_{5} \\ &\quad \mathbin{-}16 a_{1} a_{4}\mathbin{-}8 a_{1} a_{5} \\ &\quad \mathbin{-}8 a_{2} a_{4}\mathbin{-}4 a_{2} a_{5} \\ &\quad \mathbin{+}2\end{aligned}\)\par
\noindent\(\begin{aligned} w_{0}&=\frac{3}{128}\Bigl(18 a_{0}^{2}\mathbin{-}192 a_{0} a_{1} \\ &\quad \mathbin{-}72 a_{0} a_{2}\mathbin{+}23 a_{0} a_{6} \\ &\quad \mathbin{+}384 a_{1}^{2}\mathbin{+}192 a_{1} a_{2} \\ &\quad \mathbin{+}62 a_{1} a_{6}\mathbin{+}58 a_{2} a_{6} \\ &\quad \mathbin{+}a_{3} a_{7}\mathbin{-}152 a_{4} a_{6} \\ &\quad \mathbin{+}224 a_{4} a_{8}\mathbin{-}68 a_{5} a_{6} \\ &\quad \mathbin{+}112 a_{5} a_{8}\Bigr)\end{aligned}\)\par\medskip
\noindent\(\begin{aligned} f_{1}&=2 a_{0} a_{6}\mathbin{-}3 a_{0} a_{8} \\ &\quad \mathbin{+}16 a_{1} a_{6}\mathbin{-}24 a_{1} a_{8} \\ &\quad \mathbin{+}8 a_{2} a_{6}\mathbin{-}12 a_{2} a_{8} \\ &\quad \mathbin{-}2\end{aligned}\)\par
\noindent\(\begin{aligned} w_{1}&=\frac{3}{128}\Bigl(27 a_{0} a_{4}\mathbin{+}24 a_{0} a_{5} \\ &\quad \mathbin{-}48 a_{1} a_{4}\mathbin{-}60 a_{1} a_{5} \\ &\quad \mathbin{+}2 a_{3}^{2}\mathbin{+}2 a_{3} a_{7} \\ &\quad \mathbin{-}224 a_{4}^{2}\mathbin{-}224 a_{4} a_{5} \\ &\quad \mathbin{-}56 a_{5}^{2}\mathbin{+}6 a_{7}^{2} \\ &\quad \mathbin{-}58\Bigr)\end{aligned}\)\par\medskip
\noindent\(\begin{aligned} f_{2}&=-a_{0} a_{8}\mathbin{+}6 a_{1} a_{6} \\ &\quad \mathbin{-}8 a_{1} a_{8}\mathbin{+}2 a_{2} a_{6} \\ &\quad \mathbin{-}4 a_{2} a_{8}\mathbin{-}2\end{aligned}\)\par
\noindent\(\begin{aligned} w_{2}&=- \frac{3}{128}\Bigl(81 a_{0} a_{4}\mathbin{+}72 a_{0} a_{5} \\ &\quad \mathbin{-}144 a_{1} a_{4}\mathbin{-}180 a_{1} a_{5} \\ &\quad \mathbin{+}6 a_{3}^{2}\mathbin{+}6 a_{3} a_{7} \\ &\quad \mathbin{-}224 a_{4}^{2}\mathbin{-}224 a_{4} a_{5} \\ &\quad \mathbin{-}56 a_{5}^{2}\mathbin{+}18 a_{7}^{2} \\ &\quad \mathbin{-}174\Bigr)\end{aligned}\)\par\medskip
\noindent\(\displaystyle f_{4}=- 2 a_{0} a_{3}\)\par
\noindent\(\begin{aligned} w_{4}&=\frac{3}{128}\Bigl(2 a_{3} a_{6}\mathbin{-}a_{4} a_{7} \\ &\quad \mathbin{+}a_{5} a_{7}\mathbin{-}4 a_{6} a_{7}\Bigr)\end{aligned}\)\par\medskip
\noindent\(\begin{aligned} f_{7}&=- \frac{1}{3}\Bigl(3 a_{0} a_{5}\mathbin{+}18 a_{1} a_{4} \\ &\quad \mathbin{+}6 a_{1} a_{5}\mathbin{+}6 a_{2} a_{4} \\ &\quad \mathbin{+}6 a_{2} a_{5}\mathbin{-}10\Bigr)\end{aligned}\)\par
\noindent\(\begin{aligned} w_{7}&=\frac{3}{128}\Bigl(21 a_{0} a_{6}\mathbin{-}72 a_{1} a_{6} \\ &\quad \mathbin{-}3 a_{3} a_{7}\mathbin{-}64 a_{4} a_{6} \\ &\quad \mathbin{-}32 a_{5} a_{6}\Bigr)\end{aligned}\)\par\medskip
\noindent\(\displaystyle f_{9}=2 a_{7} \left(a_{1} - a_{2}\right)\)\par
\noindent\(\begin{aligned} w_{9}&=- \frac{3}{128}\Bigl(a_{3} a_{4}\mathbin{-}a_{3} a_{5} \\ &\quad \mathbin{-}6 a_{6} a_{7}\Bigr)\end{aligned}\)\par\medskip
\noindent\(\begin{aligned} f_{14}&=-48 a_{1}^{2}\mathbin{-}48 a_{1} a_{2} \\ &\quad \mathbin{-}12 a_{2}^{2}\mathbin{+}4 a_{4} a_{6} \\ &\quad \mathbin{-}6 a_{4} a_{8}\mathbin{+}2 a_{5} a_{6} \\ &\quad \mathbin{-}3 a_{5} a_{8}\mathbin{-}1\end{aligned}\)\par
\noindent\(\begin{aligned} w_{14}&=- \frac{1}{64}\Bigl(3 a_{0}^{2}\mathbin{-}288 a_{0} a_{1} \\ &\quad \mathbin{-}108 a_{0} a_{2}\mathbin{-}135 a_{0} a_{4} \\ &\quad \mathbin{-}36 a_{0} a_{5}\mathbin{-}174 a_{1}^{2} \\ &\quad \mathbin{-}516 a_{1} a_{2}\mathbin{+}432 a_{1} a_{4} \\ &\quad \mathbin{+}108 a_{1} a_{5}\mathbin{-}174 a_{2}^{2} \\ &\quad \mathbin{-}110\Bigr)\end{aligned}\)\par\medskip
\noindent\(\begin{aligned} f_{15}&=-16 a_{1}^{2}\mathbin{-}16 a_{1} a_{2} \\ &\quad \mathbin{-}4 a_{2}^{2}\mathbin{-}2 a_{4} a_{8} \\ &\quad \mathbin{-}a_{5} a_{8}\mathbin{-}1\end{aligned}\)\par
\noindent\(\begin{aligned} w_{15}&=\frac{9}{64}\Bigl(a_{0}^{2}\mathbin{-}96 a_{0} a_{1} \\ &\quad \mathbin{-}36 a_{0} a_{2}\mathbin{-}21 a_{0} a_{4} \\ &\quad \mathbin{-}58 a_{1}^{2}\mathbin{-}172 a_{1} a_{2} \\ &\quad \mathbin{+}80 a_{1} a_{4}\mathbin{+}4 a_{1} a_{5} \\ &\quad \mathbin{-}58 a_{2}^{2}\mathbin{-}26\Bigr)\end{aligned}\)\par\medskip
\noindent\(\begin{aligned} f_{17}&=-a_{0}^{2}\mathbin{+}6 a_{4} a_{6} \\ &\quad \mathbin{-}8 a_{4} a_{8}\mathbin{+}2 a_{5} a_{6} \\ &\quad \mathbin{-}4 a_{5} a_{8}\mathbin{+}4\end{aligned}\)\par
\noindent\(\begin{aligned} w_{17}&=- \frac{3}{64}\Bigl(18 a_{0} a_{4}\mathbin{+}9 a_{0} a_{5} \\ &\quad \mathbin{-}48 a_{1} a_{4}\mathbin{-}24 a_{1} a_{5} \\ &\quad \mathbin{-}20\Bigr)\end{aligned}\)\par\medskip
\noindent\(\begin{aligned} f_{18}&=- \frac{1}{2}\Bigl(2 a_{0}^{2}\mathbin{-}2 a_{3}^{2} \\ &\quad \mathbin{-}a_{4} a_{6}\mathbin{+}a_{5} a_{6} \\ &\quad \mathbin{-}4\Bigr)\end{aligned}\)\par
\noindent\(\displaystyle w_{18}=- \frac{3}{16}\)\par\medskip
\noindent\(\displaystyle f_{19}=- a_{6} \left(a_{3} + a_{7}\right)\)\par
\noindent\(\begin{aligned} w_{19}&=\frac{3}{32}\Bigl(3 a_{0} a_{7}\mathbin{-}a_{1} a_{3} \\ &\quad \mathbin{+}a_{2} a_{3}\Bigr)\end{aligned}\)\par\medskip
\noindent\(\begin{aligned} f_{25}&=- \frac{1}{2}\Bigl(2 a_{1}^{2}\mathbin{-}4 a_{1} a_{2} \\ &\quad \mathbin{+}2 a_{2}^{2}\mathbin{-}a_{4} a_{6} \\ &\quad \mathbin{+}a_{5} a_{6}\mathbin{-}2 a_{7}^{2} \\ &\quad \mathbin{+}4\Bigr)\end{aligned}\)\par
\noindent\(\displaystyle w_{25}=- \frac{9}{16}\)\par\medskip
With all omitted multipliers equal to zero, direct expansion gives
\[
\boxed{\displaystyle \sum_{i\in\{0,1,2,4,7,9,14,15,17,18,19,25\}} w_i f_i=1.}
\]
The multiplier degrees are at most two.  Thus the residual order-two
ideal is the unit ideal over $\mathbb Q[a_0,\ldots,a_8]$.
\endgroup


\section{Claim-status and verification index}\label{app:verification}

The full cross-paper claim ledger is supplement material.  Each extracted
article will carry a shorter claim table containing only its own theorem
dependencies and adjacent fail-closed boundaries.

The residual computation used in the first draft was frozen at repository
commit
\begin{center}
\texttt{e928f257c25099eb534eb34109dfc1dc6a3127a1}.
\end{center}
The global BGG bridge was added in the subsequent theorem revision based at
\begin{center}
\texttt{c471b99f5e3708e692b1c25238f6272c9e29b48f}.
\end{center}
The selected algebraic BV--BFV suspension, polarization, and pairing transfer
were frozen at
\begin{center}
\texttt{8a7e7821f8cd4af5798fd1cb7a962f1da69cdf86}.
\end{center}
The companion snapshot
records all three layers, the environment, hashes, expected-failure guards, and
command lines.  The exhaustive publication verification is run with
\begin{verbatim}
python3 symbolic/verify_conformal_paper_free.py --reproduce --timeout 1800
\end{verbatim}
The claim/evidence index for the terminal covariant theorem is the generated
proof ledger at
\begin{center}\small
\path{covariant_completion/generated/}\par
\path{covariant_H4_proof_ledger.md}.
\end{center}
It binds each required claim to its authoritative certificate digest and
reproduction command.

{\scriptsize
\begin{longtable}{>{\raggedright\arraybackslash}p{0.25\textwidth}
                >{\raggedright\arraybackslash}p{0.17\textwidth}
                >{\raggedright\arraybackslash}p{0.48\textwidth}}
\toprule
Claim & Status & Evidence / boundary\\
\midrule
\endfirsthead
\toprule
Claim & Status & Evidence / boundary\\
\midrule
\endhead
Formal pure-Weyl detour identities & Exact & Branson--Gover theorem plus
action-normalized factorization and finite scalar-block audits.\\
Global smooth deformation cohomology and curvature isomorphism & Exact
smooth theorem & The flat BGG fine resolution and
$H^\bullet(\R\times S^3)$ give $H^1_{\rm def}=H^2_{\rm def}=0$ and hence
\eqref{eq:curvature-isomorphism}.  This smooth theorem does not itself
establish the later energy-mode or field Cauchy completions.\\
Euclidean polynomial detour ranks, degrees $2$--$6$ & Exact finite
certificate & Rational matrices verify kernels, ranks, and the $15$ residual
zero modes; they are an independent convention audit rather than the source
of the global theorem.\\
Algebraic cylinder realization and stable generator coefficients & Exact
symbolic theorem & Full coordinate Weyl/Bach calculation at symbolic $n$,
nonzero $E/A/L$ pivots, explicit same-block metric preimages, character
exhaustion, invariant form, and cutoff-stable conformal brackets.\\
Quadratic minimal field-BV/raw-chain identification & Exact centered-buffer
certificate & The master-action tangent differential is constructed in full
metric and antifield variables.  Explicit inverse maps verify
$FQ_{\rm BV}=Q_{\rm raw}F$ and $GQ_{\rm raw}=Q_{\rm BV}G$ on all 2170 basis
vectors at energies two through five; gauge fixing is a separate proposition
and the dual global zero-mode sector is not included here.\\
Gauge-fixed/nonminimal equivalence & Exact centered-buffer certificate & The
stationary conformal-Landau fermion generates an invertible canonical shear.
All vector and scalar antighosts, multipliers, and antifield duals are retained;
the transported $p_{\rm gf},\jmath_{\rm gf},s_{\rm gf}$ contract them exactly.
The low-mode certificate verifies $P_ZQ_{\rm gf}=Q_{\rm gf}P_Z$ and preserves
all fifteen CKVs.\\
Canonical bulk endpoint duality and Taub codomain & Exact algebraic theorem &
Perfect endpoint pairings give
$\operatorname{coker}K^\sharp\cong(\ker K)^*$ with dual compact type
$4_{+1}\oplus7_0\oplus4_{-1}$.  Exact matrices construct $\Theta$, its
quotient section, and the direct-sum endpoint decomposition.  The quadratic
Kuranishi class evaluates to the certified Taub/moment map.  The time-slice
degree suspension is supplied by the next theorem, not by endpoint duality
alone.\\
Selected algebraic BV--BFV suspension and polarization & Exact algebraic
theorem & Canonical BFV compatibility fixes $\tau=\Theta$ and $\lambda=+1$.
The positive-frequency Lagrangian ledger gives
$\operatorname{Sym}(\W_+\oplus\W_-)\otimes\Lambda^\bullet\g^*$, represents
$b_a$ by $\iota_a$, and contracts every local and nonminimal doublet to its
vacuum.  One matter normalization and the cotangent orientation transfer the
pairing and give $G_{\rm field}=I_2$.\\
Completed residual Cartan homotopy reduction & Exact energy-mode theorem &
The maximal block-direct-sum $\overline Q$ is closed and nilpotent, every
total-degree block is finite, and bounded $\iota_D/\delta$ gives closed
range and unchanged centered cohomology.  This is a quotient only in the
full residual-gauging problem.\\
Centered vacuum and one-particle sectors & Exact & Semisimple Lie-algebra
cohomology for the vacuum; cutoff-complete absolute rank certificate for
one particle.\\
Centered two-particle sector and $G_{\rm res}=I_2$ & Exact residual theorem &
Independent relative and absolute calculations; empty incoming absolute
space and exact rank $53$.\\
Vertex descent and dynamical/topological split & Exact locally & Weight-four
descent, Chern--Weil transgression, and Euler--Lagrange quotient.  Global
theta effects are retained.\\
Energy-mode one-particle/Fock Krein completion & Exact & The normalized
$E/A/L$ basis gives an infinite-index Krein space, $\Gamma_s(\mathfrak J_1)$
gives the bosonic Fock fundamental symmetry, and weighted-shift domains give
closed Krein-adjoint conformal generators on a common core.\\
Reduced Lorentzian Cauchy--Sobolev realization & Exact cylinder theorem &
Tensor/vector curl factorization, action-derived residues, exact $E/A/L$
multiplicities, and normalized harmonics give
\eqref{eq:branch-sobolev-spaces} and a Krein-unitary map to the energy-mode
one-particle completion.  The physical TT and vector blocks are Green
hyperbolic.\\
Auxiliary $66$-to-$30$ Fourier SDR & Exact flat/Fourier theorem & The exact
ghost biwave, flat scalar symbols, support-local Stueckelberg/auxiliary
reduction, and symbol-level formal adjoints are certified.  This row does not
assert a curved scalar-wave symbol.\\
Curved auxiliary operator, retract, and current identities & Exact cylinder
theorems & The action-derived curved Hessian and four-row identities, the
all-row BV-canonical support-local SDR, and the $d+Q$ current comparison are
exact.  Thus \texttt{curved\_operator\_identity},
\texttt{curved\_deformation\_retract}, and
\texttt{curved\_current\_comparison} are true.\\
Scalar-wave witness on the $24$-field auxiliary bundle & Exact no-go theorem
& At a generic null covector $\rank E_2=11>9=\rank K_1$; hence no invertible
pointwise $J$ and first-order $C$ can make
$J^{-1}E_2+K_1C_1$ vanish on the null cone.\\
Physical null-symbol quotient and Weyl carrier & Exact symbol theorem & The
quotient has dimension two and transverse helicities $+2$ and $-2$.  The
reduced Weyl symbol is $\frac14I_2$.  No unreduced claim
$\ker\W_2=\im K_1$ is made.\\
Electric/magnetic Weyl--Cotton first-order system & Exact curved theorem &
The Cotton divergence supplies the required sixteen components.  The exact
$34\times26$ covariant table agrees with all 150 Weyl two-jets.  Its eight
primary and six secondary constraints give a formally equivalent adjusted
$26$-state system with positive symmetrizer, causal speeds, exact sourced
compatibility, and constraint propagation.  The all-level $E/A/L$ spectrum
and absence of Cotton duplication are symbolic theorems.\\
Prenormal auxiliary symbol and relative saddle diagnostics & Exact
fail-closed diagnostics & The generic principal symbol obeys
$(P_2-qI)^2=0$ with algebraic/wave/biwave Smith multiplicities $6/12/6$,
but the naive frozen lower completion has nonzero lower-order remainder.
The complete invariant ansatz has $\dim D_0=38$ and $\dim D_1=93$;
simultaneous $DP/PD$ cubic divisibility leaves a $45$-parameter family and
hence no cubic obstruction.  The complete $421$-variable sparse quadratic
PBW system is assembled; its order-two $2484\times190$ algebraic matrix has
rank $100$, and exact cokernel projection yields $2130$ nonzero polynomial
constraints with no constant obstruction.  The projected and lower-order
systems remain open.
The corrected spatial-projector curvature sign passes focused primal
coordinate-jet tests, and exact symmetrized-jet PBW inversion certifies the
composition backend.  The $48$-entry naive sorted-table defect is resolved by
the pairing-aware symmetrized-PBW adjoint: imposing $D^\sharp=D$ leaves
$44/24$ first/zeroth-order complement dimensions, a $21$-dimensional cubic
family and a $214$-parameter nonlinear sharp branch.  Its exact order-two
Schur gate has rank $52/100$ and $1050$ nonzero projected constraints, but
no constant obstruction or affine-linear consequence.  The four
bare-$\Box$ orientations are exactly inconsistent with rank $159/160$ and a
common one-row $f_{01}\to f_{00}$ obstruction.  A two-direction bilinear
split repairs its full $SO(3)$ orbit, but that fixed minimal split fails the
complete order-two system with rank $100/101$ and a one-row
$h_{22}\to f_{01}$ obstruction.  This does not decide the general
two-nontrivial-factor branch.
The smallest reciprocal relative saddle is the odd-adjoint pair $4+5$.
Its Schur complement is nonlocal.  The natural
$A_F=p_FA_{\rm eq}$, $S=A_F^\sharp$, $R=A_F^\sharp J_U$ realization has
balanced timelike rank at most $107/116$, zero temporal leading coefficient
and no positive symmetrizer.  This excludes only that realization, not larger
relative witnesses or a support-local first-order prolongation.  The complete
$162$-dimensional expanded-relative family contains three reciprocal
two-pair candidates, but its central cross-ranks are only $21/24$; the three
missing directions are rotation scalars which must be retained or locally
prolonged.\\
Curvature causal homotopy and final covariant $H^4$ & Exact covariant theorem;
two legacy implementation flags false &
Thirteen flags in \eqref{eq:open-curvature-flags} are true.  The local all-row prolongation, BV
differential, off-shell current comparison, and $356/30$ hybrid algebraic
projector are complete.  The retained $5\to10\to10\to5$ endpoint is
coefficient-complete and cyclic.  Its corrected canonical backward witness
has trace-free scalar biwave symbol.  The complete parallel invariant
two-wave-factor family is excluded by a twelve-multiplier degree-two
Nullstellensatz certificate.  Its upper curvature graph lift is exact;
the canonical middle lift is obstructed.  The minimal cyclic relative saddle
exists on all $386$ rows but has zero endpoint Schur correction, so it neither
helps nor obstructs the still-open endpoint-diagonal Green problem.  The
curved adjoint-tractor HPL compression resolves the former $48$ entries in
ten invariant channels and matches the full Bach block at orders zero through
four with normalization $-2$.  Parent transfer and the explicit trace/Weyl
shear give an exact causal homotopy on all $386$ rows.  No canonical
$D_{\rm TF}$ Green inverse or prolonged witness is claimed or needed for that
identity.  Its causal difference is a compact-to-global quasi-isomorphism;
cutoff sources recover the fifteen CKVs and dual endpoints with no duplicate
copy, and the cutoff commutator gives $SO(4,2)$ transport up to chain
homotopy.  The direct cyclic Green/current certificate identifies the causal,
prolonged, auxiliary, metric and $E/A/L$ pairings.  Consequently
\texttt{final\_covariant\_H4} is true and transports the already-complete
residual $H^4=\C^2$ and $G=I_2$ without recomputation.  The false
\texttt{prolonged\_green\_witness} and
\texttt{curvature\_causal\_green\_operators} flags are scoped to the
superseded canonical implementation;
no Hadamard theorem, group exponentiation, or uniqueness among boundary,
clock, asymptotic, or deparametrized polarizations is established.\\
Interpretation as a physical state space & Choice-dependent & Requires a
boundary and gauge principle that treats all fifteen generators, including
$D$, as gauge rather than global charges.\\
Quantum anomaly freedom and interacting probability & Not claimed & Outside
the scope of this paper.\\
\bottomrule
\end{longtable}}

\begin{thebibliography}{99}

\bibitem{BransonGover}
T.~Branson and A.~R.~Gover,
``The conformal deformation detour complex for the obstruction tensor,''
\href{https://arxiv.org/abs/math/0605192}{arXiv:math/0605192}.

\bibitem{GoverPeterson}
A.~R.~Gover and L.~J.~Peterson,
``The ambient obstruction tensor and the conformal deformation complex,''
Pacific J. Math. \textbf{226} (2006) 309--351,
\href{https://arxiv.org/abs/math/0408229}{arXiv:math/0408229}.

\bibitem{GoverSombergSoucek}
A.~R.~Gover, P.~Somberg, and V.~Sou\v{c}ek,
``Yang--Mills detour complexes and conformal geometry,''
Commun. Math. Phys. \textbf{278} (2008) 307--327,
\href{https://arxiv.org/abs/math/0606401}{arXiv:math/0606401}.

\bibitem{GreenHyperbolic}
M.~Benini, G.~Musante, and A.~Schenkel,
``Green hyperbolic complexes on Lorentzian manifolds,''
Commun. Math. Phys. (2023),
\href{https://doi.org/10.1007/s00220-023-04807-5}{doi:10.1007/s00220-023-04807-5},
\href{https://arxiv.org/abs/2207.04069}{arXiv:2207.04069}.

\bibitem{BaerGreenHyperbolic}
C.~B\"ar,
``Green-hyperbolic operators on globally hyperbolic spacetimes,''
Commun. Math. Phys. \textbf{333} (2015) 1585--1615,
\href{https://arxiv.org/abs/1310.0738}{arXiv:1310.0738}.

\bibitem{JorgensenTian}
P.~E.~T.~Jorgensen and F.~Tian,
``Unbounded operators, Lie algebras, and local representations,''
J. Math. Phys. \textbf{55} (2014) 111702,
\href{https://arxiv.org/abs/1406.6966}{arXiv:1406.6966}.

\bibitem{CapBGG}
A.~\v{C}ap,
``Overdetermined systems, conformal geometry, and the BGG complex,''
in \emph{Symmetries and Overdetermined Systems of Partial Differential
Equations}, IMA Vol. Math. Appl. \textbf{144} (Springer, 2008) 1--24,
\href{https://arxiv.org/abs/math/0610225}{arXiv:math/0610225}.

\bibitem{CapDeformations}
A.~\v{C}ap,
``Infinitesimal automorphisms and deformations of parabolic geometries,''
J. Eur. Math. Soc. \textbf{10} (2008) 415--437,
\href{https://arxiv.org/abs/math/0508535}{arXiv:math/0508535}.

\bibitem{CapHuPoincare}
A.~\v{C}ap and K.~Hu,
``Bounded Poincar\'e operators for twisted and BGG complexes,''
J. Math. Pures Appl. \textbf{179} (2023) 253--276,
\href{https://arxiv.org/abs/2304.07185}{arXiv:2304.07185}.

\bibitem{HamadaHorata}
K.-j.~Hamada and S.~Horata,
``Conformal algebra and physical states in non-critical 3-brane on
$\R\times S^3$,''
Prog. Theor. Phys. \textbf{110} (2003) 1169--1210,
\href{https://arxiv.org/abs/hep-th/0307008}{arXiv:hep-th/0307008}.

\bibitem{HamadaBuildingBlocks}
K.-j.~Hamada,
``Building blocks of physical states in a non-critical 3-brane on
$\R\times S^3$,''
Int. J. Mod. Phys. A \textbf{20} (2005) 5353--5398,
\href{https://arxiv.org/abs/hep-th/0402136}{arXiv:hep-th/0402136}.

\bibitem{HamadaCFT}
K.-j.~Hamada,
``Conformal field theory on $\R\times S^3$ from quantized gravity,''
Int. J. Mod. Phys. A \textbf{24} (2009) 3073--3110,
\href{https://arxiv.org/abs/0811.1647}{arXiv:0811.1647}.

\bibitem{HamadaBRST}
K.-j.~Hamada,
``BRST analysis of physical fields and states for 4D quantum gravity on
$\R\times S^3$,''
Phys. Rev. D \textbf{86} (2012) 124006,
\href{https://arxiv.org/abs/1202.4538}{arXiv:1202.4538}.

\bibitem{KuzenkoPonds}
S.~M.~Kuzenko and M.~Ponds,
``Conformal geometry and (super)conformal higher-spin gauge theories,''
JHEP \textbf{05} (2019) 113,
\href{https://arxiv.org/abs/1902.08010}{arXiv:1902.08010}.

\bibitem{BeccariaBekaertTseytlin}
M.~Beccaria, X.~Bekaert, and A.~A.~Tseytlin,
``Partition function of free conformal higher spin theory,''
JHEP \textbf{08} (2014) 113,
\href{https://arxiv.org/abs/1406.3542}{arXiv:1406.3542}.

\bibitem{MetsaevOrdinaryDerivative}
R.~R.~Metsaev,
``Ordinary-derivative formulation of conformal low-spin fields,''
JHEP \textbf{01} (2008) 064,
\href{https://arxiv.org/abs/0707.4437}{arXiv:0707.4437}.

\bibitem{NutmaTaronna}
T.~Nutma and M.~Taronna,
``On conformal higher spin wave operators,''
JHEP \textbf{06} (2014) 066,
\href{https://arxiv.org/abs/1404.7452}{arXiv:1404.7452}.

\bibitem{BoulangerHenneaux}
N.~Boulanger and M.~Henneaux,
``A derivation of Weyl gravity,''
Annalen Phys. \textbf{10} (2001) 935--964,
\href{https://arxiv.org/abs/hep-th/0106065}{arXiv:hep-th/0106065}.

\bibitem{BarnichBrandtHenneaux}
G.~Barnich, F.~Brandt, and M.~Henneaux,
``Local BRST cohomology in gauge theories,''
Phys. Rept. \textbf{338} (2000) 439--569,
\href{https://arxiv.org/abs/hep-th/0002245}{arXiv:hep-th/0002245}.

\bibitem{BarnichAntifield}
G.~Barnich, F.~Brandt, and M.~Henneaux,
``Local BRST cohomology in the antifield formalism: I. General theorems,''
Commun. Math. Phys. \textbf{174} (1995) 57--92,
\href{https://arxiv.org/abs/hep-th/9405109}{arXiv:hep-th/9405109}.

\bibitem{BarnichHenneauxHurthSkenderis}
G.~Barnich, M.~Henneaux, T.~Hurth, and K.~Skenderis,
``Cohomological analysis of gauged-fixed gauge theories,''
Phys. Lett. B \textbf{492} (2000) 376--384,
\href{https://arxiv.org/abs/hep-th/9910201}{arXiv:hep-th/9910201}.

\bibitem{DecaniniFolacciFKWC}
Y.~D\'ecanini and A.~Folacci,
``FKWC-bases and geometrical identities for classical and quantum field
theories in curved spacetime,''
\href{https://arxiv.org/abs/0805.1595}{arXiv:0805.1595}.

\bibitem{MartinGarciaYllanesPortugalInvar}
J.~M.~Mart\'in-Garc\'ia, D.~Yllanes, and R.~Portugal,
``The Invar tensor package: Differential invariants of Riemann,''
\href{https://arxiv.org/abs/0802.1274}{arXiv:0802.1274}.

\bibitem{CattaneoMnevReshetikhin}
A.~S.~Cattaneo, P.~Mnev, and N.~Reshetikhin,
``Classical BV theories on manifolds with boundary,''
Commun. Math. Phys. \textbf{332} (2014) 535--603,
\href{https://arxiv.org/abs/1201.0290}{arXiv:1201.0290}.

\bibitem{IyerWald}
V.~Iyer and R.~M.~Wald,
``Some properties of Noether charge and a proposal for dynamical black hole
entropy,''
Phys. Rev. D \textbf{50} (1994) 846--864,
\href{https://arxiv.org/abs/gr-qc/9403028}{arXiv:gr-qc/9403028}.

\bibitem{MetsaevBRSTBV}
R.~R.~Metsaev,
``BRST--BV approach to conformal fields,''
J. Phys. A \textbf{49} (2016) 175401,
\href{https://arxiv.org/abs/1511.01836}{arXiv:1511.01836}.

\bibitem{DneprovGrigoriev}
I.~Dneprov and M.~Grigoriev,
``Presymplectic BV--AKSZ formulation of conformal gravity,''
\href{https://doi.org/10.1140/epjc/s10052-022-11082-6}
{doi:10.1140/epjc/s10052-022-11082-6},
\href{https://arxiv.org/abs/2208.02933}{arXiv:2208.02933}.

\bibitem{BarnichBekaertGrigoriev}
G.~Barnich, X.~Bekaert, and M.~Grigoriev,
``Notes on conformal invariance of gauge fields,''
J. Phys. A \textbf{48} (2015) 505402,
\href{https://arxiv.org/abs/1506.00595}{arXiv:1506.00595}.

\bibitem{AltasTekinLinearization}
E.~Altas and B.~Tekin,
``Linearization instability for generic gravity in AdS,''
Phys. Rev. D \textbf{97} (2018) 024028,
\href{https://arxiv.org/abs/1705.10234}{arXiv:1705.10234}.

\bibitem{AltasTekinTaub}
E.~Altas and B.~Tekin,
``Second order perturbation theory in general relativity: Taub charges as
integral constraints,''
Phys. Rev. D \textbf{99} (2019) 104078,
\href{https://arxiv.org/abs/1903.11982}{arXiv:1903.11982}.

\bibitem{ChevalleyEilenberg}
C.~Chevalley and S.~Eilenberg,
``Cohomology theory of Lie groups and Lie algebras,''
Trans. Amer. Math. Soc. \textbf{63} (1948) 85--124.

\bibitem{FischlerKundu}
W.~Fischler and S.~Kundu,
``Physical effects of the gravitational $\Theta$-parameter,''
JHEP \textbf{05} (2017) 098,
\href{https://arxiv.org/abs/1612.06010}{arXiv:1612.06010}.

\end{thebibliography}

\end{document}
